% STRATIGRAPHY — Geology Upper Sixth — Cours signé OnBuch+
% Official MINESEC syllabus. Compiler avec : tectonic GEO01-stratigraphy.tex
\newcommand{\DOCMATIERE}{Geology}
\newcommand{\DOCNIVEAU}{Upper Sixth}
\newcommand{\DOCMODULE}{Upper Sixth — Geology}
\newcommand{\DOCLECON}{1}
\newcommand{\DOCTITRE}{STRATIGRAPHY}
% OnBuch+ — préambule « cours signés » ENRICHI (compilé avec tectonic / XeLaTeX)
% Système visuel « Pop » d'OnBuch+ : fond crème (jamais blanc), bordures encre
% épaisses, ombres « dures » décalées, titres Archivo Black en capitales.
% Version MATHÉMATIQUES : graphes (pgfplots/tikz), boîtes pédagogiques
% (définition, propriété, méthode, attention, le savais-tu, exercice résolu,
% exercice, TP, bilan), page de garde et signature OnBuch+ ancrée en bas.
%
% Macros attendues dans main.tex AVANT \input{preamble} :
%   \DOCMATIERE  \DOCNIVEAU  \DOCMODULE  \DOCLECON (numéro)  \DOCTITRE
\documentclass[11pt]{article}

\usepackage{fontspec}
\usepackage[english]{babel}
\usepackage[a4paper,margin=2cm,top=3.4cm,bottom=2.8cm,headheight=1.4cm,footskip=1.3cm]{geometry}
\usepackage{amsmath,amssymb,amsfonts}
\usepackage{mathtools} % \xmapsto, \coloneqq... souvent utilisés par les modèles
\usepackage{cancel} % simplifications barrées (bilans de forces)
% \abs{z} (module, valeur absolue) et \norm{u} (norme), utilisés par les modèles
\providecommand{\abs}{}\let\abs\relax\DeclarePairedDelimiter{\abs}{\lvert}{\rvert}
\providecommand{\norm}{}\let\norm\relax\DeclarePairedDelimiter{\norm}{\lVert}{\rVert}
\usepackage[table]{xcolor} % [table] : fournit aussi \rowcolors (alternance de lignes)
\usepackage{pagecolor}
\usepackage{tikz}
\usetikzlibrary{arrows.meta,positioning,calc,shapes.geometric,shapes.misc,shapes.arrows,decorations.pathmorphing,decorations.pathreplacing,decorations.markings,patterns,shadows,mindmap,trees,fit,backgrounds,matrix,intersections,angles,quotes}
\usepackage{pgfplots}
\usepackage[version=4]{mhchem} % équations nucléaires \ce{^{235}_{92}U}
\usepackage[european,siunitx]{circuitikz} % schémas électriques
\pgfplotsset{compat=1.18}
\usepgfplotslibrary{fillbetween,groupplots} % aires entre courbes (calcul intégral)
\usepackage{enumitem}
\usepackage{fancyhdr}
% [most] charge toutes les bibliothèques tcolorbox (listings, minted, raster,
% documentation...) et gonfle inutilement la mémoire TeX sur un document long
% (~300 boîtes) : on ne charge que ce qui est réellement utilisé.
\usepackage[skins,breakable]{tcolorbox}
\usepackage{graphicx}
\usepackage{colortbl}
\usepackage{booktabs}
\usepackage{tabularx}
\usepackage{diagbox} % case d'en-tête diagonale des tableaux à double entrée
% Colonnes extensibles centrée / gauche / droite (C, L, R), souvent utilisées par les modèles
\newcolumntype{C}{>{\centering\arraybackslash}X}
\newcolumntype{L}{>{\raggedright\arraybackslash}X}
\newcolumntype{R}{>{\raggedleft\arraybackslash}X}
\usepackage{array}
\usepackage{multicol}
\usepackage{siunitx}
% parse-numbers=false : accepte aussi \SI{2,5 \times 10^{-3}}{...}
\sisetup{locale=UK,output-decimal-marker={.},group-minimum-digits=4,parse-numbers=false}
\DeclareSIUnit\ohm{\ensuremath{\symup{\Omega}}} % U+2126 absent de la police
\DeclareSIUnit\division{div} % réglages d'oscilloscope (ms/div, V/div)
\DeclareSIUnit\tr{tr} % tours (vitesse de rotation, ex. \SI{1500}{\tr\per\minute})
\usepackage{qrcode}
% \degree : commande fréquemment utilisée par les modèles pour un angle ou une
% température en dehors de \SI{}{} (non fournie par les paquets chargés).
\providecommand{\degree}{^{\circ}}
% \qed (fin de démonstration), utilisé par les modèles sans amsthm
\providecommand{\qed}{\hfill$\square$}
% \barre : double barre des valeurs interdites dans un tableau de variations (tkz-tab)
\providecommand{\barre}{\ensuremath{\|}}
% environnement proof (amsthm non chargé) utilisé par les modèles
\ifdefined\proof\else\newenvironment{proof}[1][Proof]{\par\noindent\textit{#1.}\ }{\qed\par}\fi
\usepackage{lastpage}
\usepackage{hyperref}
\hypersetup{hidelinks}

% ============================================================
% Palette « Pop » OnBuch+ — mêmes hex que la classe P (plus_ui.dart)
% ============================================================
\definecolor{poporange}{HTML}{F59321}
\definecolor{popdark}{HTML}{E07A0C}
\definecolor{popcream}{HTML}{FFF6E8}
\definecolor{popink}{HTML}{1C1714}
\definecolor{poppurple}{HTML}{7A5AE0}
\definecolor{popgreen}{HTML}{2E9E6A}
\definecolor{poppink}{HTML}{E0457B}
\definecolor{popblue}{HTML}{2F7DE1}
\definecolor{popgold}{HTML}{C98A12}
\definecolor{popmuted}{HTML}{8A7F76}
\definecolor{popline}{HTML}{EADFD2}
\definecolor{popgray}{HTML}{8A7F76}
\colorlet{popgrey}{popgray}
% Alias tolérants (le modèle nomme parfois les couleurs différemment)
\colorlet{popwhite}{white}
\colorlet{popblack}{popink}
\colorlet{popred}{poppink}
\colorlet{popyellow}{popgold}
% Teintes claires pour fonds de tableaux / zones
\colorlet{popblueL}{popblue!10!white}
\colorlet{poporangeL}{poporange!14!white}
\colorlet{popgreenL}{popgreen!12!white}
\colorlet{poppurpleL}{poppurple!12!white}
\colorlet{poppinkL}{poppink!12!white}
\colorlet{popgoldL}{popgold!14!white}

\pagecolor{popcream}

% ============================================================
% Typographie
% ============================================================
\defaultfontfeatures{Ligatures=TeX}
\setmainfont{PlusJakartaSans-Medium.ttf}[Path=../../../fonts/,BoldFont=PlusJakartaSans-Bold.ttf]
% Maths en sans-serif (Fira Math) avec chiffres et lettres droites de Plus
% Jakarta, pour que les formules se fondent dans le texte.
\usepackage[math-style=ISO,bold-style=ISO]{unicode-math}
\setmathfont{FiraMath-Regular.otf}[Path=../../../fonts/]
\setmathfont{PlusJakartaSans-Medium.ttf}[Path=../../../fonts/,range={up/num,up/{Latin,latin},\mathup}]
% (icomma retiré : décimales avec point en anglais)
% Caractères mathématiques Unicode absents des polices de texte (Plus Jakarta,
% Archivo Black) : rendus par leur équivalent mathématique, en texte comme en
% formule — sinon ils s'affichent en carré vide (ex. « Arithmétique dans ℤ »).
\usepackage{newunicodechar}
\newunicodechar{ℝ}{\ensuremath{\mathbb{R}}}
\newunicodechar{ℤ}{\ensuremath{\mathbb{Z}}}
\newunicodechar{ℂ}{\ensuremath{\mathbb{C}}}
\newunicodechar{ℕ}{\ensuremath{\mathbb{N}}}
\newunicodechar{ℚ}{\ensuremath{\mathbb{Q}}}
\newunicodechar{𝔻}{\ensuremath{\mathbb{D}}}
\newunicodechar{α}{\ensuremath{\alpha}}
\newunicodechar{β}{\ensuremath{\beta}}
\newunicodechar{γ}{\ensuremath{\gamma}}
\newunicodechar{δ}{\ensuremath{\delta}}
\newunicodechar{ε}{\ensuremath{\varepsilon}}
\newunicodechar{θ}{\ensuremath{\theta}}
\newunicodechar{λ}{\ensuremath{\lambda}}
\newunicodechar{μ}{\ensuremath{\mu}}
\newunicodechar{σ}{\ensuremath{\sigma}}
\newunicodechar{τ}{\ensuremath{\tau}}
\newunicodechar{φ}{\ensuremath{\varphi}}
\newunicodechar{ψ}{\ensuremath{\psi}}
\newunicodechar{ω}{\ensuremath{\omega}}
\newunicodechar{Σ}{\ensuremath{\Sigma}}
% majuscules produites par \MakeUppercase dans les titres (α-aminés -> Α)
\newunicodechar{Α}{\ensuremath{\alpha}}
\newunicodechar{Β}{\ensuremath{\beta}}
\newunicodechar{Γ}{\ensuremath{\Gamma}}
\newunicodechar{Δ}{\ensuremath{\Delta}}
\newunicodechar{Ε}{\ensuremath{\varepsilon}}
\newunicodechar{Θ}{\ensuremath{\Theta}}
\newunicodechar{Λ}{\ensuremath{\Lambda}}
\newunicodechar{Μ}{\ensuremath{\mu}}
\newunicodechar{Τ}{\ensuremath{\tau}}
\newunicodechar{Φ}{\ensuremath{\Phi}}
\newunicodechar{Ψ}{\ensuremath{\Psi}}
\newunicodechar{Ω}{\ensuremath{\Omega}}
\newunicodechar{↦}{\ensuremath{\mapsto}}
\newunicodechar{∈}{\ensuremath{\in}}
\newunicodechar{∉}{\ensuremath{\notin}}
\newunicodechar{∧}{\ensuremath{\wedge}}
\newunicodechar{⇔}{\ensuremath{\Leftrightarrow}}
\newunicodechar{⟺}{\ensuremath{\Longleftrightarrow}}
\newunicodechar{⇒}{\ensuremath{\Rightarrow}}
\newunicodechar{⟹}{\ensuremath{\Longrightarrow}}
\newunicodechar{∘}{\ensuremath{\circ}}
\newunicodechar{∪}{\ensuremath{\cup}}
\newunicodechar{∩}{\ensuremath{\cap}}
\newunicodechar{≡}{\ensuremath{\equiv}}
\newunicodechar{⊂}{\ensuremath{\subset}}
\newunicodechar{⊥}{\ensuremath{\perp}}
\newunicodechar{∃}{\ensuremath{\exists}}
\newunicodechar{∀}{\ensuremath{\forall}}
\newunicodechar{∛}{\ensuremath{\sqrt[3]{\,}}}
\newunicodechar{∜}{\ensuremath{\sqrt[4]{\,}}}
\newunicodechar{⃗}{\ensuremath{{}^{\to}}}
\newunicodechar{ⁿ}{\ifmmode^{n}\else\textsuperscript{\ensuremath{n}}\fi}
\newunicodechar{ˣ}{\ifmmode^{x}\else\textsuperscript{\ensuremath{x}}\fi}
\newunicodechar{ᵇ}{\ifmmode^{b}\else\textsuperscript{\ensuremath{b}}\fi}
\newunicodechar{ʳ}{\ifmmode^{r}\else\textsuperscript{\ensuremath{r}}\fi}
\newunicodechar{ᵘ}{\ifmmode^{u}\else\textsuperscript{\ensuremath{u}}\fi}
\newunicodechar{ʸ}{\ifmmode^{y}\else\textsuperscript{\ensuremath{y}}\fi}
\newunicodechar{ᵃ}{\ifmmode^{a}\else\textsuperscript{\ensuremath{a}}\fi}
\newunicodechar{ᵉ}{\ifmmode^{e}\else\textsuperscript{\ensuremath{e}}\fi}
\newunicodechar{ᵏ}{\ifmmode^{k}\else\textsuperscript{\ensuremath{k}}\fi}
\newunicodechar{ᵖ}{\ifmmode^{p}\else\textsuperscript{\ensuremath{p}}\fi}
\newunicodechar{ᴺ}{\ifmmode^{N}\else\textsuperscript{\ensuremath{N}}\fi}
\newunicodechar{ᶜ}{\ifmmode^{c}\else\textsuperscript{\ensuremath{c}}\fi}
\newunicodechar{ʰ}{\ifmmode^{h}\else\textsuperscript{\ensuremath{h}}\fi}
\newunicodechar{ᵠ}{\ifmmode^{q}\else\textsuperscript{\ensuremath{q}}\fi}
\newunicodechar{ₙ}{\ifmmode_{n}\else\textsubscript{\ensuremath{n}}\fi}
\newunicodechar{ₐ}{\ifmmode_{a}\else\textsubscript{\ensuremath{a}}\fi}
\newunicodechar{ᵢ}{\ifmmode_{i}\else\textsubscript{\ensuremath{i}}\fi}
\newunicodechar{ᵣ}{\ifmmode_{r}\else\textsubscript{\ensuremath{r}}\fi}
\newunicodechar{ₖ}{\ifmmode_{k}\else\textsubscript{\ensuremath{k}}\fi}
\newunicodechar{ᵦ}{\ifmmode_{\beta}\else\textsubscript{\ensuremath{\beta}}\fi}
\newunicodechar{ₓ}{\ifmmode_{x}\else\textsubscript{\ensuremath{x}}\fi}
\newfontfamily\popdisplay{ArchivoBlack-Regular.ttf}[Path=../../../fonts/]
\newfontfamily\poplogo{SpaceGrotesk-Bold.ttf}[Path=../../../fonts/]
\newfontfamily\popsemi{PlusJakartaSans-SemiBold.ttf}[Path=../../../fonts/]
\newfontfamily\popxbold{PlusJakartaSans-ExtraBold.ttf}[Path=../../../fonts/]
\newfontfamily\popxboldls{PlusJakartaSans-ExtraBold.ttf}[Path=../../../fonts/,LetterSpace=3.0]

\setlist[enumerate]{leftmargin=1.7em,itemsep=5pt,topsep=4pt}
\setlist[itemize]{leftmargin=1.7em,itemsep=4pt,topsep=4pt,label={\color{poporange}\textbullet}}
\setlength{\parindent}{0pt}
\setlength{\parskip}{5pt}
\renewcommand{\arraystretch}{1.25}

% pgfplots : style Pop par défaut
\pgfplotsset{
  every axis/.append style={axis line style={popink,line width=1pt},tick style={popink},
    grid style={popline,line width=0.5pt},label style={font=\small},tick label style={font=\footnotesize}},
  popcurve/.style={poporange,line width=1.8pt,no markers,smooth},
}
% Même style disponible dans un \draw TikZ simple (hors pgfplots)
\tikzset{popcurve/.style={poporange,line width=1.8pt}}
\tikzset{popgrid/.style={popline,very thin}}
\colorlet{popcurve}{poporange} % le nom du style est parfois utilisé comme couleur

% ============================================================
% Pill / squiggle
% ============================================================
\newcommand{\poppill}[3]{%
  \tikz[baseline=-0.55ex]\node[draw=popink,line width=1.2pt,rounded corners=2.6pt,%
    fill=#2,inner xsep=6.5pt,inner ysep=3pt]{%
    \popxboldls\fontsize{7}{7}\selectfont\color{#3}\MakeUppercase{#1}};%
}
\newcommand{\popsquiggle}[1]{%
  \tikz[baseline=-0.4ex]{
    \draw[#1,line width=2.6pt,line cap=round]
      plot [smooth] coordinates {(0,0.09) (0.34,0.01) (0.68,0.21) (1.02,0.01) (1.36,0.10)};
  }%
}
% Mot-clé surligné (marqueur orange sous le texte)
\newcommand{\cle}[1]{{\popxbold\color{popdark}#1}}

% ============================================================
% En-tête / pied
% ============================================================
\pagestyle{fancy}
\fancyhf{}
\renewcommand{\headrulewidth}{0pt}
\renewcommand{\footrulewidth}{0pt}
\lhead{\raisebox{-0.15ex}{\poplogo\fontsize{13}{13}\selectfont\color{popink}OnBuch\color{poporange}+}}
\rhead{\poppill{\DOCMATIERE\ \textbullet\ \DOCNIVEAU\ \textbullet\ Chapter \DOCLECON}{white}{popink}}
\lfoot{\popxbold\fontsize{7.6}{7.6}\selectfont\color{popmuted}\MakeUppercase{Signed course OnBuch+}\ \textbullet\ {\popsemi\DOCTITRE}}
\rfoot{\tikz[baseline=-0.6ex]\node[rounded corners=8.5pt,fill=popink,minimum height=17pt,minimum width=17pt,inner xsep=5pt,inner sep=0]{\popxbold\fontsize{8}{8}\selectfont\color{popcream}\thepage\,/\,\pageref*{LastPage}};}

% ============================================================
% Titres
% ============================================================
\newcommand{\chaptitre}[2]{%
  \begin{tcolorbox}[enhanced,colback=poporange,colframe=popink,boxrule=2.6pt,arc=11pt,
      drop shadow=popink,left=15pt,right=15pt,top=12pt,bottom=11pt,before skip=3pt,after skip=18pt]
    \poppill{#2}{popink}{popcream}\par\vspace{9pt}
    {\raggedright\hyphenpenalty=10000\exhyphenpenalty=10000\popdisplay\fontsize{22}{24}\selectfont\color{popcream}\MakeUppercase{#1}\par}
  \end{tcolorbox}
}
% Section numérotée : pastille orange avec le numéro + titre Archivo
\newcounter{popsec}
\newcommand{\coursec}[1]{%
  \stepcounter{popsec}\setcounter{popsub}{0}%
  \par\vspace{16pt}\noindent
  \begin{minipage}{\linewidth}
  \tikz[baseline=-0.7ex]\node[draw=popink,line width=1.6pt,fill=poporange,rounded corners=4pt,minimum size=22pt,inner sep=0]{\popdisplay\fontsize{12}{12}\selectfont\color{popcream}\thepopsec};%
  \hspace{8pt}{\popdisplay\fontsize{15}{17}\selectfont\color{popink}\MakeUppercase{#1}}\par
  \vspace{4pt}\noindent{\color{poporange}\rule{36pt}{3.6pt}}
  \end{minipage}\par\nopagebreak\vspace{8pt}
}
\newcounter{popsub}
\newcommand{\courssub}[1]{%
  \stepcounter{popsub}%
  \par\vspace{9pt}\noindent{\popxbold\fontsize{11.5}{13}\selectfont\color{popink}{\color{poporange}\thepopsec.\thepopsub}\hspace{6pt}#1}\par\nopagebreak\vspace{4pt}
}

% ============================================================
% Boîtes pédagogiques — même recette : fond blanc, bordure épaisse colorée,
% ombre dure encre, badge plein. Titre optionnel [#1].
% ============================================================
\tcbset{popbase/.style={breakable,enhanced,colback=white,boxrule=2.3pt,arc=9pt,
  shadow={3pt}{-3pt}{0pt}{popink},left=14pt,right=14pt,top=11pt,bottom=11pt,before skip=11pt,after skip=15pt,
  fontupper=\popsemi\fontsize{10.6}{14.5}\selectfont\color{popink},
  fontlower=\popsemi\fontsize{10.6}{14.5}\selectfont\color{popink}}}
% Coche dessinée (le glyphe ✓ n'existe pas dans les polices du document)
\newcommand{\popcheck}[1]{\tikz[baseline=-0.5ex]\draw[#1,line width=1.6pt,line cap=round,line join=round](-0.13,0.0)--(-0.04,-0.09)--(0.13,0.1);}
\providecommand{\coche}{\popcheck{popgreen}}
\providecommand{\resultat}[1]{\par\smallskip\noindent\fcolorbox{popgreen}{popgreenL}{\parbox{\dimexpr\linewidth-2\fboxsep-2\fboxrule\relax}{\textbf{Result.} #1}}\par\smallskip}
\newcommand{\popboxhead}[3]{\poppill{#1}{#2}{white}\ifx\relax#3\relax\else\hspace{6pt}{\popxbold\fontsize{10.6}{12}\selectfont\color{popink}#3}\fi\par\vspace{7pt}}

\newtcolorbox{aretenir}[1][]{popbase,colframe=popgreen,before upper={\popboxhead{Key points}{popgreen}{#1}}}
\newtcolorbox{exemplebox}[1][]{popbase,colframe=poporange,before upper={\popboxhead{Exemple}{poporange}{#1}}}
\newtcolorbox{definition}[1][]{popbase,colframe=popblue,before upper={\popboxhead{Definition}{popblue}{#1}}}
\newtcolorbox{propriete}[1][]{popbase,colframe=poppurple,before upper={\popboxhead{Property}{poppurple}{#1}}}
\newtcolorbox{methode}[1][]{popbase,colframe=popgold,before upper={\popboxhead{Method}{popgold}{#1}}}
\newtcolorbox{attention}[1][]{popbase,colframe=poppink,before upper={\popboxhead{Warning}{poppink}{#1}}}
\newtcolorbox{experience}[1][]{popbase,colframe=popblue,colback=popblueL,before upper={\popboxhead{Experiment}{popblue}{#1}}}
% « Le savais-tu ? » : carte violette pleine (contexte, culture, Cameroun)
\newtcolorbox{savaistu}[1][]{popbase,colback=poppurple,colframe=popink,
  fontupper=\popsemi\fontsize{10.4}{14.2}\selectfont\color{white},
  before upper={\poppill{Did you know?}{white}{poppurple}\ifx\relax#1\relax\else\hspace{6pt}{\popxbold\color{white}#1}\fi\par\vspace{7pt}}}
% Exercice résolu : énoncé en haut, \tcblower puis la solution sur fond crème
\newcounter{popexo}\newcounter{popexr}
\newtcolorbox{exoresolu}[1][]{popbase,colframe=poporange,colbacklower=popcream,
  segmentation style={popink,line width=1.2pt,dashed},
  before upper={\stepcounter{popexr}\popboxhead{Worked example \thepopexr}{poporange}{#1}},
  before lower={\poppill{Solution}{popink}{popcream}\par\vspace{6pt}}}
% Exercice (non résolu en place) — la correction est dans la partie Corrigés
\newtcolorbox{exercice}[1][]{popbase,colframe=popink,
  before upper={\stepcounter{popexo}\popboxhead{Exercise \thepopexo}{popink}{#1}}}
% Corrigé
\newtcolorbox{corrige}[1][]{popbase,colframe=popgreen,colback=popgreenL,
  before upper={\popboxhead{Solution}{popgreen}{#1}}}
% Objectifs / prérequis / bilan
\newtcolorbox{objectifs}{popbase,colframe=popink,colback=poporangeL,
  before upper={\popboxhead{Chapter objectives}{popink}{}}}
\newtcolorbox{prerequis}{popbase,colframe=popink,colback=popblueL,
  before upper={\popboxhead{Prerequisites}{popblue}{}}}
\newtcolorbox{bilan}{popbase,colframe=popink,colback=white,boxrule=2.8pt,
  before upper={\popboxhead{Fiche bilan}{poporange}{L'essentiel à connaître pour l'examen}}}
% Figure encadrée avec légende
\newtcolorbox{popfigure}[1][]{enhanced,breakable=false,colback=white,colframe=popline,boxrule=1.4pt,arc=8pt,
  left=10pt,right=10pt,top=10pt,bottom=8pt,before skip=10pt,after skip=12pt,halign=center,
  lower separated=false,halign lower=center,
  fontlower=\popsemi\fontsize{9}{11.5}\selectfont\color{popmuted},
  #1}
\newcommand{\legende}[1]{\par\vspace{6pt}{\popsemi\fontsize{9}{11.5}\selectfont\color{popmuted}#1}}

% ============================================================
% Signature OnBuch+
% ============================================================
\newcommand{\poplogoinline}[1]{{\poplogo\fontsize{#1}{#1}\selectfont\color{popink}OnBuch\color{poporange}+}}
\newcommand{\signatureonbuch}{%
  \begin{tcolorbox}[enhanced,colback=popink,colframe=popink,boxrule=2.2pt,arc=10pt,
      drop shadow=poporange,left=14pt,right=14pt,top=10pt,bottom=10pt,before skip=0pt,after skip=0pt]
    \tikz[baseline=-0.7ex]\node[circle,fill=popgreen,draw=popcream,line width=1.3pt,minimum size=22pt,inner sep=0]{\popcheck{white}};%
    \hspace{9pt}\begin{minipage}[c]{0.62\linewidth}
      {\popxbold\fontsize{12}{13}\selectfont\color{popcream}Signed\ \textbullet\ {\poplogo OnBuch\color{poporange}+}}\par\vspace{2pt}
      {\raggedright\popsemi\fontsize{8}{10.5}\selectfont\color{popline}\MakeUppercase{OnBuch+ teaching team}\\\MakeUppercase{Follows the official MINESEC syllabus}\par}
    \end{minipage}\hfill
    {\poplogo\fontsize{22}{22}\selectfont\color{popcream}OnBuch\color{poporange}+}
  \end{tcolorbox}%
}

% ============================================================
% Page de garde
% #1 titre · #2 matière · #3 niveau · #4 module · #5 n° leçon · #6 durée ·
% #7 description / objectifs (texte ou itemize)
% ============================================================
\newcommand{\pagegarde}[7]{%
  \thispagestyle{empty}
  \newgeometry{a4paper,margin=1.7cm,top=1.6cm,bottom=1.4cm}%
  \begin{tikzpicture}[remember picture,overlay]
    \fill[poporange!18] ([xshift=-3.2cm,yshift=-3.6cm]current page.north east) circle (4.6cm);
    \fill[poppurple!14] ([xshift=2.4cm,yshift=7.6cm]current page.south west) circle (3.8cm);
  \end{tikzpicture}%
  {\poplogo\fontsize{30}{30}\selectfont\color{popink}OnBuch\color{poporange}+}\hfill
  \raisebox{0.6ex}{\poppill{Signed course \textbullet\ Premium}{popink}{popcream}}\par
  \vspace{1pt}\popsquiggle{poppurple}\par
  \vspace{8pt}
  \begin{tcolorbox}[enhanced,colback=poporange,colframe=popink,boxrule=2.8pt,arc=13pt,
      drop shadow=popink,left=18pt,right=18pt,top=12pt,bottom=13pt,after skip=10pt]
    \poppill{#2 \textbullet\ #3}{popink}{popcream}\hspace{5pt}\poppill{#4}{white}{popink}\par\vspace{8pt}
    {\popdisplay\fontsize{14}{14}\selectfont\color{popink}CHAPTER #5}\par\vspace{4pt}
    {\raggedright\hyphenpenalty=10000\exhyphenpenalty=10000\popdisplay\fontsize{25}{27}\selectfont\color{popcream}\MakeUppercase{#1}\par}
    \vspace{8pt}
    {\popxbold\fontsize{9.5}{9.5}\selectfont\color{popink}\MakeUppercase{Suggested time: #6}}
  \end{tcolorbox}
  \begin{tcolorbox}[enhanced,colback=white,colframe=popink,boxrule=2.2pt,arc=10pt,
      drop shadow=popink,left=16pt,right=16pt,top=10pt,bottom=10pt,after skip=8pt]
    \poppill{In this chapter}{poppurple}{white}\par\vspace{5pt}
    \popsemi\fontsize{9.3}{12.6}\selectfont\color{popink}#7
  \end{tcolorbox}
  \noindent\begin{minipage}[t]{0.487\linewidth}
    \begin{tcolorbox}[enhanced,colback=popgreen,colframe=popink,boxrule=2.2pt,arc=10pt,
        drop shadow=popink,left=11pt,right=11pt,top=10pt,bottom=10pt,height=3.1cm,valign=center]
      \tcbox[colback=white,colframe=popink,boxrule=1.3pt,arc=5pt,boxsep=0pt,nobeforeafter,
        left=3pt,right=3pt,top=3pt,bottom=3pt,valign=center]{\qrcode[height=1.9cm]{https://play.google.com/store/apps/details?id=com.luvvix.onbuch&hl=en}}%
      \hspace{8pt}\begin{minipage}[c]{0.5\linewidth}
        {\popdisplay\fontsize{11}{12.5}\selectfont\color{white}\MakeUppercase{Download OnBuch}}\par\vspace{4pt}
        {\raggedright\popsemi\fontsize{8.4}{11}\selectfont\color{white}Free \textbullet\ Google Play\\AI tutor, past papers, results\par}
      \end{minipage}
    \end{tcolorbox}
  \end{minipage}\hfill
  \begin{minipage}[t]{0.487\linewidth}
    \begin{tcolorbox}[enhanced,colback=poppurple,colframe=popink,boxrule=2.2pt,arc=10pt,
        drop shadow=popink,left=11pt,right=11pt,top=10pt,bottom=10pt,height=3.1cm,valign=center]
      \tikz[baseline=-0.7ex]\node[circle,fill=white,draw=popink,line width=1.3pt,minimum size=1.6cm,inner sep=0]{\popdisplay\fontsize{24}{24}\selectfont\color{poppurple}+};%
      \hspace{8pt}\begin{minipage}[c]{0.6\linewidth}
        {\popdisplay\fontsize{11}{12.5}\selectfont\color{white}\MakeUppercase{Go OnBuch+}}\par\vspace{4pt}
        {\raggedright\popsemi\fontsize{8.4}{11}\selectfont\color{white}All signed courses, unlimited AI tutor, PDF sheets — OnBuch+ tab in the app\par}
      \end{minipage}
    \end{tcolorbox}
  \end{minipage}
  \vspace{8pt}
  \popcontenu
  \vfill
  \signatureonbuch
  \restoregeometry
  \newpage
}

% Rangée « Ce cours contient » (page de garde)
\newcommand{\poptile}[3]{% #1 couleur #2 gros texte #3 légende
  \begin{tcolorbox}[enhanced,colback=#1,colframe=popink,boxrule=2pt,arc=9pt,shadow={2.5pt}{-2.5pt}{0pt}{popink},
      left=6pt,right=6pt,top=9pt,bottom=9pt,halign=center,valign=center,height=1.75cm,width=\linewidth,nobeforeafter]
    {\popdisplay\fontsize{12.5}{13}\selectfont\color{white}\MakeUppercase{#2}}\par\vspace{3pt}
    {\centering\popsemi\fontsize{7.8}{9.5}\selectfont\color{white}#3\par}
  \end{tcolorbox}}
\newcommand{\popcontenu}{%
  \par\noindent{\popxboldls\fontsize{7.5}{7.5}\selectfont\color{popmuted}THIS SIGNED COURSE CONTAINS}\par\vspace{7pt}
  \noindent\begin{minipage}[t]{0.235\linewidth}\poptile{popblue}{Course}{complete, illustrated, step by step}\end{minipage}\hfill
  \begin{minipage}[t]{0.235\linewidth}\poptile{poporange}{Methods}{and worked examples}\end{minipage}\hfill
  \begin{minipage}[t]{0.235\linewidth}\poptile{poppink}{Exercises}{exam style + solutions}\end{minipage}\hfill
  \begin{minipage}[t]{0.235\linewidth}\poptile{popgreen}{Summary sheet}{the essentials to remember}\end{minipage}\par}

% Sommaire
\newtcolorbox{sommaire}{enhanced,colback=white,colframe=popink,boxrule=2.2pt,arc=10pt,
  drop shadow=popink,left=18pt,right=18pt,top=13pt,bottom=13pt,before skip=6pt,after skip=20pt,
  before upper={\poppill{Contents}{poppurple}{white}\par\vspace{9pt}\popsemi\fontsize{10.6}{14.5}\selectfont\color{popink}}}

% Signature de fin de cours — poussée tout en bas de la dernière page
\newcommand{\findecours}{%
  \par\vfill\nopagebreak
  \begin{center}\popsquiggle{poporange}\end{center}\vspace{4pt}
  \signatureonbuch
}
\AtBeginDocument{\def\llbracket{\mathopen{[\mkern-3.5mu[}}\def\rrbracket{\mathclose{]\mkern-3.5mu]}}} % intervalles d'entiers (glyphe ⟦ absent de la police)
% Glyphes absents de Fira Math : redéfinis à partir de glyphes disponibles
\AtBeginDocument{%
  \def\setminus{\mathbin{\backslash}}\let\smallsetminus\setminus
  \def\vdots{\mathord{\vbox{\baselineskip3.2pt\lineskiplimit0pt\kern4pt\hbox{.}\hbox{.}\hbox{.}}}}%
  \def\ddots{\mathinner{\mkern1mu\raise6.5pt\hbox{.}\mkern2mu\raise3.6pt\hbox{.}\mkern2mu\raise0.7pt\hbox{.}\mkern1mu}}%
  \def\triangle{\mathord{\scalebox{0.9}{$\symup{\Delta}$}}}%
  \def\nabla{\mathord{\raisebox{\depth}{\scalebox{1}[-1]{$\symup{\Delta}$}}}}%
  \def\checkmark{\popcheck{popdark}}%
  \def\star{\mathbin{\ast}}\def\bigstar{\mathord{\ast}}%
  \def\diamond{\mathbin{\scalebox{0.6}{$\Diamond$}}}%
  \def\bigcup{\mathop{\mathchoice{\raisebox{-0.3em}{\scalebox{1.7}{$\cup$}}}{\scalebox{1.25}{$\cup$}}{\cup}{\cup}}\displaylimits}%
  \def\bigcap{\mathop{\mathchoice{\raisebox{-0.3em}{\scalebox{1.7}{$\cap$}}}{\scalebox{1.25}{$\cap$}}{\cap}{\cap}}\displaylimits}%
  \def\lesseqqgtr{\mathrel{\lessgtr}}\def\bigoplus{\mathop{\oplus}\displaylimits}%
}
\providecommand{\ssi}{if and only if} % abréviation fréquente
\AtBeginDocument{\providecommand{\wideparen}{\overparen}} % arc de cercle

% Fonctions usuelles que les modèles utilisent sans que amsmath les fournisse
\providecommand{\arsinh}{\operatorname{arsinh}}\providecommand{\arcosh}{\operatorname{arcosh}}
\providecommand{\artanh}{\operatorname{artanh}}\providecommand{\arsech}{\operatorname{arsech}}
\providecommand{\arcsch}{\operatorname{arcsch}}\providecommand{\arcoth}{\operatorname{arcoth}}
\providecommand{\arcsinh}{\operatorname{arsinh}}\providecommand{\arccosh}{\operatorname{arcosh}}
\providecommand{\arctanh}{\operatorname{artanh}}\providecommand{\sech}{\operatorname{sech}}
\providecommand{\csch}{\operatorname{csch}}\providecommand{\arcsec}{\operatorname{arcsec}}
\providecommand{\arccsc}{\operatorname{arccsc}}\providecommand{\arccot}{\operatorname{arccot}}
\providecommand{\sgn}{\operatorname{sgn}}\providecommand{\lcm}{\operatorname{lcm}}
\providecommand{\Var}{\operatorname{Var}}\providecommand{\Cov}{\operatorname{Cov}}

%% FIGFIX-BEGIN v3 — correctifs globaux des figures (docs/onbuch-generation/tools/figfix)
\usepackage{adjustbox}\usepackage{etoolbox}
% toute figure TikZ/pgfplots plus large que la ligne est réduite à la largeur disponible
\BeforeBeginEnvironment{tikzpicture}{\begin{adjustbox}{max width=\linewidth}}
\AfterEndEnvironment{tikzpicture}{\end{adjustbox}}
% légendes pgfplots lisibles par-dessus les courbes
\pgfplotsset{every axis/.append style={legend style={fill=white,fill opacity=0.92,text opacity=1,draw=black!15,inner sep=3pt}}}
% étiquettes d'arêtes (syntaxe edge["..."]) avec fond blanc
\tikzset{every edge quotes/.append style={fill=white,inner sep=1.2pt}}
% bulles de cartes mentales : texte à la ligne dans la bulle, bulle un peu plus grande
\tikzset{every concept/.append style={text width=2.0cm,align=center,minimum size=2.3cm,inner sep=2pt}}
%% FIGFIX-END
\begin{document}

\pagegarde{STRATIGRAPHY}{Geology}{Upper Sixth}{Upper Sixth — Geology}{1}{22 periods}{This chapter equips students with the fundamental tools of stratigraphy: principles, the geologic time scale, dating techniques, contact relationships, event sequencing, and correlation methods. Mastery of these concepts is essential for interpreting sedimentary basins and succeeding in the GCE Advanced Level Geology examination.
\vspace{4pt}
\begin{multicols}{2}\small
\begin{itemize}
\item Define the scope and importance of stratigraphy in geological investigations.
\item State and apply the fundamental principles of stratigraphy (superposition, original horizontality, lateral continuity, cross‑cutting relationships, faunal succession).
\item Construct and interpret the geologic time scale, distinguishing eons, eras, periods, epochs, and ages.
\item Explain the principles, applications, and limitations of relative and absolute dating methods.
\item Calculate ages using radiometric decay equations and interpret isotopic data.
\item Identify and classify stratigraphic contacts (conformable, unconformable, disconformity, angular unconformity, nonconformity).
\end{itemize}\end{multicols}}

\begin{sommaire}
\begin{multicols}{2}
\begin{enumerate}
\item 1. Introduction \& Scope of Stratigraphy
\item 2. Fundamental Stratigraphic Principles
\item 3. The Geologic Time Scale
\item 4. Overview of Dating Methods
\item 5. Radiometric Dating in Detail
\item 6. Stratigraphic Contacts and Unconformities
\item 7. Sequencing Geologic Events
\item 8. Correlation of Stratigraphic Successions
\item 9. Integration, Revision \& Exam Practice
\item Integration activity
\item Methods and worked examples
\item Exercises
\item Solutions to the exercises
\item Summary sheet
\end{enumerate}
\end{multicols}
\end{sommaire}

\begin{objectifs}
\begin{itemize}
\item Define the scope and importance of stratigraphy in geological investigations.
\item State and apply the fundamental principles of stratigraphy (superposition, original horizontality, lateral continuity, cross‑cutting relationships, faunal succession).
\item Construct and interpret the geologic time scale, distinguishing eons, eras, periods, epochs, and ages.
\item Explain the principles, applications, and limitations of relative and absolute dating methods.
\item Calculate ages using radiometric decay equations and interpret isotopic data.
\item Identify and classify stratigraphic contacts (conformable, unconformable, disconformity, angular unconformity, nonconformity).
\item Reconstruct a sequence of geologic events from a cross‑section or outcrop description.
\item Correlate lithostratigraphic, biostratigraphic, and chronostratigraphic units across different localities.
\item Integrate multiple stratigraphic datasets to solve a realistic geological problem.
\item Apply stratigraphic reasoning to answer GCE‑style structured and essay questions.
\end{itemize}
\end{objectifs}

\begin{prerequis}
\begin{itemize}
\item Basic rock classification (igneous, sedimentary, metamorphic) from Lower Sixth.
\item Understanding of sedimentary processes and depositional environments.
\item Fundamentals of radioactive decay and half‑life from chemistry/physics.
\item Map reading and cross‑section construction skills.
\item Familiarity with fossil identification and index fossil concept.
\end{itemize}
\end{prerequis}

\newpage

\coursec{Introduction \& Scope of Stratigraphy}

\begin{experience}[The geologist's puzzle]
Imagine you are a geologist hired by a mining company in the Adamawa region to locate a new bauxite deposit. The only clues are the layered rocks exposed along a river cut. How would you read the Earth's history written in those layers to predict where the ore lies? You notice a thick reddish layer at the bottom, a grey fossil-bearing layer above it, and a pale sandy layer at the top. Which layer is the oldest? What environment did each layer form in? And in which layer should you dig for bauxite, which forms by intense tropical weathering of older rocks? Every one of these questions can be answered --- but only if you master the science that reads rock layers the way a historian reads the pages of a book. That science is \cle{stratigraphy}, and it is the foundation upon which the whole of historical geology is built.
\end{experience}

The problem posed above is the central problem of this chapter: \emph{given a stack of rock layers, how do we establish their order, their ages, their environments of formation, and their economic potential?} By the end of this section you will know what stratigraphy is, how it is subdivided, how it was born, and why it is indispensable to geologists in Cameroon and worldwide.

\courssub{What is Stratigraphy?}

When you walk along the road cuts between Bafoussam and Bamenda, or along the sea cliffs near Limbe, you see that rocks are not random: they occur in \emph{layers}, called \cle{strata} (singular: \emph{stratum}). These layers are pages of an archive. Each layer records the conditions --- river, sea, desert, volcano --- that existed when it was deposited. Stratigraphy is the science of reading that archive.

\begin{definition}[Stratigraphy]
\cle{Stratigraphy} (from Latin \emph{stratum}, ``layer'', and Greek \emph{graphein}, ``to write/describe'') is the branch of geology concerned with the \textbf{order}, \textbf{relative position}, \textbf{description}, \textbf{age determination} and \textbf{correlation} of rock layers (strata), and with the interpretation of the geological history they record.
\end{definition}

Notice the five verbs hidden in this definition: to \emph{order} strata, to \emph{position} them, to \emph{describe} them, to \emph{date} them, and to \emph{correlate} them (match them from place to place). Each of these tasks corresponds to a subdivision of stratigraphy.

\begin{definition}[Subdivisions of stratigraphy]
\begin{itemize}
\item \cle{Lithostratigraphy}: the classification of strata into units based on their \textbf{physical rock characters} (lithology, texture, colour, mineral composition), independent of age or fossils. Its basic unit is the \cle{formation}.
\item \cle{Biostratigraphy}: the classification and correlation of strata based on their \textbf{fossil content}. Its basic unit is the \cle{biozone} (or simply \emph{zone}).
\item \cle{Chronostratigraphy}: the classification of strata into units that represent rocks formed during a \textbf{specific interval of geologic time}. Its units (e.g.\ \cle{system}, \cle{series}, \cle{stage}) are bodies of rock with time significance.
\item \cle{Magnetostratigraphy}: the classification and correlation of strata based on the record of \textbf{past reversals of the Earth's magnetic field} preserved in the rocks (remanent magnetism).
\end{itemize}
\end{definition}

\begin{attention}[Do not confuse rock units with time units]
A very common GCE mistake is to mix up \textbf{lithostratigraphic} units (rock bodies: group, formation, member, bed) with \textbf{chronostratigraphic} units (rock bodies defined by time: system, series, stage) and with \textbf{geologic-time} units (pure time: period, epoch, age). Remember: a \emph{formation} is a mappable rock body; a \emph{system} is the rock record of a \emph{period}. You may walk on the Cretaceous \emph{System}, but you live in the Holocene \emph{Epoch}. A formation can be diachronous --- deposited at different times in different places --- whereas a stage boundary is, by definition, the same age everywhere.
\end{attention}

\courssub{The Hierarchy of Lithostratigraphic Units}

Because rock bodies occur at many scales, lithostratigraphy organises them into a hierarchy, from the thinnest distinguishable layer to vast complexes of related strata.

\begin{definition}[Lithostratigraphic hierarchy]
\begin{itemize}
\item \cle{Bed}: the smallest formal unit; a single distinguishable layer (e.g.\ one sandstone bed).
\item \cle{Member}: a subdivision of a formation with a distinct lithology.
\item \cle{Formation}: the \textbf{fundamental unit}; a body of rock homogeneous enough to be mapped at surface or traced in the subsurface. Formations are named after a geographic locality (e.g.\ a formation named after a town near its type section).
\item \cle{Group}: two or more adjacent, genetically related formations.
\item \cle{Supergroup}: an assemblage of related groups.
\end{itemize}
\end{definition}

\begin{popfigure}
\begin{center}
\begin{tikzpicture}[scale=1, every node/.style={font=\small}]
% nested rectangles from supergroup (outer) to bed (inner)
\draw[fill=popblueL, draw=popblue, thick] (0,0) rectangle (10,6.4);
\node[anchor=north west, text=popdark, font=\small\bfseries] at (0.15,6.3) {SUPERGROUP};
\draw[fill=white, draw=popgreen, thick] (0.5,0.5) rectangle (9.5,5.3);
\node[anchor=north west, text=popdark, font=\small\bfseries] at (0.65,5.2) {GROUP};
\draw[fill=popblueL!60, draw=poporange, thick] (1.0,1.0) rectangle (9.0,4.2);
\node[anchor=north west, text=popdark, font=\small\bfseries] at (1.15,4.1) {FORMATION (fundamental unit)};
\draw[fill=white, draw=poppurple, thick] (1.5,1.5) rectangle (8.5,3.3);
\node[anchor=north west, text=popdark, font=\small\bfseries] at (1.65,3.2) {MEMBER};
\draw[fill=poppink!30, draw=poppink, thick] (2.0,1.9) rectangle (8.0,2.6);
\node[text=popdark, font=\small\bfseries] at (5.0,2.25) {BED (smallest formal unit)};
% arrow showing increasing scale
\draw[->, very thick, popdark] (10.4,1.2) -- (10.4,5.8) node[midway, right, align=left, text width=3.2cm, font=\small] {increasing thickness and map area};
\end{tikzpicture}
\end{center}
\legende{Hierarchy of lithostratigraphic units: each unit nests inside the next larger one, from the individual bed up to the supergroup.}
\end{popfigure}

\courssub{Historical Development of Stratigraphy}

Stratigraphy did not appear overnight; it grew from simple observations into a rigorous, internationally codified science.

\textbf{Nicolas Steno (1669).} Working in Italy, the Danish physician Nicolas Steno (Niels Stensen) studied the layered rocks of Tuscany and stated the three founding principles of stratigraphy: the \cle{law of superposition} (younger layers lie on older ones), the \cle{principle of original horizontality} (layers are deposited flat), and the \cle{principle of lateral continuity} (layers extend sideways until they thin out or meet a barrier). We shall study these in detail in Section 2.

\begin{propriete}[Law of Superposition (Steno, 1669)]
In an \textbf{undeformed} sequence of sedimentary (or layered volcanic) rocks, each layer is \textbf{older than the layer above it and younger than the layer below it}: the oldest layers are at the bottom, the youngest at the top. This law gives \emph{relative} ages only, not numerical ages. (Its justification by simple reasoning about deposition is examinable: sediments can only settle on top of a pre-existing surface.)
\end{propriete}

\textbf{William Smith (1815).} The English canal engineer William ``Strata'' Smith noticed that the same rock layers, exposed in different parts of England, always contained the \emph{same assemblages of fossils}, and always occurred in the \emph{same order}. From this he formulated the \cle{principle of faunal succession} and, in 1815, published the first geological map of an entire country (England and Wales). His discovery made it possible to \emph{correlate} strata across long distances --- the birth of biostratigraphy.

\textbf{Charles Lyell (1830--1833).} In his \emph{Principles of Geology}, Lyell championed \cle{uniformitarianism}: the idea that the slow processes we observe today (erosion, sedimentation, volcanism) operated in the same way through the past, so ``the present is the key to the past''. This justified interpreting ancient strata by comparison with modern environments, and implied that Earth history was immensely long.

\textbf{Modern stratigraphic codes.} During the twentieth century, stratigraphy became codified and international. National and international \emph{stratigraphic codes} fixed the rules for naming and defining units (type sections, priority of names). Today the \cle{International Commission on Stratigraphy} (ICS), a body of the International Union of Geological Sciences (IUGS), coordinates the definition of global chronostratigraphic boundaries and publishes the \cle{International Chronostratigraphic Chart}, the official framework of geologic time used by geologists worldwide. Each boundary is fixed at a physical reference point in a rock section called a \cle{GSSP} (Global Boundary Stratotype Section and Point), informally a ``golden spike''.

\begin{popfigure}
\begin{center}
\begin{tikzpicture}[scale=1]
\draw[->, very thick, popdark] (0,0) -- (13.2,0);
% milestones
\foreach \x/\yr/\txt in {1.0/1669/{Steno: superposition, horizontality, continuity}, 4.2/1815/{Smith: faunal succession; first geological map}, 7.4/1830--1833/{Lyell: uniformitarianism}, 10.6/20th c./{Stratigraphic codes; ICS and International Chart}}{
  \draw[thick, popblue] (\x,-0.15) -- (\x,0.15);
  \node[below, font=\small\bfseries, text=popblue] at (\x,-0.2) {\yr};
  \node[above, font=\small, align=center, text width=2.9cm] at (\x,0.25) {\txt};
}
\node[font=\small\itshape, text=popmuted] at (6.6,-1.15) {Key milestones in the birth of stratigraphy};
\end{tikzpicture}
\end{center}
\legende{Timeline of the major historical milestones in the development of stratigraphy.}
\end{popfigure}

\begin{savaistu}[Stratigraphy beneath our feet in Cameroon]
Cameroon's sedimentary basins --- the Douala/Kribi--Campo basin in the south and the Logone Birni basin in the far north --- were unravelled by stratigraphy. Oil exploration in the Rio del Rey basin depends on correlating sandstone reservoir formations and shale seals from well to well, exactly the lithostratigraphic and biostratigraphic reasoning you are learning here. Even the volcanic rocks of the Cameroon Volcanic Line are dated and ordered using stratigraphic principles combined with radiometric dating.
\end{savaistu}

\courssub{Why Stratigraphy Matters}

Stratigraphy is not an academic luxury; it is a working tool with direct economic and scientific value.

\begin{itemize}
\item \textbf{Resource exploration.} Oil, gas, coal, bauxite, iron, limestone and groundwater all occur in specific layers. Stratigraphy tells the explorer \emph{which} layer to target and \emph{where} that layer continues underground. Our Adamawa bauxite, for example, forms as a weathering horizon on older rocks: finding the right ancient land surface (a stratigraphic contact) is the key to the ore.
\item \textbf{Basin analysis.} Sedimentary basins (like the Douala basin) fill with strata that record subsidence, sea-level changes and sediment supply. Reading this fill lets geologists reconstruct the basin's history and predict where reservoir rocks, source rocks and seals lie --- the essence of petroleum geology.
\item \textbf{Palaeoenvironmental reconstruction.} Fossils and sedimentary structures in strata reveal past climates, seas, rivers and deserts. Stratigraphy turns a cliff face into a history book of environments, and provides the framework for dating past events --- including climate changes relevant to our own future.
\end{itemize}

\courssub{Reading a Stratigraphic Column}

The standard document of stratigraphy is the \cle{stratigraphic column}: a vertical drawing of the succession of layers, oldest at the bottom, with thicknesses to scale, standard symbols for rock types, colours, and fossil content.

\begin{methode}[How to read a stratigraphic column legend]
\begin{enumerate}
\item Start at the \textbf{bottom}: by superposition, this is the oldest unit.
\item Read the \textbf{lithology symbols}: brick pattern for limestone, dots for sandstone, dashes for shale, triangles or ``v'' shapes for volcanic rocks. Check the legend first --- symbols vary between authors.
\item Read the \textbf{thickness scale}: the vertical axis is drawn to scale (e.g.\ 1 cm = 10 m), so you can measure true bed thicknesses.
\item Note the \textbf{colours}: colours usually follow international convention (e.g.\ blue tones for marine limestones, yellow for sandstones) but always confirm with the legend.
\item Identify \textbf{fossils and ages} written beside the column: these give the biostratigraphic and chronostratigraphic framework.
\item Mark the \textbf{contacts} between units: sharp, gradational, or wavy (unconformity) lines --- they will be studied in Section 6.
\end{enumerate}
\end{methode}

\begin{exoresolu}[Reading a river-cut section]
Statement. In a river cut in the Adamawa region, a geologist observes from bottom to top: (A) 5 m of grey shale with marine shells; (B) 3 m of cross-bedded sandstone; (C) 2 m of red lateritic horizon rich in aluminium oxides. (1) Place the units in chronological order. (2) Which subdivision(s) of stratigraphy are used to order the units, and which to interpret their environments? (3) Where should the mining company sample for bauxite?
\tcblower
\textbf{Solution.}
\begin{enumerate}
\item By the \textbf{law of superposition}, the undeformed sequence is oldest at the bottom: A (shale) $\rightarrow$ B (sandstone) $\rightarrow$ C (laterite). Chronological order of formation: A, then B, then C.
\item \textbf{Lithostratigraphy} is used to describe and order the rock units (shale, sandstone, laterite as distinct beds/formations). The marine shells allow \textbf{biostratigraphy} (fossil-based dating and correlation). Environmental interpretation (marine shale, fluvial or coastal sandstone, tropical weathering surface) uses sedimentology together with the stratigraphic framework.
\item Bauxite forms by intense tropical weathering, so the target is unit \textbf{C}, the lateritic horizon --- and, crucially, the company should trace this horizon \emph{laterally} (principle of lateral continuity) to find where it thickens into an economic deposit.
\end{enumerate}
\end{exoresolu}

\begin{aretenir}[Key points of Section 1]
\begin{itemize}
\item Stratigraphy is the science of the order, position, description, dating and correlation of rock layers.
\item Its main subdivisions are \textbf{lithostratigraphy} (rock character), \textbf{biostratigraphy} (fossils), \textbf{chronostratigraphy} (time-rock units) and \textbf{magnetostratigraphy} (magnetic reversals).
\item Lithostratigraphic hierarchy: bed $\rightarrow$ member $\rightarrow$ \textbf{formation} (fundamental unit) $\rightarrow$ group $\rightarrow$ supergroup.
\item Founders: \textbf{Steno} (superposition, 1669), \textbf{Smith} (faunal succession, 1815), \textbf{Lyell} (uniformitarianism, 1830s); today the \textbf{ICS} maintains the International Chronostratigraphic Chart.
\item Stratigraphy underpins resource exploration, basin analysis and palaeoenvironmental reconstruction.
\item Never confuse lithostratigraphic units (rocks) with chronostratigraphic units (time-rock) or geologic-time units (pure time).
\end{itemize}
\end{aretenir}

\begin{exercice}[Check your understanding]
\begin{enumerate}
\item Define stratigraphy and name its four subdivisions, giving the basis of classification of each.
\item A student writes: ``The Cretaceous formation contains many fossils.'' Identify and correct the terminological error.
\item Explain why the formation is called the \emph{fundamental} lithostratigraphic unit.
\item State the law of superposition and give the condition under which it may fail.
\end{enumerate}
\end{exercice}

\begin{corrige}[Exercise 1]
\begin{enumerate}
\item Stratigraphy is the branch of geology dealing with the order, relative position, description, dating and correlation of strata. Subdivisions: lithostratigraphy (rock character), biostratigraphy (fossil content), chronostratigraphy (time interval of formation), magnetostratigraphy (magnetic reversal record).
\item ``Cretaceous'' is a \emph{chronostratigraphic/time} term (System/Period), not a lithostratigraphic one. Correct: ``The Cretaceous \emph{System} contains many fossils'' or ``The formation, dated as Cretaceous, contains many fossils''.
\item Because it is the smallest unit that is homogeneous and distinct enough to be \emph{mapped} and traced in the field; members and beds are its subdivisions, groups its assemblages.
\item In an undeformed sequence, each bed is younger than the one below and older than the one above. It fails where strata have been \emph{overturned} by folding or thrusting; younging indicators (graded bedding, cross-bedding, mud cracks) must then be used.
\end{enumerate}
\end{corrige}

\coursec{Fundamental Stratigraphic Principles}

In the previous section we defined stratigraphy as the science of layered rocks and saw that its ultimate goal is to reconstruct the history of the Earth from the rock record. But how can a geologist, standing in front of a road cut near Bafoussam or a cliff on the flanks of Mount Cameroon, \emph{read} that history? The answer lies in a small set of remarkably powerful principles, most of them established between the 17th and 19th centuries by pioneers such as Nicolaus Steno, James Hutton and William Smith. These principles are the grammar of stratigraphy: once you master them, any outcrop becomes a readable text. In this section we state each principle rigorously, justify it, illustrate it, and learn to apply it.

\courssub{Principle of Original Horizontality}

\textbf{Intuition.} Pour a bucket of sand into a basin of water and watch it settle: it forms a flat, horizontal layer. Sediments deposited under gravity behave the same way, whether in the ocean, a lake such as Lake Nyos, or a river floodplain in the Benue trough.

\begin{definition}[Principle of Original Horizontality (Steno, 1669)]
Layers of sediment are originally deposited in a horizontal or near-horizontal position, because deposition is controlled by gravity. Consequently, any sedimentary succession that is today tilted, folded or overturned must have been deformed \emph{after} deposition.
\end{definition}

This simple idea has an enormous consequence: \textbf{tilted strata are evidence of a tectonic event younger than the strata themselves}. When you see the steeply dipping sandstones around the Benue trough, you immediately know that compressional tectonics affected the region after the sediments were laid down.

\begin{attention}[Exceptions to strict horizontality]
Two situations produce beds that were \emph{never} perfectly horizontal:
\begin{itemize}
\item \cle{Depositional (primary) dip}: cross-bedding in dunes and deltas, and the slopes of reef mounds or volcano flanks, can reach $30^\circ$ or more. These dips are \emph{primary} and do not indicate tectonism.
\item \cle{Heavily deformed terranes}: in strongly folded or thrusted regions, beds may be vertical or overturned. Never assume horizontality there without a structural analysis (way-up criteria such as graded bedding, mud cracks, ripple marks).
\end{itemize}
\end{attention}

\courssub{Principle of Lateral Continuity and Facies Changes}

\textbf{Intuition.} Imagine a layer of mud deposited on the floor of a wide lagoon. If a valley is later eroded through the lagoon, the mud layer appears on both valley walls. You do not need to see the layer continuously to know the two outcrops were once connected.

\begin{definition}[Principle of Lateral Continuity]
A sedimentary layer extends laterally in all directions until it thins out (pinches out) at the edge of the depositional basin, or until it passes laterally into a different sediment type (a \cle{facies change}) where the depositional environment changed.
\end{definition}

The notion of \cle{facies} is essential: at any one time, different sediments accumulate side by side — sand near a shore, mud offshore, limestone on a clear-water platform. A single time plane therefore cuts across several rock types, and a single rock type can be of different ages in different places (the facies is said to be \emph{diachronous} when it migrates with time, e.g.\ during a marine transgression).

\begin{attention}[Common mistake]
Do not correlate two beds simply because they have the same lithology (e.g.\ two sandstones). A beach sand is young near the present shore and older inland after a transgression. Lithology alone is not a time marker.
\end{attention}

\courssub{Principle of Cross-Cutting Relationships}

\textbf{Intuition.} To slice a cake, the cake must already exist. Likewise, a feature that cuts across rocks must be younger than the rocks it cuts.

\begin{propriete}[Principle of Cross-Cutting Relationships]
Any geological feature (igneous intrusion, fault, mineral vein, erosional surface) that cuts across another rock body or structure is \emph{younger} than everything it cuts. Conversely, a rock body is older than any feature that cuts it. The proof is direct observation of geometric relations at the contact; it is examinable as a reasoning exercise on a cross-section.
\end{propriete}

Typical applications:
\begin{itemize}
\item A basaltic \cle{dyke} cutting gneisses of the Congo craton is younger than the gneisses.
\item A \cle{fault} displacing sedimentary beds is younger than the youngest bed it offsets.
\item A quartz \cle{vein} cross-cutting a granite post-dates the granite.
\end{itemize}

If a dyke is truncated by a fault, the dyke is older than the fault: the principle applies recursively and allows a complete relative chronology to be built.

\courssub{Principle of Faunal Succession and Index Fossils}

\textbf{Intuition.} William Smith, surveying canal cuttings in England around 1800, noticed that each layer contained a characteristic assemblage of fossils, and that these assemblages always succeeded one another in the same order, wherever he looked. Life evolves and never repeats itself; therefore fossil content is a clock.

\begin{definition}[Principle of Faunal Succession]
Fossil organisms succeed one another through time in a definite, irreversible order. Each interval of geological time is characterised by a distinctive fossil assemblage, so strata containing the same assemblage can be recognised and correlated from one region to another.
\end{definition}

\begin{definition}[Index fossil]
An \cle{index fossil} is a fossil that is ideal for dating and correlating rocks because it combines three qualities:
\begin{enumerate}
\item \textbf{Wide geographic distribution} — it can be found in many regions;
\item \textbf{Narrow time range} — the species lived for a short interval only;
\item \textbf{Easy identification} — it is abundant, well preserved and distinctive.
\end{enumerate}
Classic examples: ammonites (Mesozoic), trilobites (Palaeozoic), graptolites (Ordovician--Silurian), foraminifera (Cenozoic).
\end{definition}

\courssub{Principle of Inclusions and Baked Contacts}

\textbf{Intuition.} If you find pebbles of granite inside a conglomerate, the granite must have existed, been eroded and fragmented \emph{before} the conglomerate formed. Fragments are always older than the rock that contains them.

\begin{definition}[Principle of Inclusions]
Fragments (\cle{inclusions}) of one rock contained within another rock are older than the enclosing rock. This applies to pebbles in conglomerates, xenoliths in granites, and clasts in volcanic breccias.
\end{definition}

A related field criterion is the \cle{baked contact}: where hot magma intrudes cold country rock, the surrounding rock is thermally metamorphosed (``baked'') over a contact aureole. A baked margin on the \emph{country rock} proves the intrusion is younger; conversely, if pebbles of an already-baked rock occur in an overlying conglomerate, the intrusion is older than the conglomerate. These two criteria together let you distinguish an intrusive contact from a depositional one.

\courssub{Synthetic Illustration: the Five Principles on One Cross-Section}

\begin{popfigure}
\begin{tikzpicture}[scale=1.0, font=\small]
% horizontal beds
\fill[popblueL] (0,0) rectangle (4,0.8);
\fill[popgreenL] (0,0.8) rectangle (4,1.6);
\fill[popblueL] (0,1.6) rectangle (4,2.4);
\draw (0,0) rectangle (4,2.4);
\node[align=center] at (2,-0.5) {(a) Horizontality};
\node at (2,0.4) {bed 1};
\node at (2,1.2) {bed 2};
\node at (2,2.0) {bed 3};
% lateral continuity / pinch-out
\fill[popgreenL] (5,0) rectangle (9,0.7);
\fill[popblueL] (5,0.7) -- (9,0.7) -- (7.6,1.5) -- (5,1.5) -- cycle;
\fill[popgreenL] (5,1.5) rectangle (9,2.3);
\draw (5,0) rectangle (9,2.3);
\node[align=center] at (7,-0.5) {(b) Pinch-out};
\node at (6.2,1.1) {sand};
% cross-cutting dyke
\fill[popblueL] (10,0) rectangle (14,2.4);
\fill[poporange] (11.4,0) -- (11.9,0) -- (12.6,2.4) -- (12.1,2.4) -- cycle;
\draw (10,0) rectangle (14,2.4);
\node[align=center] at (12,-0.5) {(c) Dyke (younger)};
\node at (12.4,1.2) {\textcolor{white}{dyke}};
% faunal succession
\fill[popgreenL] (15,0) rectangle (19,0.8);
\fill[popblueL] (15,0.8) rectangle (19,1.6);
\fill[popgreenL] (15,1.6) rectangle (19,2.4);
\draw (15,0) rectangle (19,2.4);
\node at (17,0.4) {trilobites};
\node at (17,1.2) {ammonites};
\node at (17,2.0) {foraminifera};
\node[align=center] at (17,-0.5) {(d) Faunal succession};
\end{tikzpicture}
\legende{The fundamental principles at a glance: (a) originally horizontal beds; (b) lateral continuity ending in a pinch-out; (c) a cross-cutting dyke younger than the host beds; (d) an ordered succession of fossil assemblages.}
\end{popfigure}

\courssub{Worked Example: Interpreting an Outcrop Sketch}

\begin{exoresolu}[Reading a simple outcrop]
A road cut shows, from bottom to top: (i) folded schists; (ii) a horizontal conglomerate containing pebbles of schist and of a granite; (iii) horizontal sandstone beds with ammonites; (iv) a basalt dyke cutting the conglomerate and the sandstone but stopping at the top of the sandstone; (v) a soil cover. Establish the relative chronology.
\tcblower
\textbf{Solution.} Apply the principles one by one:
\begin{enumerate}
\item \textbf{Schists}: the oldest unit present. Their folding (deformation) is younger than the schists but older than everything that lies undeformed above (original horizontality).
\item \textbf{Conglomerate}: rests on the eroded schists (an unconformity). By the \emph{principle of inclusions}, the schist and granite pebbles are older than the conglomerate; the granite therefore existed and was exposed at the surface before conglomerate deposition.
\item \textbf{Sandstone with ammonites}: younger than the conglomerate (superposition) and Mesozoic in character (faunal succession, ammonites as index fossils).
\item \textbf{Basalt dyke}: by \emph{cross-cutting relationships}, it is younger than the conglomerate and the sandstone. Since it does not cut the soil, it is older than the soil.
\item \textbf{Soil}: youngest feature.
\end{enumerate}
\textbf{Chronology (oldest $\rightarrow$ youngest):} schists $\rightarrow$ folding and erosion $\rightarrow$ granite (existed before conglomerate) $\rightarrow$ conglomerate $\rightarrow$ sandstone $\rightarrow$ basalt dyke $\rightarrow$ soil.
\end{exoresolu}

\courssub{Method: a Checklist for Any Unknown Outcrop}

\begin{methode}[Applying the five principles systematically]
\begin{enumerate}
\item \textbf{List the units and features}: beds, intrusions, faults, veins, erosion surfaces.
\item \textbf{Check way-up and horizontality}: are beds horizontal, tilted or overturned? Look for graded bedding or mud cracks.
\item \textbf{Apply superposition and lateral continuity}: order the undisturbed beds from oldest (bottom) to youngest (top); trace each bed laterally and note pinch-outs and facies changes.
\item \textbf{Apply inclusions}: any fragments (pebbles, xenoliths) are older than their host rock.
\item \textbf{Apply cross-cutting relationships}: every dyke, fault or vein is younger than what it cuts; rank several cutting features by who cuts whom.
\item \textbf{Use fossils}: identify index fossils to assign relative ages and to correlate with other sections.
\item \textbf{Write the full chronology} from oldest to youngest, stating the principle used at each step.
\end{enumerate}
\end{methode}

\begin{attention}[Horizontality trap]
In heavily deformed terranes (e.g.\ the Pan-African belts of northern Cameroon), beds may be vertical or upside down. Never apply superposition blindly: first establish the younging direction with sedimentary way-up criteria. Assuming horizontality without structural analysis is one of the most frequent causes of inverted chronologies in the GCE practical paper.
\end{attention}

\begin{aretenir}[Key points]
\begin{itemize}
\item Sediments are deposited horizontally; tilted beds record later tectonism.
\item Beds extend laterally until they pinch out or change facies; lithology alone is not a time marker.
\item Whatever cuts is younger; whatever is included is older.
\item Fossil assemblages succeed in a fixed order; index fossils (wide distribution, short range, easy identification) date and correlate strata.
\item Baked contacts prove an intrusion post-dates its country rock.
\item Combine all five principles, in the checklist order, to decode any outcrop.
\end{itemize}
\end{aretenir}

\begin{exercice}[Outcrop chronology]
A cliff section shows limestone beds (with trilobites) cut by a dolerite dyke; the dyke and the limestone are both truncated by an erosion surface overlain by a conglomerate containing dolerite pebbles; a fault cuts the limestone and the dyke but not the conglomerate. Establish the complete relative chronology, justifying each step with the appropriate principle.
\end{exercice}

\begin{corrige}[Exercise 1]
Oldest $\rightarrow$ youngest: \textbf{limestone} (base of the exposed succession; trilobites indicate Palaeozoic) $\rightarrow$ \textbf{dolerite dyke} (cross-cuts the limestone, so younger) $\rightarrow$ \textbf{fault} (cuts limestone and dyke, so younger than both) $\rightarrow$ \textbf{erosion surface} (truncates dyke and fault, so younger still) $\rightarrow$ \textbf{conglomerate} (lies above the erosion surface; its dolerite pebbles confirm, by the principle of inclusions, that the dyke existed and was eroded before conglomerate deposition). Note the double proof that the dyke is older than the conglomerate: the dyke is truncated by the erosion surface \emph{and} its pebbles occur in the conglomerate.
\end{corrige}

\coursec{The Geologic Time Scale}

In the previous section we established the \emph{fundamental principles} of stratigraphy --- superposition, horizontality, lateral continuity, cross-cutting relationships, faunal succession. These principles allow us to place rocks and events in \emph{relative} order: bed A is older than bed B, fault F is younger than both. But a relative order is not enough. To communicate globally, geologists need a \emph{standard calendar of Earth history}: a single, internationally agreed framework in which every slice of time has a name, a rank and a numerical age. That framework is the \cle{Geologic Time Scale} (GTS). It is the single most frequently examined diagram in the whole GCE Advanced Level Geology paper, and mastering it is non-negotiable.

\courssub{From Relative Order to a Global Calendar}

Imagine the history of Cameroon written without any dates: ``the German period came before the French mandate, which came before independence.'' The order is correct, but you could not say \emph{how long} each episode lasted or \emph{when exactly} it began. Stratigraphy faced exactly this problem in the nineteenth century. Geologists could stack strata in order, but they needed named compartments of time.

The solution was to divide the rock record into packages bounded by \emph{worldwide recognisable events} --- typically mass extinctions or abrupt faunal changes --- and to give each package a name. Later (twentieth century), radiometric dating attached numerical ages to the boundaries. The modern GTS is therefore a \emph{hybrid}: its compartments are defined by rocks and fossils, and its boundary ages are calibrated by isotopic dating.

\begin{aretenir}[Two linked hierarchies]
There are always \textbf{two parallel scales}:
\begin{itemize}
\item \cle{Chronostratigraphic} (bodies of rock): Eonothem $\rightarrow$ Erathem $\rightarrow$ \textbf{System} $\rightarrow$ Series $\rightarrow$ Stage.
\item \cle{Geochronologic} (intervals of time): Eon $\rightarrow$ Era $\rightarrow$ \textbf{Period} $\rightarrow$ Epoch $\rightarrow$ Age.
\end{itemize}
Rocks belong to a \emph{System}; time belongs to a \emph{Period}. We say ``the Jurassic \textbf{Period}'' (time) but ``the Jurassic \textbf{System}'' (the rocks deposited during it).
\end{aretenir}

\begin{attention}[Period, System, Series --- do not mix the ranks]
A classic GCE trap: candidates write ``the Cretaceous Series'' when they mean the whole Cretaceous time interval. \textbf{Series} is the chronostratigraphic rank \emph{below} System (equivalent to an \emph{Epoch} in time), e.g.\ the Upper Cretaceous Series. \textbf{Period} (time) and \textbf{System} (rock) are equivalent in rank but used in different contexts. Rule of thumb: \emph{time words} (Period, Epoch, Age) answer ``when?''; \emph{rock words} (System, Series, Stage) answer ``which rocks?''.
\end{attention}

\courssub{The Structure of the Scale: Eons to Ages}

The scale is a nested hierarchy, like a set of boxes within boxes:

\begin{center}
\begin{tabularx}{\linewidth}{|l|X|l|}
\hline
\rowcolor{popblueL} \textbf{Rank (time)} & \textbf{Typical duration} & \textbf{Example} \\
\hline
Eon & Hundreds of Ma to $\sim$2 Ga & Phanerozoic \\
\hline
Era & Tens to hundreds of Ma & Mesozoic \\
\hline
Period & Tens of Ma & Jurassic \\
\hline
Epoch & A few to tens of Ma & Pleistocene \\
\hline
Age & $\sim$1 Ma or less & Calabrian \\
\hline
\end{tabularx}
\end{center}

\begin{exemplebox}[Zooming into one address]
Just as a letter addressed to a student in Buea reads \emph{Cameroon $\rightarrow$ South-West Region $\rightarrow$ Fako Division $\rightarrow$ Buea}, a fossil ammonite might be ``addressed'' in time as: \textbf{Phanerozoic Eon $\rightarrow$ Mesozoic Era $\rightarrow$ Cretaceous Period $\rightarrow$ Late Cretaceous Epoch $\rightarrow$ Campanian Age}. Each rank narrows the interval. For the GCE you must know the scale down to \textbf{Period} (with Cenozoic Epochs as a bonus).
\end{exemplebox}

\courssub{The Precambrian: Hadean, Archean, Proterozoic}

Nearly \textbf{88\%} of Earth history precedes the Cambrian. This vast interval, formerly lumped as ``the Precambrian'', is now formally divided:

\begin{itemize}
\item \cle{Hadean} Eon (4\,567--4\,000 Ma): from the formation of the Earth (about 4.57 Ga) to the oldest preserved rocks. Intense meteorite bombardment, magma ocean, formation of the first crust and oceans. No formal rock record --- hence the name, from Hades.
\item \cle{Archean} Eon (4\,000--2\,500 Ma): oldest continental crust (cratons), first undisputed life (prokaryotes, stromatolites). The Congo craton, which underlies much of southern Cameroon, contains Archean rocks.
\item \cle{Proterozoic} Eon (2\,500--541 Ma): oxygenation of the atmosphere, first eukaryotes, then multicellular (Ediacaran) life; ends with the appearance of abundant shelly fossils that marks the Cambrian. The Proterozoic is subdivided into three eras (Paleoproterozoic, Mesoproterozoic, Neoproterozoic) but knowledge of these subdivisions is not required for the GCE examination.
\end{itemize}

\begin{propriete}[The Great Oxidation Event (GOE)]
Around \textbf{2.4--2.3 Ga} (early Proterozoic), photosynthetic cyanobacteria raised atmospheric \ce{O2} dramatically. Geological evidence: \emph{banded iron formations} (BIFs) stop forming in abundance as dissolved iron was oxidised and precipitated, and \emph{red beds} (continental sediments stained by iron oxides) appear. The GOE made aerobic life --- and ultimately animals --- possible. (Its causes and timing are discussed qualitatively at GCE level; no numerical derivation is examinable.)
\end{propriete}

\begin{attention}[``Precambrian'' is informal]
``Precambrian'' is a convenient informal term for everything before 541 Ma, but it is \textbf{not} a formal unit. In an exam answer, name the formal eons: Hadean, Archean, Proterozoic. Also note the Archean--Proterozoic boundary is fixed \emph{chronometrically} at 2\,500 Ma, not by a fossil event.
\end{attention}

\courssub{The Phanerozoic: Eras, Periods and Numerical Ages}

The \cle{Phanerozoic} Eon (``visible life'', 541 Ma to present) is divided into three eras, separated by the two greatest mass extinctions:

\begin{itemize}
\item \cle{Paleozoic} (541--252 Ma): ``ancient life''. Cambrian explosion of shelly faunas; Ordovician radiation; Silurian colonisation of land by plants; Devonian ``age of fishes'' and first forests; Carboniferous coal swamps; Permian supercontinent Pangaea. Ends with the \textbf{Permian--Triassic extinction} ($\sim$252 Ma), the largest known, wiping out most marine species.
\item \cle{Mesozoic} (252--66 Ma): ``middle life'', the age of reptiles. Triassic recovery and first dinosaurs and mammals; Jurassic giant dinosaurs and first birds; Cretaceous flowering plants and chalk seas. Ends with the \textbf{Cretaceous--Paleogene (K/Pg) extinction} at \textbf{66.0 Ma} --- an asteroid impact plus massive volcanism --- which eliminated non-avian dinosaurs and ammonites.
\item \cle{Cenozoic} (66 Ma--present): ``recent life'', the age of mammals and birds. Divided into Paleogene, Neogene and Quaternary periods, with epochs (Paleocene, Eocene, Oligocene, Miocene, Pliocene, Pleistocene, Holocene). The Holocene began $\sim$11\,700 years ago after the last ice age.
\end{itemize}

\begin{center}
\begin{tabularx}{\linewidth}{|l|l|X|}
\hline
\rowcolor{popblueL} \textbf{Period} & \textbf{Base age (Ma)} & \textbf{Key milestone} \\
\hline
Quaternary & 2.58 & Ice ages; humans \\
\hline
Neogene & 23.0 & Grasslands, hominids \\
\hline
Paleogene & 66.0 & Mammal radiation after K/Pg \\
\hline
Cretaceous & 145 & Flowering plants; ends 66.0 Ma \\
\hline
Jurassic & 201 & Dinosaurs dominant; first birds \\
\hline
Triassic & 252 & Recovery; first dinosaurs, mammals \\
\hline
Permian & 299 & Pangaea; ends in greatest extinction \\
\hline
Carboniferous & 359 & Coal forests, first reptiles \\
\hline
Devonian & 419 & Fishes, first forests \& tetrapods \\
\hline
Silurian & 444 & Land plants, jawed fish \\
\hline
Ordovician & 485 & Marine radiation \\
\hline
Cambrian & 541 & Explosion of shelly life \\
\hline
\end{tabularx}
\end{center}

\begin{methode}[Mnemonic for the order of periods]
Learn the Phanerozoic periods in order (oldest $\rightarrow$ youngest) with:
\begin{center}
\emph{``Camels Often Sit Down Carefully; Perhaps Their Joints Creak --- Pretty Nasty Quakes''}
\end{center}
\textbf{C}ambrian, \textbf{O}rdovician, \textbf{S}ilurian, \textbf{D}evonian, \textbf{C}arboniferous, \textbf{P}ermian, \textbf{T}riassic, \textbf{J}urassic, \textbf{C}retaceous, \textbf{P}aleogene, \textbf{N}eogene, \textbf{Q}uaternary. Recite it until automatic --- a mis-ordered time scale loses marks instantly.
\end{methode}

\begin{popfigure}
\begin{center}
\begin{tikzpicture}
\begin{axis}[
  xbar, xmin=0, xmax=560,
  width=11cm, height=9.5cm,
  xlabel={Millions of years before present (Ma)},
  symbolic y coords={Quaternary,Neogene,Paleogene,Cretaceous,Jurassic,Triassic,Permian,Carboniferous,Devonian,Silurian,Ordovician,Cambrian},
  ytick=data, y tick label style={font=\small},
  nodes near coords, nodes near coords style={font=\scriptsize},
  bar width=6pt, xmajorgrids=true, grid style={popline!40},
]
\addplot[fill=popblue, draw=popdark] coordinates {
 (541,Cambrian) (485,Ordovician) (444,Silurian) (419,Devonian)
 (359,Carboniferous) (299,Permian) (252,Triassic) (201,Jurassic)
 (145,Cretaceous) (66,Paleogene) (23,Neogene) (2.58,Quaternary)};
\end{axis}
\end{tikzpicture}
\end{center}
\legende{Phanerozoic time scale drawn to scale: each bar runs from the present back to the base age of the period. Note how the Cretaceous ($\sim$79 Myr) dwarfs the Quaternary ($\sim$2.6 Myr) --- a common exam point.}
\end{popfigure}

\begin{attention}[Reading the chart correctly]
Two frequent errors: (1) assuming all periods have equal duration --- they do not; the Cretaceous lasted about 79 Myr, the Silurian only about 25 Myr; (2) forgetting that \textbf{larger Ma numbers are older}. The base of the Cambrian (541 Ma) is \emph{older} than the base of the Triassic (252 Ma).
\end{attention}

\courssub{GSSPs: The ``Golden Spikes''}

How is a boundary made official? Not by vote alone, but by choosing one \emph{physical reference point} in one rock section on Earth.

\begin{definition}[GSSP -- Global Boundary Stratotype Section and Point]
A \cle{GSSP} is an internationally agreed reference point within a specific, continuously deposited rock section that \emph{defines} the lower boundary of a chronostratigraphic unit (Stage, Series, System). It is popularly called a \cle{golden spike} because geologists symbolically drive a marker into the outcrop. The boundary is defined by a \emph{primary marker} --- usually the first appearance of a fossil species, a magnetic reversal, or a geochemical signal --- and all sections worldwide are then correlated to it.
\end{definition}

\begin{exemplebox}[Famous GSSPs]
\begin{itemize}
\item \textbf{Base of the Cambrian (541 Ma)}: Fortune Head, Newfoundland (Canada); marked by the first appearance of the trace fossil \emph{Treptichnus pedum}.
\item \textbf{Cretaceous--Paleogene boundary (66.0 Ma)}: El Kef, Tunisia; marked by a global \emph{iridium anomaly} and shocked quartz from the Chicxulub impact.
\item \textbf{Base of the Quaternary (2.58 Ma)}: Monte San Nicola, Sicily; marked by a marine climate (glacial) signal.
\end{itemize}
\end{exemplebox}

\begin{popfigure}
\begin{center}
\begin{tikzpicture}[scale=1.0]
% rock section
\fill[popgold!35] (0,0) rectangle (6,2.2);
\fill[popblue!30] (0,2.2) rectangle (6,4.6);
\draw[thick,popdark] (0,0) rectangle (6,4.6);
\foreach \y in {0.5,1.0,1.5} {\draw[popdark!60] (0,\y)--(6,\y);}
\foreach \y in {2.7,3.2,3.7,4.2} {\draw[popdark!60] (0,\y)--(6,\y);}
% boundary
\draw[very thick,poporange] (0,2.2)--(6,2.2);
\node[anchor=west,popdark,font=\small] at (6.1,2.2) {GSSP};
% golden spike
\draw[fill=popgold,draw=popdark] (3,2.2) circle (0.09);
\draw[popdark,thick] (3,2.2)--(3,1.55);
\draw[fill=popgold,draw=popdark] (2.85,2.2) rectangle (3.15,2.42);
\node[anchor=south,popdark,font=\small,align=center,text width=3cm] at (3,2.55) {``golden spike''};
% labels
\node[popdark,font=\small] at (3,1.1) {Older System (below)};
\node[popdark,font=\small] at (3,3.4) {Younger System (above)};
\node[anchor=east,popdark,font=\scriptsize,align=right,text width=2.6cm] at (-0.15,2.2) {primary marker:\\ first appearance of index fossil};
\draw[popdark,->] (-0.1,2.2)--(0.4,2.2);
\end{tikzpicture}
\end{center}
\legende{A GSSP ``golden spike'': the boundary between two Systems is fixed at a precise point in a continuous section, defined by a primary marker (e.g.\ a fossil first appearance).}
\end{popfigure}

\begin{propriete}[The time scale is a living document]
The GTS is \emph{continuously refined} by the International Commission on Stratigraphy (ICS) as dating improves: boundary ages are revised (e.g.\ the K/Pg boundary is now placed at \textbf{66.0 Ma}, not 65.5 Ma as in older books). \textbf{For exams, always use the latest ICS values given in your syllabus or question paper}, and quote ages to the precision taught (whole Ma for the GCE). This property is not examinable as a proof, but quoting an outdated age can cost marks.
\end{propriete}

\begin{savaistu}[Cameroon on the time scale]
Cameroon displays rocks from across the scale: Archean--Proterozoic basement of the Congo craton in the south, Cretaceous sediments of the Benue trough in the north, and Cenozoic volcanics of the \textbf{Cameroon Volcanic Line}, including Mount Cameroon --- whose eruptions (e.g.\ 1999, 2000) are events of the \emph{Holocene}, the youngest Age of the youngest Epoch. Lake Nyos, site of the 1986 gas disaster, occupies a young volcanic maar on this same line.
\end{savaistu}

\courssub{Reading and Constructing a Time-Scale Chart for the Exam}

\begin{methode}[Constructing a simplified time scale]
\begin{enumerate}
\item Draw a vertical bar; mark the \textbf{present (0 Ma) at the top} and \textbf{4\,600 Ma at the bottom}.
\item Mark the three great divides: \textbf{2\,500 Ma} (Archean/Proterozoic), \textbf{541 Ma} (Precambrian/Cambrian = base of Phanerozoic), \textbf{252 Ma} (Paleozoic/Mesozoic), \textbf{66 Ma} (Mesozoic/Cenozoic).
\item Subdivide the Phanerozoic into the 12 periods using the mnemonic; add base ages from the table above.
\item Check proportions: the Precambrian must occupy about \textbf{88\%} of the bar --- if your Phanerozoic looks half the chart, it is wrong.
\end{enumerate}
\end{methode}

\begin{exoresolu}[Placing events on the time scale]
Statement: A trilobite-bearing shale in the Benue region is dated at 470 Ma. (a) To which eon, era and period does it belong? (b) A basalt of the Cameroon Volcanic Line is dated at 30 Ma. Which era and period? (c) Which rock is older, and by how much?
\tcblower
\textbf{Solution.}
\begin{itemize}
\item (a) 470 Ma lies between 485 and 444 Ma $\Rightarrow$ \textbf{Ordovician Period}, \textbf{Paleozoic Era}, \textbf{Phanerozoic Eon}. (Consistent with trilobites, typical Paleozoic fossils.)
\item (b) 30 Ma lies between 66 and 23 Ma $\Rightarrow$ \textbf{Paleogene Period}, \textbf{Cenozoic Era}.
\item (c) The shale is older: $470 - 30 = 440$ Myr older.
\end{itemize}
Method: always locate the age \emph{between two boundary values} from the table, then read off period, era, eon.
\end{exoresolu}

\begin{exercice}[Time-scale practice]
(a) Give the full ``time address'' (eon, era, period) of a dinosaur bone dated at 100 Ma and of a stromatolite dated at 2\,000 Ma. (b) State the numerical ages of the two great mass-extinction boundaries. (c) Why is the base of the Cambrian defined by a trace fossil rather than a numerical age alone?
\end{exercice}

\begin{corrige}[Exercise 1]
(a) 100 Ma: between 145 and 66 Ma $\Rightarrow$ Cretaceous Period, Mesozoic Era, Phanerozoic Eon. 2\,000 Ma: between 2\,500 and 541 Ma $\Rightarrow$ Proterozoic Eon (subdivided into Paleoproterozoic, Mesoproterozoic, Neoproterozoic eras, but not required at this level). (b) Permian--Triassic at 252 Ma; Cretaceous--Paleogene at 66.0 Ma. (c) Because the GSSP concept defines boundaries by a \emph{physical, correlatable event} in a rock section (first appearance of \emph{Treptichnus pedum}); the numerical age (541 Ma) is a secondary, revisable calibration obtained by radiometric dating, whereas the golden spike is permanent.
\end{corrige}

\begin{aretenir}[Key points of Section 3]
\begin{itemize}
\item Hierarchy: Eon $\rightarrow$ Era $\rightarrow$ Period $\rightarrow$ Epoch $\rightarrow$ Age (time); System $\rightarrow$ Series $\rightarrow$ Stage (rock).
\item Precambrian = Hadean + Archean + Proterozoic ($\sim$88\% of history); GOE at $\sim$2.4 Ga.
\item Learn the four anchor ages: \textbf{541, 252, 66.0 Ma} and \textbf{2\,500 Ma}; then the 12 period bases.
\item Boundaries are defined by \textbf{GSSPs} (``golden spikes''), calibrated numerically by radiometric dating.
\item Use the latest ICS ages; larger Ma = older.
\end{itemize}
\end{aretenir}

\coursec{Overview of Dating Methods}

In the previous section we built the \cle{geologic time scale}: a ladder of eons, eras, periods and epochs whose rungs were first defined by \emph{relative} criteria (superposition, fossils) and later calibrated in millions of years. But where do those numbers --- 66 Ma for the end of the dinosaurs, 541 Ma for the base of the Cambrian --- actually come from? To answer that question we must now look at the \emph{toolbox} of the geologist: the family of \cle{dating methods}. This section gives you the map of that toolbox; the next section will zoom into radiometric dating, the most powerful tool of all.

\courssub{Relative versus Absolute Dating}

Imagine you find two old photographs of the town of Buea. In the first, the mountain road is unpaved; in the second, it is tarred. Even without any date written on the photos, you know which one is \emph{older}: the unpaved road came first. That is \cle{relative dating}: placing events in order. If, however, one photo carries the stamp ``March 1978'', you have an \cle{absolute} (or \emph{numerical}) \cle{age}: a number attached to a time scale.

\begin{definition}[Relative and absolute dating]
\cle{Relative dating} determines the \emph{order} of geologic events (which rock, structure or fossil is older or younger) without giving an age in years. It rests on the stratigraphic principles (superposition, cross-cutting relationships, included fragments) and on fossil succession.

\cle{Absolute} (\cle{numerical}) \cle{dating} assigns an age in years (ka = thousands of years, Ma = millions of years) to a rock, mineral or object, usually by measuring a physical or chemical ``clock'' such as the radioactive decay of an isotope.
\end{definition}

\begin{attention}[A frequent confusion]
``Absolute'' does \emph{not} mean ``exact''. Every numerical age carries an \emph{uncertainty}. An age of $125.3 \pm 0.7$ Ma is absolute but not exact; a relative statement (``the lava flow is younger than the sandstone it crosses'') can be perfectly certain. Precision and certainty are different things.
\end{attention}

\courssub{Relative Dating Techniques}

Beyond the basic principles of Steno and the cross-cutting rules, modern stratigraphy has developed several specialised relative-dating techniques. Each uses a different kind of ``signal'' recorded in the rocks.

\textbf{1. Biostratigraphy.} This is the oldest and still the most widely used technique. It exploits the \emph{law of faunal succession}: fossil species appear and disappear in a definite, irreversible order. \cle{Index fossils} --- species that were widespread geographically but short-lived in time (many ammonites, graptolites, foraminifera) --- allow beds to be placed in a \cle{biozone} and correlated between distant outcrops. The Cretaceous limestones of the Douala and Rio del Rey basins, for example, are zoned using planktonic foraminifera.

\textbf{2. Magnetostratigraphy.} When lava cools or sediment settles, magnetic minerals (like magnetite) align with the Earth's magnetic field and \emph{freeze} its direction. The Earth's field has reversed polarity many times (normal $\leftrightarrow$ reversed). The resulting pattern of normal and reversed \cle{magnetozones} in a succession is like a barcode: it can be matched to the global \cle{Geomagnetic Polarity Time Scale} (GPTS). Magnetostratigraphy is especially powerful for dating the lava flows of the \cle{Cameroon Volcanic Line} and for correlating deep-sea sediment cores.

\textbf{3. Chemostratigraphy.} The chemical and isotopic composition of seawater has varied through time (e.g.\ the ratios $^{87}$Sr/$^{86}$Sr and $\delta^{13}$C recorded in marine carbonates). A measured curve of isotope values through a succession can be matched to the global reference curve, giving relative ages --- very useful in rocks \emph{without} fossils, such as some Precambrian carbonates of the Congo craton margin.

\textbf{4. Sequence stratigraphy.} Sea-level rises and falls leave characteristic packages of sediment (\cle{sequences}) bounded by erosion surfaces (\cle{sequence boundaries}). Because major sea-level changes are often global (eustatic), sequence boundaries can be approximately synchronous and serve as correlation horizons between basins --- a technique heavily used in petroleum exploration of the Rio del Rey and Douala basins.

\begin{center}
\begin{tabularx}{\linewidth}{|l|X|X|}
\hline
\rowcolor{popblueL} \textbf{Technique} & \textbf{Signal used} & \textbf{Best suited to} \\
\hline
Biostratigraphy & Index fossils, biozones & Fossiliferous sedimentary rocks \\
\hline
Magnetostratigraphy & Polarity reversals & Lavas, fine sediments, any age \\
\hline
Chemostratigraphy & Isotope / element curves & Marine carbonates, barren rocks \\
\hline
Sequence stratigraphy & Sea-level cycles, boundaries & Basin-fill sedimentary successions \\
\hline
\end{tabularx}
\end{center}

\begin{aretenir}[Relative dating in one sentence]
All relative techniques answer the question ``\emph{older or younger than what?}'' --- they order events and correlate rocks, but only an absolute clock can say ``\emph{how many years ago?}''.
\end{aretenir}

\courssub{Families of Absolute Dating Methods}

Absolute clocks fall into four main families. You should know what each measures and over what time range it works.

\textbf{1. Radiometric dating (parent--daughter decay).} A radioactive \cle{parent} isotope decays at a constant rate (half-life $T_{1/2}$) into a stable \cle{daughter} isotope. Measuring the parent/daughter ratio in a mineral gives the time elapsed since the system closed. Examples: K--Ar and $^{40}$Ar/$^{39}$Ar (volcanic rocks), U--Pb on zircon (very old rocks), Rb--Sr, Sm--Nd, U-series (e.g.\ U--Th for carbonates up to $\sim$500 ka), and \ce{^{14}C} for young organic material. Range: from a few thousand years ($^{14}$C) to the age of the Earth (U--Pb). This family is treated in detail in Section 5.

\textbf{2. Luminescence dating (OSL, TL).} Quartz and feldspar grains buried in sediment trap electrons in crystal defects, accumulating energy from natural radioactivity. Exposure to light (or heat) ``resets'' the clock. Measuring the stored luminescence gives the time since the grains were last exposed to sunlight --- i.e.\ since burial. Range: roughly 100 years to $\sim$300 ka. Ideal for dating dunes, river terraces and archaeological sediments that contain no dateable volcanic material.

\textbf{3. Cosmogenic nuclide dating.} Cosmic rays striking rocks at the Earth's surface produce rare isotopes such as \ce{^{10}Be} and \ce{^{26}Al} \emph{in situ}. Their concentration measures how long a surface (a moraine, a lava flow top, an exposed fault scarp) has been exposed to the sky. Range: roughly $10^3$ to a few $10^6$ years. Useful for landscape evolution, e.g.\ exposure ages of young flows on Mount Cameroon.

\textbf{4. Dendrochronology (tree-ring dating).} Counting annual growth rings of trees gives exact calendar ages, but only back a few thousand years (about 12 ka with overlapping sequences). Its great geological role is to \emph{calibrate} the radiocarbon clock.

\begin{savaistu}[Dating the Cameroon Volcanic Line]
The volcanoes of the Cameroon Volcanic Line, from Mount Cameroon on the coast to the Adamawa plateau, have been dated mainly by the K--Ar and $^{40}$Ar/$^{39}$Ar methods on feldspar and whole-rock samples. These ages show that volcanic activity has occurred along the line for tens of millions of years, with Mount Cameroon remaining active today --- it erupted several times during the 20th century.
\end{savaistu}

\courssub{Closure Temperature: When the Clock Starts Ticking}

Here is the key intuition. A radiometric clock does \emph{not} date the moment a mineral formed in a deep magma chamber; it dates the moment the mineral became a \emph{closed box} for parent and daughter atoms. While the rock is hot, atoms (especially argon, a gas) can diffuse in and out of the crystal, so no daughter product accumulates. Only when the rock cools enough does the mineral trap its daughter atoms, and the clock starts.

\begin{definition}[Closure temperature]
The \cle{closure temperature} ($T_c$) of an isotopic system in a given mineral is the temperature below which that mineral becomes a \emph{closed system}: parent and daughter isotopes can no longer escape or enter by diffusion, so the daughter product begins to accumulate and the radiometric clock starts.
\end{definition}

Different systems have different closure temperatures. Approximate values (worth knowing as orders of magnitude):

\begin{center}
\begin{tabularx}{\linewidth}{|l|X|l|}
\hline
\rowcolor{popblueL} \textbf{System / mineral} & \textbf{Typical use} & $T_c$ (approx.) \\
\hline
U--Pb in zircon & Crystallisation of granites, very old rocks & $> 900\ ^\circ$C \\
\hline
K--Ar / Ar--Ar in hornblende & Cooling of deep crustal rocks & $\sim 500\ ^\circ$C \\
\hline
K--Ar / Ar--Ar in biotite & Later cooling stages & $\sim 300\ ^\circ$C \\
\hline
Fission track in apatite & Near-surface cooling, uplift & $\sim 110\ ^\circ$C \\
\hline
\end{tabularx}
\end{center}

Two consequences follow. First, a single granite can yield \emph{different} ages from different minerals --- not an error, but a \emph{cooling history}: the zircon records crystallisation, the biotite records later cooling. Second, if a rock is \emph{reheated} above $T_c$ during a later metamorphic event, the clock is \emph{reset} (daughter products escape) and the measured age dates the reheating, not the original formation.

\begin{propriete}[The closed-system assumption]
A radiometric age is only as reliable as the assumption that the system remained \emph{closed}: no parent or daughter isotope was added or lost since closure, and the initial daughter content is known or negligible. Weathering, metamorphism or fluid circulation that opens the system produces ages that are too young, too old, or geologically meaningless. (The mathematical treatment of this assumption is examinable in Section 5.)
\end{propriete}

\courssub{Choosing the Right Dating Method}

No single method dates everything. The geologist must match the method to the sample, asking four questions:

\begin{methode}[Selecting a dating method]
\begin{enumerate}
\item \textbf{What is the rock type?} Igneous rocks $\rightarrow$ radiometric (K--Ar, U--Pb). Sedimentary rocks can rarely be dated directly $\rightarrow$ date bracketing lava flows or ash beds, or use luminescence on the sediment itself. Metamorphic rocks $\rightarrow$ systems with high $T_c$ (U--Pb, Sm--Nd).
\item \textbf{What is the expected age?} Younger than $\sim$50 ka with organic matter $\rightarrow$ \ce{^{14}C}. Thousands to hundreds of ka $\rightarrow$ luminescence, U-series, Ar--Ar. Millions to billions of years $\rightarrow$ K--Ar, U--Pb, Rb--Sr.
\item \textbf{Which minerals are available?} No zircon $\rightarrow$ no U--Pb. No K-bearing mineral $\rightarrow$ no K--Ar. No quartz/feldspar $\rightarrow$ no OSL.
\item \textbf{What precision is required?} $^{40}$Ar/$^{39}$Ar and U--Pb give the best analytical precision; choose the system whose uncertainty is small compared with the age differences you need to resolve.
\end{enumerate}
\end{methode}

\begin{popfigure}
\begin{tikzpicture}[
  font=\small,
  box/.style={draw=popblue, thick, rounded corners=2pt, fill=popblueL, text width=3.1cm, align=center, minimum height=1.05cm},
  dec/.style={draw=poporange, thick, fill=white, text width=2.9cm, align=center, minimum height=1.0cm},
  arr/.style={->, thick, popdark},
  lbl/.style={font=\footnotesize\itshape, popdark}
]
\node[dec] (q1) at (0,0) {Rock type?};
\node[dec] (q2) at (-4.3,-2.2) {Sedimentary: organic matter present?};
\node[dec] (q3) at (4.3,-2.2) {Igneous / metamorphic: expected age?};
\node[box] (c14) at (-6.6,-4.6) {\ce{^{14}C} if $< 50$ ka};
\node[box] (osl) at (-2.1,-4.6) {OSL / TL on quartz or feldspar};
\node[box] (ar) at (2.1,-4.6) {K--Ar or $^{40}$Ar/$^{39}$Ar (ka--Ma)};
\node[box] (upb) at (6.6,-4.6) {U--Pb on zircon (Ma--Ga)};
\draw[arr] (q1) -- (q2) node[lbl, midway, left=1pt] {sedimentary};
\draw[arr] (q1) -- (q3) node[lbl, midway, right=1pt] {igneous / meta.};
\draw[arr] (q2) -- (c14) node[lbl, midway, left=1pt] {yes};
\draw[arr] (q2) -- (osl) node[lbl, midway, right=1pt] {no};
\draw[arr] (q3) -- (ar) node[lbl, midway, left=1pt] {young};
\draw[arr] (q3) -- (upb) node[lbl, midway, right=1pt] {old};
\end{tikzpicture}
\legende{Decision tree for selecting an absolute dating method from sample characteristics (simplified).}
\end{popfigure}

\courssub{Reporting an Age with its Uncertainty}

A measurement without an uncertainty is scientifically useless. Analytical results are always reported as an age \emph{plus or minus} an uncertainty at a stated confidence level.

\begin{methode}[How to report an age]
\begin{enumerate}
\item Give the \textbf{central value} and the \textbf{uncertainty} with the same number of decimal places: $125.3 \pm 0.7$ Ma.
\item State the \textbf{confidence level}: ``$2\sigma$'' means there is about a 95\% probability that the true age lies within the interval; ``$1\sigma$'' corresponds to about 68\%.
\item State the \textbf{method} and the \textbf{mineral} dated: e.g.\ ``$^{40}$Ar/$^{39}$Ar on hornblende''.
\item Interpret the result: here the true age lies, with $\sim$95\% confidence, between 124.6 and 126.0 Ma (Early Cretaceous).
\end{enumerate}
\end{methode}

\begin{attention}[The carbon-14 trap]
\ce{^{14}C} has a half-life of only about 5\,730 years. After roughly 8--9 half-lives ($\sim$50 ka) so little \ce{^{14}C} remains that ages become unreliable. \textbf{Never} use carbon-14 to date a dinosaur bone (66 Ma old), a granite, or any rock: it works only on \emph{organic material younger than about 50\,000 years}. This is a classic GCE trap.
\end{attention}

\begin{exoresolu}[Choosing a method for a Cameroonian basalt]
\textbf{Statement.} A geologist collects a fresh basalt flow from the flank of Mount Cameroon. Field relations show it overlies a soil containing charcoal. She wishes to date the eruption. (a) Which absolute method(s) could date the basalt itself? (b) Which method could date the buried soil? (c) Why would \ce{^{14}C} be unsuitable for the basalt?
\tcblower
\textbf{Solution.}
\begin{enumerate}
\item[(a)] The basalt is an igneous rock containing K-bearing minerals (feldspar, glass): the \textbf{K--Ar or $^{40}$Ar/$^{39}$Ar method} is appropriate. The clock started when the lava cooled below the closure temperature of argon in those minerals --- essentially the eruption date.
\item[(b)] The soil contains \textbf{charcoal} (organic carbon). If the eruption is younger than $\sim$50 ka, \textbf{radiocarbon (\ce{^{14}C}) dating} of the charcoal gives the age of the soil just before burial --- a maximum age for the flow. If no charcoal were present, \textbf{OSL} on buried quartz grains could date the soil's last exposure to sunlight.
\item[(c)] Basalt contains \emph{no organic carbon}: \ce{^{14}C} dating measures the decay of carbon incorporated by living organisms. Moreover, if the flow were older than $\sim$50 ka, essentially all \ce{^{14}C} would have decayed anyway. Carbon-14 can never date a rock directly.
\end{enumerate}
\end{exoresolu}

\begin{exercice}[Matching methods to samples]
For each sample, propose the most suitable dating method and justify your choice in one sentence: (1) a zircon-bearing granite from the Precambrian basement north of Yaound\'e; (2) a sand dune estimated at 20\,000 years old, with no fossils; (3) a fossilised tree trunk from a Miocene lignite; (4) a Holocene lava flow of the Cameroon Volcanic Line.
\end{exercice}

\begin{corrige}[Exercise 1]
\begin{enumerate}
\item \textbf{U--Pb on zircon}: very old rock, and zircon's very high closure temperature preserves the crystallisation age even through later events.
\item \textbf{OSL (luminescence)}: quartz grains in the dune record time since burial; the age (20 ka) fits the OSL range, and no organic matter is needed.
\item \textbf{No direct method is ideal}: the trunk is Miocene (millions of years), far beyond \ce{^{14}C}; one must date \emph{bracketing} volcanic ash beds by K--Ar/Ar--Ar, or correlate the lignite biostratigraphically.
\item \textbf{$^{40}$Ar/$^{39}$Ar} (or K--Ar) on the lava: it dates cooling of the flow; charcoal beneath it could be cross-checked by \ce{^{14}C} if younger than 50 ka.
\end{enumerate}
\end{corrige}

\begin{aretenir}[Key points of this section]
\begin{itemize}
\item \textbf{Relative dating} orders events (biostratigraphy, magnetostratigraphy, chemostratigraphy, sequence stratigraphy); \textbf{absolute dating} gives ages in years.
\item The four absolute families are \textbf{radiometric}, \textbf{luminescence}, \textbf{cosmogenic nuclides} and \textbf{dendrochronology}, each with its own age range and target material.
\item The \textbf{closure temperature} defines when a mineral's isotopic clock starts; reheating can reset it.
\item A radiometric age is valid only if the system stayed \textbf{closed}.
\item Always report an age with its \textbf{uncertainty} (e.g.\ $125.3 \pm 0.7$ Ma, $2\sigma$) --- and never use \ce{^{14}C} beyond $\sim$50 ka or on rocks.
\end{itemize}
\end{aretenir}

\coursec{Radiometric Dating in Detail}

In the previous section we surveyed the landscape of dating methods and saw that radiometric dating is the only family of techniques capable of giving us \emph{absolute} ages in millions of years. We now open the hood of this remarkable engine. Our guiding question is simple: \emph{how can a handful of atoms inside a mineral grain tell us that a granite cooled 500 million years ago?} The answer rests on one of the most reliable clocks in nature — the spontaneous, unchanging decay of radioactive isotopes.

\courssub{The Radioactive Decay Law and Half-Life}

\textbf{Intuition first.} Imagine a very large bag of popcorn placed over a fire. You cannot predict \emph{which} kernel will pop next, but you can predict with great confidence that about half of the kernels will have popped after a certain time. Radioactive atoms behave exactly like this: each unstable (\emph{parent}) nucleus has a fixed probability of decaying per unit time into a stable (\emph{daughter}) nucleus, and this probability is unaffected by temperature, pressure, or chemical bonds. That is what makes radioactivity a trustworthy clock — no geological furnace or deep burial can speed it up or slow it down.

\begin{definition}[Radioactive decay and the decay constant]
\cle{Radioactive decay} is the spontaneous transformation of an unstable parent isotope (P) into a daughter isotope (D), often with emission of particles or radiation. The \cle{decay constant} $\lambda$ is the probability that a given nucleus decays per unit time; it is characteristic of each isotopic system.
\end{definition}

Because every nucleus decays independently with the same probability, the number of parent atoms $N$ decreases in proportion to the number remaining:

\[
\frac{dN}{dt} = -\lambda N
\]

Integrating from an initial quantity $N_0$ at $t = 0$ gives the \cle{decay law}:

\begin{propriete}[Radioactive decay law]
\[
N = N_0\, e^{-\lambda t}
\]
where $N$ is the number of parent atoms remaining after time $t$, $N_0$ the initial number, and $\lambda$ the decay constant (in $\mathrm{yr^{-1}}$ or $\mathrm{Ma^{-1}}$). The derivation by separation of variables is examinable.
\end{propriete}

The \cle{half-life} $t_{1/2}$ is the time needed for half of the parent atoms to decay. Setting $N = N_0/2$:

\[
\frac{N_0}{2} = N_0 e^{-\lambda t_{1/2}} \;\Longrightarrow\; \ln \tfrac{1}{2} = -\lambda t_{1/2} \;\Longrightarrow\; \boxed{t_{1/2} = \frac{\ln 2}{\lambda} \approx \frac{0.693}{\lambda}}
\]

Since a geologist measures the \emph{ratio} of daughter to parent rather than $N_0$, we rewrite the law in a practical form. The daughter produced is $D^* = N_0 - N = N\left(e^{\lambda t} - 1\right)$, so:

\[
\frac{D^*}{N} = e^{\lambda t} - 1 \qquad\Longrightarrow\qquad t = \frac{1}{\lambda}\ln\left(1 + \frac{D^*}{N}\right)
\]

This is the master equation of radiometric dating: measure the daughter-to-parent ratio today, know $\lambda$, and the age $t$ falls out.

\begin{attention}[Half-life is not "half the age"]
A very common mistake is to think that after one half-life the clock is "half used up" and after two half-lives it is finished. In fact the parent decreases \emph{exponentially}: after two half-lives one quarter remains, after three one eighth, and so on — the parent never mathematically reaches zero. Also, $t_{1/2}$ is a constant of the isotope; it does \emph{not} depend on how much material you have.
\end{attention}

\courssub{The Major Isotopic Systems}

Nature offers several parent–daughter pairs, each with its own half-life, suitable minerals, and geological niche. A good geochronologist chooses the system the way a mechanic chooses a spanner — the right tool for the job.

\begin{center}
\rowcolors{2}{popblueL}{white}
\begin{tabularx}{\linewidth}{|l|l|X|X|}
\hline
\rowcolor{popblueL}\textbf{System} & \textbf{Decay} & \textbf{Half-life / typical minerals} & \textbf{Best used for} \\
\hline
U--Pb & \ce{^{238}U -> ^{206}Pb}; \ce{^{235}U -> ^{207}Pb} & $t_{1/2}$: 4\,468 Ma and 704 Ma; zircon, monazite, apatite & Very old rocks; the gold standard for Precambrian terranes \\
\hline
K--Ar / Ar--Ar & \ce{^{40}K -> ^{40}Ar} & $t_{1/2}$ = 1\,250 Ma; micas, K-feldspar, hornblende & Volcanic rocks (e.g.\ Mount Cameroon lavas), cooling ages \\
\hline
Rb--Sr & \ce{^{87}Rb -> ^{87}Sr} & $t_{1/2}$ $\approx$ 48\,800 Ma; micas, K-feldspar, whole rocks & Ancient granites and metamorphic suites; isochron dating \\
\hline
Sm--Nd & \ce{^{147}Sm -> ^{143}Nd} & $t_{1/2}$ $\approx$ 106\,000 Ma; garnet, pyroxenes, whole rocks & Mafic/ultramafic rocks; crustal residence ages \\
\hline
Lu--Hf & \ce{^{176}Lu -> ^{176}Hf} & $t_{1/2}$ $\approx$ 37\,000 Ma; garnet, zircon & Garnet-bearing metamorphic rocks; crust evolution \\
\hline
Re--Os & \ce{^{187}Re -> ^{187}Os} & $t_{1/2}$ $\approx$ 41\,600 Ma; sulphides, organic-rich shales & Ore deposits, mantle rocks, black shales \\
\hline
\end{tabularx}
\end{center}

\begin{savaistu}[Dating the Cameroon Volcanic Line]
The volcanic mounts of the Cameroon Volcanic Line — from Mount Cameroon on the coast to the Adamawa plateau — have been dated largely by the \cle{K--Ar} and \cle{Ar--Ar} methods, because young basaltic lavas contain potassium-bearing minerals and glass that trap argon well. These ages revealed that the line became progressively active over tens of millions of years, with Mount Cameroon remaining active to the present day.
\end{savaistu}

\courssub{The Isochron Method: Eliminating the Unknown Initial Daughter}

\textbf{The problem.} Our master equation assumes that \emph{all} the daughter isotope present today was produced by decay. But when a mineral crystallises, it usually already contains some daughter isotope — the \emph{initial daughter} $D_0$. If we ignore it, we overestimate the age. The isochron method is the elegant graphical solution.

Take the Rb--Sr system. \ce{^{87}Rb} decays to \ce{^{87}Sr}, and we normalise everything by the stable, non-radiogenic isotope \ce{^{86}Sr}:

\[
\left(\frac{\ce{^{87}Sr}}{\ce{^{86}Sr}}\right)_{\text{today}} = \left(\frac{\ce{^{87}Sr}}{\ce{^{86}Sr}}\right)_{0} + \left(\frac{\ce{^{87}Rb}}{\ce{^{86}Sr}}\right)_{\text{today}} \left(e^{\lambda t} - 1\right)
\]

This is the equation of a straight line $y = b + mx$ where $y = \ce{^{87}Sr}/\ce{^{86}Sr}$, $x = \ce{^{87}Rb}/\ce{^{86}Sr}$, the intercept $b$ is the initial ratio, and the slope is $m = e^{\lambda t} - 1$.

\begin{definition}[Isochron]
An \cle{isochron} is a straight line on an isotope-ratio diagram (e.g.\ $\ce{^{87}Sr}/\ce{^{86}Sr}$ versus $\ce{^{87}Rb}/\ce{^{86}Sr}$) along which all points correspond to samples of the \emph{same age} that shared the \emph{same initial daughter isotope ratio}.
\end{definition}

\begin{propriete}[Slope of the isochron]
The slope of an isochron equals
\[
m = e^{\lambda t} - 1 \qquad\Longrightarrow\qquad t = \frac{1}{\lambda}\ln(1 + m).
\]
The y-intercept gives the initial daughter ratio $(\ce{^{87}Sr}/\ce{^{86}Sr})_0$ — a bonus piece of petrogenetic information. No proof is required at GCE level, but you must be able to \emph{use} the relation.
\end{propriete}

\begin{methode}[Constructing and reading a Rb--Sr isochron]
\begin{enumerate}
\item Collect several cogenetic samples (different minerals or whole rocks from the same intrusion).
\item Measure $\ce{^{87}Rb}/\ce{^{86}Sr}$ (x-axis) and $\ce{^{87}Sr}/\ce{^{86}Sr}$ (y-axis) for each.
\item Plot the points; if they fall on a straight line, the system stayed closed — the line is an isochron.
\item Measure the slope $m$ from the graph, then compute $t = \dfrac{1}{\lambda}\ln(1+m)$.
\item Read the initial ratio from the y-intercept.
\end{enumerate}
\end{methode}

\begin{popfigure}
\begin{center}
\begin{tikzpicture}[scale=1.05, font=\small]
% axes
\draw[->, thick, popdark] (0,0) -- (7.2,0) node[below left] {$\ce{^{87}Rb}/\ce{^{86}Sr}$};
\draw[->, thick, popdark] (0,0) -- (0,5.2) node[below left, rotate=90, anchor=south east] {$\ce{^{87}Sr}/\ce{^{86}Sr}$};
% isochron
\draw[very thick, popblue] (0,0.8) -- (6.6,4.6) node[above right, popblue] {isochron};
% initial ratio marker
\fill[poporange] (0,0.8) circle (2.2pt) node[left, popdark] {$(\ce{^{87}Sr}/\ce{^{86}Sr})_0$};
% data points
\foreach \x/\y in {1.2/1.35, 2.4/2.0, 3.6/2.65, 4.8/3.3, 6.0/4.0}{
  \fill[poppurple] (\x,\y) circle (2.2pt);
}
% slope triangle
\draw[dashed, popmuted] (4.8,3.3) -- (6.0,3.3) -- (6.0,4.0);
\node[popmuted, below] at (5.4,3.3) {$\Delta x$};
\node[popmuted, right] at (6.0,3.65) {$\Delta y$};
\node[align=center, popdark] at (4.6,0.7) {slope $m = \dfrac{\Delta y}{\Delta x} = e^{\lambda t}-1$};
\end{tikzpicture}
\end{center}
\legende{Schematic Rb--Sr isochron diagram. Cogenetic samples define a straight line whose slope gives the age and whose intercept gives the initial strontium isotope ratio.}
\end{popfigure}

\courssub{The Concordia--Discordia Diagram for U--Pb Zircons}

Uranium--lead dating enjoys a unique privilege: it contains \emph{two} independent clocks in the same mineral, \ce{^{238}U -> ^{206}Pb} and \ce{^{235}U -> ^{207}Pb}. If a zircon has remained a perfectly closed system since crystallisation, both clocks must give the \emph{same} age — the ages are \emph{concordant}.

\begin{definition}[Concordia]
The \cle{Concordia} is the theoretical curve on a plot of $\ce{^{206}Pb}/\ce{^{238}U}$ (y-axis) versus $\ce{^{207}Pb}/\ce{^{235}U}$ (x-axis) along which the two U--Pb ages agree exactly. Each point on the curve corresponds to one specific age.
\end{definition}

Real zircons often lose some of their radiogenic lead (for example during a later heating event). Lead loss moves a data point \emph{off} the Concordia, and a suite of zircons that all suffered the same disturbance plot along a straight line called the \cle{Discordia}. The Discordia intersects the Concordia at two points:

\begin{itemize}
\item the \textbf{upper intercept} gives the \emph{crystallisation age} of the zircons;
\item the \textbf{lower intercept} approximates the \emph{age of the disturbance} (the lead-loss event).
\end{itemize}

\begin{popfigure}
\begin{center}
\begin{tikzpicture}
\begin{axis}[
  width=11cm, height=8cm,
  xlabel={$\ce{^{207}Pb}/\ce{^{235}U}$},
  ylabel={$\ce{^{206}Pb}/\ce{^{238}U}$},
  xmin=0, xmax=16, ymin=0, ymax=0.7,
  domain=0.05:15, samples=120,
  grid=both, grid style={popline, very thin},
  legend style={at={(0.03,0.97)}, anchor=north west, font=\small},
  tick label style={font=\small}, label style={font=\small}
]
% Concordia curve (parametric via table-like expression using age parameter is complex; use approximate parametric points)
\addplot[very thick, popblue] coordinates {
(0.052,0.0057) (0.111,0.0121) (0.179,0.0303) (0.302,0.0606)
(0.464,0.0909) (0.675,0.1515) (0.947,0.2121) (1.295,0.2727)
(1.740,0.3333) (2.303,0.3939) (3.016,0.4545) (3.923,0.5152)
(5.090,0.5758) (6.617,0.6364) (8.70,0.68) (11.6,0.70) (15.5,0.70)};
\addlegendentry{Concordia}
% Discordia line
\addplot[very thick, poporange, dashed] coordinates {(1.0,0.09) (8.0,0.63)};
\addlegendentry{Discordia}
% data points along discordia
\addplot[only marks, mark=*, mark size=2.2pt, poppurple] coordinates {
(2.0,0.20) (3.5,0.31) (5.0,0.42) (6.5,0.52)};
\addlegendentry{zircon analyses}
% intercept markers
\addplot[only marks, mark=square*, mark size=2.4pt, popgreen] coordinates {(1.0,0.09) (8.0,0.63)};
\node[popdark, font=\small, anchor=west] at (axis cs:8.2,0.63) {upper intercept: crystallisation age};
\node[popdark, font=\small, anchor=west] at (axis cs:1.1,0.05) {lower intercept: disturbance};
\end{axis}
\end{tikzpicture}
\end{center}
\legende{Concordia--Discordia diagram. Undisturbed zircons lie on the Concordia; lead loss displaces analyses along the Discordia, whose intercepts date the crystallisation and the later disturbance.}
\end{popfigure}

\begin{attention}[Always check concordance]
Always check for concordance between the $\ce{^{238}U}/\ce{^{206}Pb}$ and $\ce{^{235}U}/\ce{^{207}Pb}$ ages before accepting a U--Pb result. If the two ages disagree (the point lies off the Concordia), the system has been disturbed and a single "age" is meaningless — you must use the Discordia intercepts instead.
\end{attention}

\courssub{Common Sources of Error}

Radiometric dating is powerful but not magical. The examiner expects you to know the four classic pitfalls:

\begin{itemize}
\item \textbf{Lead loss (U--Pb):} Pb is a relatively mobile element; metamorphic heating can expel part of the accumulated \ce{^{206}Pb} and \ce{^{207}Pb}, giving discordant, too-young ages. Detected by the Concordia diagram.
\item \textbf{Inheritance:} a zircon crystal may grow around an \emph{older} core (inherited from the source rock or an earlier crystal). Dating a mixed grain yields a spurious intermediate age. Modern practice dates single growth zones.
\item \textbf{Metamorphic resetting:} during a strong metamorphic event, isotopic clocks (especially K--Ar and Rb--Sr) may be partially or wholly reset, so the measured age records the \emph{metamorphism}, not the original crystallisation.
\item \textbf{Excess argon (K--Ar, Ar--Ar):} sometimes \ce{^{40}Ar} from surrounding fluids is trapped in the mineral at crystallisation. This "excess argon" mimics radiogenic argon and produces ages that are \emph{too old} — a notorious problem in very young volcanic rocks.
\end{itemize}

\begin{attention}[Closed system assumption]
Every radiometric age assumes the mineral behaved as a \cle{closed system}: no parent or daughter entered or left after crystallisation (except by decay). Whenever you quote an age, ask yourself whether this assumption is defensible — this is exactly what examiners probe in essay questions.
\end{attention}

\courssub{Practice Calculation: Age of a Zircon from $\ce{^{238}U}/\ce{^{206}Pb}$}

\begin{methode}[Age from a measured parent/daughter ratio]
\begin{enumerate}
\item Write the decay law in daughter/parent form: $\dfrac{D^*}{N} = e^{\lambda t} - 1$.
\item Identify the measured ratio. If given $\ce{^{238}U}/\ce{^{206}Pb} = N/D^*$, invert it.
\item Compute $t = \dfrac{1}{\lambda}\ln\!\left(1 + \dfrac{D^*}{N}\right)$, using $\lambda = \dfrac{\ln 2}{t_{1/2}}$.
\item State the age with units (Ma) and comment on concordance if a second system is available.
\end{enumerate}
\end{methode}

\begin{exoresolu}[Dating a zircon from a Cameroonian basement granite]
A zircon extracted from a granite of the Precambrian basement north of the Benue trough yields a measured ratio $\ce{^{238}U}/\ce{^{206}Pb} = 12.0$. Take $t_{1/2}(\ce{^{238}U}) = 4.468 \times 10^{9}$ yr and $\ln 2 = 0.693$. Calculate the age of the zircon, assuming no initial lead and a closed system.
\tcblower
\textbf{Step 1 — Decay constant:}
\[
\lambda = \frac{\ln 2}{t_{1/2}} = \frac{0.693}{4.468 \times 10^{9}\ \mathrm{yr}} = 1.551 \times 10^{-10}\ \mathrm{yr^{-1}}.
\]

\textbf{Step 2 — Daughter/parent ratio:}
\[
\frac{D^*}{N} = \frac{\ce{^{206}Pb}}{\ce{^{238}U}} = \frac{1}{12.0} = 0.0833.
\]

\textbf{Step 3 — Age equation:}
\[
t = \frac{1}{\lambda}\ln\left(1 + \frac{D^*}{N}\right)
= \frac{1}{1.551\times 10^{-10}}\ln(1.0833).
\]

Since $\ln(1.0833) = 0.0800$:
\[
t = \frac{0.0800}{1.551\times 10^{-10}} = 5.16 \times 10^{8}\ \mathrm{yr} \approx 516\ \mathrm{Ma}.
\]

\textbf{Conclusion:} the zircon crystallised about \textbf{516 million years ago} (early Palaeozoic, Cambrian). Before accepting the result, we would verify that the $\ce{^{235}U}/\ce{^{207}Pb}$ age is concordant.
\end{exoresolu}

\begin{exercice}[Check your mastery]
A biotite from a schist gives a Rb--Sr isochron with slope $m = 0.0143$. Given $\lambda(\ce{^{87}Rb}) = 1.42 \times 10^{-11}\ \mathrm{yr^{-1}}$, calculate the age of the rock. Then state two reasons why a K--Ar age on the same rock might come out \emph{younger}.
\end{exercice}

\begin{corrige}[Exercise 1]
$t = \dfrac{1}{\lambda}\ln(1+m) = \dfrac{\ln(1.0143)}{1.42\times10^{-11}} = \dfrac{0.0142}{1.42\times10^{-11}} = 1.0\times10^{9}$ yr $= 1000$ Ma. A K--Ar age could be younger because (i) metamorphic heating caused argon loss (partial resetting), or (ii) the K--Ar system records later \emph{cooling} below the argon closure temperature rather than crystallisation.
\end{corrige}

\begin{aretenir}[Key points of Section 5]
\begin{itemize}
\item Decay law: $N = N_0 e^{-\lambda t}$; half-life $t_{1/2} = \ln 2/\lambda$; practical form $t = \frac{1}{\lambda}\ln(1 + D^*/N)$.
\item The \cle{isochron} method removes the unknown initial daughter: slope $= e^{\lambda t}-1$, intercept $=$ initial ratio.
\item U--Pb zircon dating is self-checking via the \cle{Concordia}; disturbed samples define a \cle{Discordia} whose upper intercept dates crystallisation.
\item Beware lead loss, inheritance, metamorphic resetting and excess argon — all violate the closed-system assumption.
\end{itemize}
\end{aretenir}

With the machinery of absolute dating now in hand, we are ready to examine what happens when rock successions meet: the \emph{stratigraphic contacts and unconformities} that record gaps and dramas in Earth's history — the subject of the next section.

\coursec{Stratigraphic Contacts and Unconformities}

In the previous section, we saw how radiometric dating gives us \emph{absolute} ages for rocks. But long before geologists could measure ages in millions of years, they could already read the \emph{relative} history of a region simply by looking at how rock units touch one another. The surfaces that separate rock units --- the \cle{stratigraphic contacts} --- are the punctuation marks of the stratigraphic record: some mark continuous, quiet deposition; others mark dramatic interruptions of uplift, folding and erosion. In this section we learn to recognise, classify and interpret these contacts, a skill that is systematically tested at the GCE Advanced Level, both in theory and in map- and section-based exercises.

\courssub{What is a Stratigraphic Contact?}

\begin{definition}[Stratigraphic contact]
A \cle{stratigraphic contact} is the surface that separates two distinct rock units (formations, members, or any mappable bodies of rock). It may correspond to a bedding plane, an erosional surface, or an intrusive or faulted boundary.
\end{definition}

Think of a layered cake: the interface between two layers is a contact. If the baker poured the second layer immediately after the first, the interface is smooth and continuous --- this is a \cle{conformable} contact. If, instead, the first cake was sliced, tilted and partly eaten before the second layer was added, the interface records a \emph{gap in time and a period of disturbance} --- this is an \cle{unconformable} contact.

Contacts are fundamental because they are the surfaces we trace on geological maps and the boundaries we measure in stratigraphic sections. Their \emph{nature} (conformable or not) tells us whether deposition was continuous or interrupted.

\courssub{Conformable Contacts}

A contact is \cle{conformable} when the strata above and below were deposited in an essentially continuous sequence, without significant erosion or interruption of sedimentation. Beds above and below are parallel, and palaeontological content shows a regular succession. Three varieties are distinguished:

\begin{itemize}
\item \textbf{Gradational contact:} one rock type passes progressively into the other (e.g. shale passing upward into siltstone then sandstone as the sea shallows). There is no single sharp surface; the boundary is placed by convention where a chosen lithological change occurs.
\item \textbf{Sharp (abrupt) conformable contact:} lithology changes suddenly (e.g. limestone directly overlain by shale) but without any evidence of erosion or time gap. Such contacts reflect a rapid change in depositional conditions, for instance a sudden marine transgression.
\item \textbf{Paraconformity:} a special case discussed below, in which the contact looks conformable (beds are parallel) but actually conceals a time gap.
\end{itemize}

\begin{attention}[Conformable does not always mean ``no time gap'']
A sharp conformable contact can look dramatic yet represent continuous deposition, while a perfectly parallel contact (paraconformity) can hide millions of years of non-deposition. Never judge a contact by its geometry alone: look for \emph{evidence} of erosion or missing time.
\end{attention}

\courssub{Unconformities: Definition and Types}

\begin{definition}[Unconformity]
An \cle{unconformity} is a surface of erosion or non-deposition that separates younger strata from older rocks and represents a significant \cle{gap in the geological record} (a \emph{hiatus}). It records an interval during which deposition stopped and/or previously deposited rocks were eroded.
\end{definition}

Four main types are recognised, and you must be able to define, sketch and identify each of them.

\begin{definition}[Angular unconformity]
An \cle{angular unconformity} is an erosional surface on which younger, roughly horizontal strata rest upon older strata that have been \emph{tilted or folded} and truncated by erosion. The beds above and below the surface are \textbf{not parallel}.
\end{definition}

An angular unconformity records a full tectonic cycle: (1) deposition of the older sequence; (2) tectonic deformation (tilting/folding) and uplift; (3) erosion levelling the landscape; (4) subsidence and deposition of the younger sequence.

\begin{definition}[Disconformity]
A \cle{disconformity} is an unconformity in which the strata above and below the erosional surface are \emph{parallel}, but the surface itself shows clear evidence of erosion (irregular topography, channels, basal conglomerate, paleosol). It records uplift and erosion \emph{without} significant deformation.
\end{definition}

\begin{definition}[Nonconformity]
A \cle{nonconformity} is an unconformity in which sedimentary strata rest directly on eroded \emph{igneous or metamorphic} (crystalline) basement. It records deep erosion that stripped away the cover rocks before renewed deposition. In Cameroon, sediments of the Douala basin resting on the Precambrian basement of the Congo craton illustrate this relationship.
\end{definition}

\begin{propriete}[The paraconformity: a hidden gap]
A \cle{paraconformity} is an unconformity in which the beds above and below are parallel and the surface shows \emph{no obvious erosional features}; the time gap is detected mainly through a \emph{fossil (biostratigraphic) break} --- one or more biozones are missing between the two units. It typically records a long period of non-deposition (e.g. a submarine hiatus) rather than active erosion.
\end{propriete}

\begin{popfigure}
\begin{tikzpicture}[scale=0.9, every node/.style={font=\small}]
% ---- Panel 1: Angular unconformity ----
\begin{scope}[shift={(0,0)}]
\draw[fill=popblueL] (-0.2,0) rectangle (3.4,1.0);
\foreach \y in {0.25,0.5,0.75}{\draw[popblue!60] (-0.2,\y) -- (3.4,\y);}
\draw[fill=popmuted!40] (-0.2,-1.2) -- (3.4,-1.2) -- (3.4,0) -- (-0.2,0) -- cycle;
\foreach \x in {0,0.5,1.0,1.5,2.0,2.5,3.0}{\draw[popdark!50] (\x,-1.2) -- (\x+0.55,0);}
\draw[very thick,poporange] (-0.2,0) -- (3.4,0);
\node[align=center] at (1.6,-1.55) {\textbf{(a) Angular unconformity}};
\node[align=center,font=\footnotesize] at (1.6,1.25) {horizontal young beds};
\node[align=center,font=\footnotesize] at (1.6,-0.65) {tilted old beds};
\end{scope}
% ---- Panel 2: Disconformity ----
\begin{scope}[shift={(4.6,0)}]
\draw[fill=popblueL] (-0.2,0.15) rectangle (3.4,1.0);
\foreach \y in {0.4,0.65,0.9}{\draw[popblue!60] (-0.2,\y) -- (3.4,\y);}
\draw[fill=popmuted!40] (-0.2,-1.2) rectangle (3.4,0.15);
\foreach \y in {-0.95,-0.6,-0.25}{\draw[popdark!50] (-0.2,\y) -- (3.4,\y);}
\draw[very thick,poporange] (-0.2,0.15) .. controls (0.6,0.3) and (1.2,-0.05) .. (1.8,0.15) .. controls (2.4,0.32) and (2.9,0.0) .. (3.4,0.15);
\node[align=center] at (1.6,-1.55) {\textbf{(b) Disconformity}};
\node[align=center,font=\footnotesize] at (1.6,1.25) {parallel beds, wavy erosion};
\end{scope}
% ---- Panel 3: Nonconformity ----
\begin{scope}[shift={(9.2,0)}]
\draw[fill=popblueL] (-0.2,0.15) rectangle (3.4,1.0);
\foreach \y in {0.4,0.65,0.9}{\draw[popblue!60] (-0.2,\y) -- (3.4,\y);}
\draw[fill=poppink!35] (-0.2,-1.2) rectangle (3.4,0.15);
\foreach \p in {(0.3,-0.4),(1.1,-0.8),(2.0,-0.35),(2.8,-0.75),(0.7,-0.95),(2.5,-0.55)}{\node[font=\tiny,popdark] at \p {+};}
\draw[very thick,poporange] (-0.2,0.15) .. controls (0.8,0.3) and (1.6,0.0) .. (2.4,0.2) .. controls (2.9,0.3) .. (3.4,0.15);
\node[align=center] at (1.6,-1.55) {\textbf{(c) Nonconformity}};
\node[align=center,font=\footnotesize] at (1.6,-0.55) {crystalline basement};
\end{scope}
% ---- Panel 4: Paraconformity ----
\begin{scope}[shift={(13.8,0)}]
\draw[fill=popblueL] (-0.2,0) rectangle (3.4,1.0);
\foreach \y in {0.25,0.5,0.75}{\draw[popblue!60] (-0.2,\y) -- (3.4,\y);}
\draw[fill=popmuted!40] (-0.2,-1.2) rectangle (3.4,0);
\foreach \y in {-0.95,-0.6,-0.25}{\draw[popdark!50] (-0.2,\y) -- (3.4,\y);}
\draw[very thick,poporange,dashed] (-0.2,0) -- (3.4,0);
\node[align=center] at (1.6,-1.55) {\textbf{(d) Paraconformity}};
\node[align=center,font=\footnotesize] at (1.6,1.25) {flat surface, fossil gap only};
\end{scope}
\end{tikzpicture}
\legende{The four main types of unconformity (orange surface). (a) Angular unconformity: truncated tilted beds below. (b) Disconformity: parallel beds with an irregular erosional surface. (c) Nonconformity: sediments on eroded crystalline basement. (d) Paraconformity: parallel beds, no visible erosion; the gap is revealed by missing fossil zones.}
\end{popfigure}

\courssub{Recognition Criteria in the Field and on Maps}

How do we \emph{prove} that a contact is an unconformity and not simply a facies change? We look for diagnostic evidence:

\begin{itemize}
\item \textbf{Basal conglomerate:} a conglomerate at the base of the younger unit containing pebbles derived from the underlying rocks --- direct proof that the older unit was lithified and eroded before the younger one was deposited.
\item \textbf{Paleosol (fossil soil):} a horizon of weathering, root traces or lateritic alteration just below the contact, proving exposure at the Earth's surface.
\item \textbf{Truncation:} beds, folds, dykes or faults in the older unit are cut off by the contact surface; on a geological map, the contact cuts across the outcrop pattern of the older unit.
\item \textbf{Fossil break:} missing biozones between the two units (the key criterion for a paraconformity).
\item \textbf{Age gap:} radiometric or biostratigraphic ages show a hiatus between the two units.
\item \textbf{Relief on the surface:} channels, potholes or an irregular, wavy geometry (disconformity).
\end{itemize}

\begin{methode}[Field checklist for identifying an unconformity]
\begin{enumerate}
\item \textbf{Locate the contact} and describe its geometry: flat, wavy, channelled? Are beds above and below parallel?
\item \textbf{Examine the base of the upper unit:} search for a basal conglomerate; identify the pebbles --- do they match the rocks below?
\item \textbf{Examine the top of the lower unit:} look for weathering, paleosol, root traces, iron staining.
\item \textbf{Check for truncation:} are beds, folds, dykes or faults below cut by the surface?
\item \textbf{Sample for fossils} above and below: is any biozone missing?
\item \textbf{Classify the contact} (angular unconformity, disconformity, nonconformity, paraconformity) and \textbf{estimate the hiatus} if ages are available.
\end{enumerate}
\end{methode}

\begin{attention}[Do not confuse a facies change with an unconformity]
A lateral or vertical \cle{facies change} (e.g. sandstone passing to shale) can produce a sharp contact with \emph{no time gap at all}. Before calling a surface an unconformity, you must demonstrate \emph{erosion or non-deposition}: basal conglomerate, paleosol, truncation, or a fossil/age gap. Sharpness of a contact alone proves nothing.
\end{attention}

\courssub{Significance of Unconformities}

\begin{itemize}
\item \textbf{Basin history:} each unconformity records a tectonic event --- uplift, deformation, sea-level fall. Counting and dating unconformities allows the geologist to reconstruct the subsidence/uplift history of a sedimentary basin and to subdivide the fill into \emph{sequences}.
\item \textbf{Missing record:} the hiatus means part of Earth history is simply absent at that locality; correlation with more complete sections elsewhere is needed.
\item \textbf{Hydrocarbon traps:} unconformities are major exploration targets. Weathered rocks just below the surface are often porous (good \emph{reservoirs}), while fine-grained sediments deposited above form a \emph{seal}; truncation of tilted reservoir beds against the unconformity creates \emph{stratigraphic traps}. Several traps of this type are known in the Rio del Rey and Douala basins of Cameroon.
\item \textbf{Mineralisation and aquifers:} weathering profiles above unconformities can concentrate lateritic ores, and the permeable zone along the surface can channel groundwater.
\end{itemize}

\begin{savaistu}[Hutton's unconformity]
The concept was born in 1788 when James Hutton, the father of modern geology, examined the outcrop at Siccar Point (Scotland), where near-horizontal Devonian sandstones rest on vertical Silurian greywackes. Hutton realised that an immense span of time --- deposition, tilting, mountain building, erosion, then renewed deposition --- was recorded in that single surface. This convinced him that the Earth was vastly older than commonly believed.
\end{savaistu}

\courssub{Sketching and Labelling Contacts on a Measured Section}

In the GCE practical-style questions you may be asked to draw a stratigraphic column from field data. Proceed as follows:

\begin{methode}[Drawing a measured section with contacts]
\begin{enumerate}
\item Draw a vertical axis with a scale (e.g. 1 cm = 10 m) and mark thicknesses upward from the base.
\item For each bed, draw a rectangle whose width is constant and whose height equals the bed thickness; ornament it with the standard lithological symbol (dots for sandstone, dashes for shale, bricks for limestone).
\item Draw each \textbf{contact} according to its nature: a simple line for a conformable contact; a bold wavy line for a disconformity; a bold line with truncated, differently oriented beds below for an angular unconformity.
\item Label each unit (name, lithology, thickness, age) and each contact (type of contact, evidence observed).
\item Add a title, scale and north/location information if it is a field section.
\end{enumerate}
\end{methode}

\begin{exoresolu}[Interpreting a contact in the field]
\textbf{Statement.} Along a road cut, a geologist observes horizontal red sandstones resting on steeply dipping grey shales. At the base of the sandstone there is a 30 cm bed of conglomerate whose pebbles are made of grey shale identical to that below. Fossils indicate that the shale is Ordovician and the sandstone is Cretaceous. (a) Name the type of contact. (b) List the evidence supporting your answer. (c) Reconstruct the geological history.
\tcblower
\textbf{Solution.}
\begin{enumerate}
\item \textbf{(a)} The contact is an \cle{angular unconformity}: the beds above (horizontal sandstone) and below (steeply dipping shale) are not parallel, and an erosional surface separates them.
\item \textbf{(b)} Evidence: (i) angular discordance between horizontal and dipping strata; (ii) truncation of the dipping shale beds by the contact; (iii) a \emph{basal conglomerate} containing pebbles reworked from the underlying shale, proving the shale was lithified and eroded before sandstone deposition; (iv) a large \emph{time gap} (Ordovician to Cretaceous) with missing systems.
\item \textbf{(c)} History: (1) deposition of the grey shales in the Ordovician; (2) tectonic deformation tilting the shales, with uplift; (3) prolonged erosion levelling the surface; (4) renewed deposition of red sandstones in the Cretaceous, beginning with a basal conglomerate.
\end{enumerate}
\end{exoresolu}

\begin{exercice}[Classify and justify]
For each situation, name the type of contact and justify: (1) horizontal Cenozoic sands resting on eroded Precambrian gneiss of the Congo craton; (2) parallel Cretaceous limestones and marls separated by a surface with a lateritic paleosol; (3) parallel Eocene and Oligocene shales with two missing planktonic foraminifera zones but no erosional feature.
\end{exercice}

\begin{corrige}[Exercise 1]
\begin{enumerate}
\item \textbf{Nonconformity:} sedimentary rocks rest on eroded crystalline (metamorphic) basement.
\item \textbf{Disconformity:} beds are parallel but the paleosol proves exposure and erosion/non-deposition.
\item \textbf{Paraconformity:} beds parallel, no erosional surface; the hiatus is revealed only by the missing fossil zones.
\end{enumerate}
\end{corrige}

\begin{aretenir}[Key points]
\begin{itemize}
\item A \cle{stratigraphic contact} separates two rock units; it is \emph{conformable} (continuous deposition) or \emph{unconformable} (hiatus).
\item Four unconformity types: \cle{angular unconformity} (tilted/folded beds below), \cle{disconformity} (parallel beds, erosional surface), \cle{nonconformity} (sediments on crystalline basement), \cle{paraconformity} (parallel beds, gap shown only by fossils).
\item Diagnostic evidence: basal conglomerate, paleosol, truncation, fossil/age gap.
\item Unconformities record tectonic and sea-level events and can form stratigraphic hydrocarbon traps.
\item Never call a contact an unconformity without evidence of erosion or non-deposition.
\end{itemize}
\end{aretenir}

\coursec{Sequencing Geologic Events}

In Section 6 we learned to recognise and interpret stratigraphic contacts --- conformable surfaces, unconformities, intrusive contacts and fault contacts. Each contact, taken alone, tells us something about the order in which two rock bodies were emplaced. But real outcrops, and real GCE examination questions, rarely present a single contact: they present a \emph{puzzle} in which sedimentation, folding, intrusion, erosion and faulting are all tangled together. The skill we now develop is the one that examiners reward most highly: reading a complex cross-section and reconstructing its \cle{geologic history} --- a complete, ordered narrative of events from oldest to youngest.

\courssub{From contacts to chronology: the logic of relative dating}

\begin{experience}[A roadside cutting near Bafoussam]
Imagine a road cutting exposing, from bottom to top: folded schists, a granite that cuts across them, a horizontal sandstone resting on the eroded top of both, and a dolerite dyke that cuts the sandstone but stops at a fault which offsets everything. Even without a single numerical age, you can already say: the schists are oldest, the granite came next, then erosion, then the sandstone, then the dyke, then the fault. You have just built a relative chronology using nothing but the principles of superposition, cross-cutting relationships and unconformities.
\end{experience}

The whole of this section rests on a small toolkit of principles already met in Sections 2 and 6, now applied \emph{simultaneously}:

\begin{itemize}
\item \textbf{Superposition:} in an undisturbed sedimentary pile, the lowest bed is oldest.
\item \textbf{Original horizontality:} tilted or folded strata were deposited flat, so deformation post-dates deposition.
\item \textbf{Cross-cutting relationships:} whatever cuts (dyke, fault, vein) is younger than whatever is cut.
\item \textbf{Unconformities:} an erosional surface separates older, deformed and eroded rocks below from younger rocks above.
\item \textbf{Inclusions:} fragments enclosed in a rock are older than the rock containing them (e.g.\ schist xenoliths in granite).
\end{itemize}

\begin{propriete}[Age of an intrusion]
An intrusion is \textbf{younger than all the rocks it cuts} (cross-cutting) but \textbf{older than all the rocks that overlie it unconformably} or that contain its eroded pebbles. The same logic applies to faults and to any cross-cutting body. (Proof not examinable; the property follows directly from the cross-cutting and superposition principles.)
\end{propriete}

\courssub{The Event List technique}

Faced with a complicated section, students who ``just look and guess'' almost always miss an event or invert two of them. The rigorous approach is systematic.

\begin{methode}[The Event List technique]
\begin{enumerate}
\item \textbf{Inventory.} List every observable feature on the section: each sedimentary unit, each fold, each intrusion, each fault, each unconformity (erosional surface). Give each a label (A, B, C\ldots).
\item \textbf{Local ordering.} For each pair of features in contact, apply one principle to decide which is older. Write down the evidence (``dyke D cuts unit B $\Rightarrow$ D younger than B'').
\item \textbf{Global ordering.} Chain the local decisions into a single ordered list from oldest to youngest. Check for contradictions (e.g.\ a fault claimed to be both older and younger than the same unit).
\item \textbf{Insert absolute ages.} Where radiometric dates are given, place them in the list and use them to \emph{bracket} the undated events.
\item \textbf{Write the narrative.} Convert the ordered list into a prose geologic history, oldest event first.
\end{enumerate}
\end{methode}

\begin{definition}[Bracketing age]
The \cle{bracketing age} of a rock unit or event is the age \emph{interval} constrained by the age of the youngest dated unit \emph{below} it (or cut by it) and the age of the oldest dated unit \emph{above} it (or cutting it). The true age lies between the two bounds.
\end{definition}

\begin{exemplebox}[Bracketing a lava flow]
A basalt flow rests on sandstone dated by fossils to the Cretaceous and is overlain by a tuff dated radiometrically at \SI{60}{Ma}. The basalt is therefore younger than Cretaceous and older than \SI{60}{Ma}: its bracketing age is ``Late Cretaceous to earliest Palaeogene''. Along the Cameroon Volcanic Line, lavas are commonly bracketed between dated flows above and below.
\end{exemplebox}

\courssub{A multi-event cross-section}

\begin{popfigure}
\begin{tikzpicture}[scale=1.05, every node/.style={font=\small}]
% Basement folded strata
\fill[popblueL] (0,0) rectangle (8,2.2);
\draw[popblue, thick] (0,0.4) .. controls (1,1.1) and (1.6,-0.1) .. (2.6,0.5)
  .. controls (3.6,1.1) and (4.2,-0.1) .. (5.2,0.5)
  .. controls (6.2,1.1) and (6.8,-0.1) .. (8,0.6);
\draw[popblue, thick] (0,1.1) .. controls (1,1.8) and (1.6,0.6) .. (2.6,1.2)
  .. controls (3.6,1.8) and (4.2,0.6) .. (5.2,1.2)
  .. controls (6.2,1.8) and (6.8,0.6) .. (8,1.3);
\draw[popblue, thick] (0,1.8) .. controls (1,2.4) and (1.6,1.3) .. (2.6,1.9)
  .. controls (3.6,2.4) and (4.2,1.3) .. (5.2,1.9)
  .. controls (6.2,2.4) and (6.8,1.3) .. (8,2.0);
% Granite intrusion
\fill[poppink!40] (2.6,0) -- (2.9,0) -- (3.4,2.2) -- (3.0,2.2) -- cycle;
\draw[popdark] (2.6,0) -- (3.0,2.2);
\draw[popdark] (2.9,0) -- (3.4,2.2);
% Unconformity surface
\draw[poporange, very thick] (0,2.2) -- (8,2.2);
% Horizontal cover strata
\fill[popgreenL] (0,2.2) rectangle (8,3.1);
\fill[popgold!35] (0,3.1) rectangle (8,3.9);
\draw[popmuted] (0,2.65) -- (8,2.65);
\draw[popmuted] (0,3.1) -- (8,3.1);
\draw[popmuted] (0,3.55) -- (8,3.55);
% Dyke cutting cover
\fill[poppurple!35] (5.6,2.2) -- (5.9,2.2) -- (6.2,3.9) -- (5.9,3.9) -- cycle;
\draw[popdark] (5.6,2.2) -- (5.9,3.9);
\draw[popdark] (5.9,2.2) -- (6.2,3.9);
% Fault
\draw[popink, very thick] (6.9,0) -- (7.5,3.9);
% Labels
\node[align=center] at (1.1,1.5) {folded\\strata (A)};
\node[align=center] at (3.15,0.7) {granite\\(B)};
\node[align=center] at (1.6,2.62) {sandstone (C)};
\node[align=center] at (1.6,3.5) {limestone (D)};
\node[align=center] at (6.0,3.15) {dyke (E)};
\node[align=center] at (7.6,1.6) {fault\\(F)};
\node[poporange, align=center] at (4.6,2.42) {unconformity (U)};
\end{tikzpicture}
\legende{Multi-event cross-section for annotation. Events: deposition of A, folding of A, intrusion of granite B, erosion (unconformity U), deposition of C then D, intrusion of dyke E, faulting F.}
\end{popfigure}

\begin{exoresolu}[Worked example: reading the cross-section]
Using the cross-section above, establish the complete sequence of geologic events from oldest to youngest. Granite B is dated at \SI{550}{Ma} and dyke E at \SI{120}{Ma}. What is the bracketing age of the unconformity U and of the folding?
\tcblower
\textbf{Step 1 --- Inventory.} Features: folded strata A; granite B; unconformity U; sandstone C; limestone D; dyke E; fault F.

\textbf{Step 2 --- Local ordering.}
\begin{itemize}
\item A is folded $\Rightarrow$ deposition of A, then folding (original horizontality).
\item B cuts the \emph{folded} fabric of A $\Rightarrow$ B is younger than A \emph{and} younger than the folding.
\item U truncates both A and B $\Rightarrow$ the erosion surface is younger than B.
\item C and D overlie U $\Rightarrow$ C then D (superposition), both younger than U.
\item E cuts C and D $\Rightarrow$ E younger than D.
\item F offsets everything, including E $\Rightarrow$ F is the youngest event.
\end{itemize}

\textbf{Step 3 --- Ordered list (oldest $\to$ youngest):} deposition of A $\to$ folding of A $\to$ intrusion of B (\SI{550}{Ma}) $\to$ uplift and erosion (U) $\to$ deposition of C $\to$ deposition of D $\to$ intrusion of E (\SI{120}{Ma}) $\to$ faulting F.

\textbf{Step 4 --- Bracketing.} The folding affects A but predates B (\SI{550}{Ma}): folding is therefore \emph{older than} \SI{550}{Ma} (its lower bound is the undated age of A). The unconformity U post-dates B (\SI{550}{Ma}) and pre-dates C; if C is, say, fossil-dated to the Jurassic, then U is bracketed between \SI{550}{Ma} and the Jurassic. The fault F is younger than \SI{120}{Ma} because it offsets dyke E.

\textbf{Step 5 --- Narrative.} ``Sediments of unit A were deposited in horizontal layers, then compressed and folded. A granite magma (B) intruded the folded rocks at \SI{550}{Ma}. The region was uplifted and deeply eroded, producing unconformity U. A marine transgression deposited sandstone C then limestone D. A dolerite dyke (E) intruded at \SI{120}{Ma}, and finally fault F displaced the whole pile.''
\end{exoresolu}

\courssub{Integrating absolute ages with the relative sequence}

Radiometric dates (Section 5) do not replace the relative sequence --- they \emph{anchor} it. The strategy is always the same:

\begin{enumerate}
\item Build the relative sequence first, from field relationships alone.
\item Slot each dated rock into its position in the sequence.
\item Every undated event then inherits a bracketing age from its dated neighbours.
\end{enumerate}

\begin{exemplebox}[Anchoring a sequence]
Suppose in the section above we additionally know that limestone D contains Early Cretaceous fossils. Then: dyke E (\SI{120}{Ma}, Early Cretaceous) is consistent with cutting D; the unconformity U is bracketed between \SI{550}{Ma} (granite below) and the Early Cretaceous (limestone above); and fault F is younger than \SI{120}{Ma}. Notice how each new date tightens the brackets.
\end{exemplebox}

\begin{attention}[Faults and intrusions: read the contact carefully]
A fault that \textbf{offsets} an intrusion must be \textbf{younger} than the intrusion. But a fault that is \textbf{sealed} (covered and undisturbed) by an overlying sedimentary unit is \textbf{older} than that unit. The same fault can therefore be bracketed: younger than everything it cuts, older than everything that seals it. Never assume a fault is ``recent'' just because it looks sharp on a diagram.
\end{attention}

\begin{attention}[The two classic pitfalls]
\begin{itemize}
\item \textbf{Assuming simultaneous deformation.} Folding, intrusion and faulting on one section are almost never one event. Each structure must be placed separately in the sequence using its own cross-cutting evidence. Two folds with different orientations may record two distinct orogenic phases separated by an unconformity.
\item \textbf{Ignoring hidden contacts.} A contact may be concealed beneath younger cover or below the base of the section. If dyke E disappears beneath unit C, you cannot claim it cuts C --- check whether it is truncated by the unconformity (older) or intrudes beneath it (younger). Always ask: \emph{what happens at the contact, even where I cannot see it?}
\end{itemize}
\end{attention}

\courssub{Guided practice: three cross-sections of increasing difficulty}

\begin{exercice}[Level 1 --- Simple pile with one intrusion]
A section shows, from bottom to top: shale (unit 1), sandstone (unit 2), both cut by a basalt dyke (unit 3), and a conglomerate (unit 4) resting unconformably on the eroded top of units 1, 2 and 3. The dyke is dated at \SI{200}{Ma}. (a) List the events oldest to youngest. (b) Give the bracketing age of the unconformity.
\end{exercice}

\begin{exercice}[Level 2 --- Folding, two intrusions, one fault]
Folded greywackes (A) are intruded by granite (G1, dated \SI{600}{Ma}). An unconformity (U1) separates these from horizontal sandstones (B). A second granite (G2) cuts B but is truncated by a second unconformity (U2) overlain by limestones (C). A fault (F) offsets A, G1, B and G2 but not C. Establish the full sequence and bracket the age of G2 and of F.
\end{exercice}

\begin{exercice}[Level 3 --- The full puzzle]
Using the annotated cross-section of this section, but suppose additionally that: (i) a second dyke E$'$ cuts the granite B but is truncated by unconformity U; (ii) fault F is itself sealed by a horizontal basalt flow (G) dated at \SI{30}{Ma}. Reconstruct the complete geologic history, bracket the age of dyke E$'$, of the folding, and of fault F, and explain what the existence of E$'$ tells us about the duration of magmatic activity in the region.
\end{exercice}

\begin{corrige}[Exercise Level 1]
(a) Events, oldest to youngest: deposition of shale 1 $\to$ deposition of sandstone 2 $\to$ intrusion of basalt dyke 3 (\SI{200}{Ma}) $\to$ uplift and erosion (unconformity) $\to$ deposition of conglomerate 4. (b) The unconformity post-dates the dyke (\SI{200}{Ma}) and pre-dates conglomerate 4; its bracketing age is therefore between \SI{200}{Ma} and the age of unit 4 (younger than \SI{200}{Ma}).
\end{corrige}

\begin{corrige}[Exercise Level 2]
Sequence: deposition of A $\to$ folding of A $\to$ intrusion of G1 (\SI{600}{Ma}) $\to$ erosion (U1) $\to$ deposition of B $\to$ intrusion of G2 $\to$ erosion (U2) $\to$ deposition of C $\to$ faulting F. Bracketing: G2 is younger than B (which it cuts) and older than U2 and C (which truncate and overlie it); since B is younger than \SI{600}{Ma}, G2 is younger than \SI{600}{Ma} and older than C. Fault F is younger than G2 (it offsets it) and older than C (which seals it, since F does not offset C); F is therefore bracketed between the intrusion of G2 and the deposition of C.
\end{corrige}

\begin{corrige}[Exercise Level 3]
Full history: deposition of A $\to$ folding of A $\to$ intrusion of granite B (\SI{550}{Ma}) $\to$ intrusion of dyke E$'$ (cuts B) $\to$ uplift and erosion (U, truncating B and E$'$) $\to$ deposition of C $\to$ deposition of D $\to$ intrusion of dyke E (\SI{120}{Ma}) $\to$ faulting F $\to$ eruption of basalt G (\SI{30}{Ma}) sealing F. Bracketing: E$'$ is younger than \SI{550}{Ma} and older than U (hence older than C); the folding is older than \SI{550}{Ma}; fault F is bracketed between \SI{120}{Ma} (it offsets E) and \SI{30}{Ma} (it is sealed by G). Dyke E$'$ shows that magmatism was \emph{episodic and long-lived}: at least two intrusive episodes separated by a major erosion interval, a pattern typical of repeatedly reactivated magmatic provinces such as the Cameroon Volcanic Line.
\end{corrige}

\begin{aretenir}[Key points of Section 7]
\begin{itemize}
\item Reconstruct a geologic history with the \cle{Event List}: inventory, local ordering, global ordering, insert absolute ages, write the narrative.
\item An intrusion (or fault) is younger than everything it cuts and older than everything that overlies or seals it.
\item A \cle{bracketing age} constrains an undated event between the youngest dated unit below and the oldest dated unit above.
\item Absolute dates anchor the relative sequence; they never replace field relationships.
\item Beware the two traps: assuming simultaneous deformation, and ignoring hidden or concealed contacts.
\end{itemize}
\end{aretenir}

\begin{savaistu}[Did you know?]
The same logic you have just used on paper was applied in the field to the rocks around Mount Cameroon: geologists recognised that some lava flows are cut by dykes, while others blanket and seal older faulted terrain --- allowing the recent volcanic and tectonic history of the Cameroon Volcanic Line to be ordered in time even before precise radiometric dates were available. Relative chronology remains the geologist's first and most powerful tool.
\end{savaistu}

\coursec{Correlation of Stratigraphic Successions}

In the previous section we learned how to read a single outcrop or borehole and reconstruct the sequence of events that produced it. Real geological mapping, however, rarely deals with isolated sections. To understand the regional geology of the Douala Basin, the Cameroon Volcanic Line, or the Congo Craton margin, we must demonstrate that the sandstone at the base of a hill near Kribi is the same unit as the sandstone cropping out 50 km away near Edea, or that a limestone in the Mamfe Basin is time-equivalent to a shale in the Rio del Rey Basin. This is the science of \textbf{correlation}: the demonstration of equivalence between stratigraphic units in different geographic areas. Correlation transforms local successions into a regional stratigraphic framework, allowing us to build paleogeographic maps, locate hydrocarbon reservoirs, and calibrate the Geologic Time Scale.

\courssub{Principles and Objectives of Correlation}

Correlation rests on the fundamental principle that \emph{layers which are the same age and formed in the same depositional environment will share similar physical, paleontological, and chemical characteristics}. The objective is to establish \emph{equivalence} of three distinct types:

\begin{itemize}
\item \textbf{Lithostratigraphic equivalence}: same rock type, same position in the succession (same formation).
\item \textbf{Biostratigraphic equivalence}: same fossil assemblage, hence same relative age (same biozone).
\item \textbf{Chronostratigraphic equivalence}: same absolute time interval (same stage, e.g., Cenomanian).
\end{itemize}

A successful correlation integrates all three. In the Douala Basin, for instance, the \emph{Mundeck Formation} (lithostratigraphy) contains \emph{Orbitolina} foraminifera (biostratigraphy) that place it firmly in the Aptian--Albian (chronostratigraphy). When we find the same trio in a well log offshore, we have correlated the onshore outcrop with the subsurface.

\begin{definition}[Marker Bed]
A \cle{marker bed} (or key bed) is a distinctive, laterally extensive layer -- volcanic ash (bentonite), limestone, coal seam, or a unique conglomerate -- that can be recognized across a basin and used as a precise correlation horizon. Its lateral continuity and unique lithology make it an isochronous surface (deposited at the same time).
\end{definition}

\courssub{Lithostratigraphic Correlation}

Lithostratigraphic correlation matches rock bodies based on \emph{lithology} (composition, texture, colour), \emph{thickness}, \emph{sedimentary structures}, and \emph{vertical sequence}. It is the first step in the field.

\begin{methode}[Lithostratigraphic Correlation Workflow]
\begin{enumerate}
\item Describe each section in detail (litholog): bed thickness, grain size, sedimentary structures, fossil content.
\item Identify distinctive marker beds (e.g., a 2 m thick oolitic limestone, a bentonite layer).
\item Match the vertical succession of lithologies between sections (pattern matching).
\item Trace marker beds laterally on the ground or in seismic/well logs.
\item Define formations and members where the lithology is consistent and mappable.
\end{enumerate}
\end{methode}

\emph{Example -- Douala Basin}: The \emph{Logbaba Formation} (Cretaceous) is recognized by its basal conglomerate, followed by cross-bedded sandstones and a distinctive red shale marker bed. In the field, you walk from the type section at Logbaba towards the coast, noting the red shale at the same stratigraphic level, confirming lateral continuity.

\begin{attention}[Facies Changes]
Lithostratigraphic correlation fails when a formation changes laterally into a different rock type (facies change). A sandstone may pinch out into a shale, or a reef limestone may grade into basinal marl. Correlating the sandstone with the shale \emph{as the same formation} is a classic error. Always seek biostratigraphic or chronostratigraphic control across facies boundaries.
\end{attention}

\courssub{Biostratigraphic Correlation}

Fossils provide the most precise tool for correlating sedimentary rocks of Phanerozoic age. The principle is \emph{faunal/floral succession}: fossil assemblages succeed each other in a definite, irreversible order.

\begin{propriete}[Biostratigraphic Zones and FAD/LAD]
Biostratigraphic zones (biozones) are bodies of rock defined by their fossil content. The most useful for correlation are \cle{range zones} defined by the \cle{First Appearance Datum (FAD)} and \cle{Last Appearance Datum (LAD)} of an \cle{index fossil} -- a species that is geographically widespread, environmentally tolerant, rapidly evolving, abundant, and easily preserved/identified. The zone represents the time interval between the FAD and LAD of that taxon.
\end{propriete}

Other biozone types include \emph{assemblage zones} (characteristic association of taxa), \emph{abundance zones} (acme zones), and \emph{interval zones} (between two biohorizons). A \cle{range chart} plots the stratigraphic ranges of multiple taxa; the zones of overlap allow high-resolution correlation.

\begin{exemplebox}[Correlation using Ammonites in the Mamfe Basin]
Two sections 30 km apart in the Mamfe Basin show different lithologies: Section A has sandy shales, Section B has calcareous mudstones. Both yield ammonites. Section A yields \emph{Desmoceras latidorsatum} (FAD: Middle Albian) and \emph{Mortoniceras} sp. (LAD: Late Albian). Section B yields \emph{Desmoceras latidorsatum} and \emph{Stoliczkaia} (FAD: Late Albian). The overlap of \emph{D. latidorsatum} proves the two sections share a Middle--Late Albian interval, despite the facies difference.
\end{exemplebox}

\courssub{Chronostratigraphic Correlation}

Chronostratigraphic correlation aligns local successions to the \emph{Global Geologic Time Scale} via \cle{Global Boundary Stratotype Sections and Points (GSSPs)} -- the "golden spikes" that define the base of each stage. Three powerful tools allow this alignment:

\begin{enumerate}
\item \textbf{Magnetostratigraphy}: Correlation of the pattern of magnetic polarity reversals (chrons) recorded in the rocks to the Global Polarity Time Scale (GPTS). A sedimentary section yielding a pattern of normal/reverse intervals (e.g., N-R-N-R) can be matched uniquely to the GPTS, providing absolute age control.
\item \textbf{Chemostratigraphy}: Correlation using global geochemical signals. The most famous is the \ensuremath{\mathrm{\delta^{13}C}} curve. Major excursions (e.g., the \emph{Weissert Event} in the Valanginian, the \emph{OAE1a} in the Aptian) are global isochrons. Measuring \ensuremath{\mathrm{\delta^{13}C}} in bulk carbonate or organic matter from a section in the Douala Basin and matching the curve to the global reference instantly ties the section to the absolute time scale.
\item \textbf{Astrochronology}: Where Milankovitch cycles (precession, obliquity, eccentricity) are recorded in rhythmic bedding (e.g., limestone--marl couplets), spectral analysis can provide a floating astronomical time scale that is anchored to radiometric dates.
\end{enumerate}

\begin{savaistu}[The Aptian OAE1a in Cameroon]
Oceanic Anoxic Event 1a (OAE1a, $\sim$120 Ma) is marked by a prominent positive \ensuremath{\mathrm{\delta^{13}C}} excursion and widespread black shales. In the Douala Basin, the \emph{Mundeck Formation} contains organic-rich shales correlated to OAE1a via chemostratigraphy, confirming an Early Aptian age and linking Cameroon's stratigraphy to a global climate perturbation.
\end{savaistu}

\courssub{Geophysical Correlation: Seismic and Well Logs}

In subsurface exploration (oil, gas, groundwater), correlation relies on geophysical logs and seismic reflection profiles.

\begin{itemize}
\item \textbf{Well logs}: Gamma-ray (GR) logs distinguish shales (high GR) from clean sandstones/carbonates (low GR). A distinctive GR pattern (e.g., a "serrated" coarsening-upward motif) serves as a marker. Resistivity, sonic, and density logs provide additional correlation horizons.
\item \textbf{Seismic reflectors}: A seismic line images acoustic impedance contrasts. Continuous, high-amplitude reflectors often correspond to major stratigraphic boundaries (unconformities, flooding surfaces, volcanic layers). \cle{Seismic stratigraphy} correlates these reflectors across the basin, tying well logs to the seismic grid.
\end{itemize}

\begin{methode}[Integrated Well--Seismic Correlation]
\begin{enumerate}
\item Pick formation tops on well logs (GR, resistivity, sonic).
\item Generate synthetic seismograms from sonic/density logs.
\item Match synthetic wiggles to seismic reflectors at the well location.
\item Extend reflector interpretation across the seismic grid.
\item Map the correlated horizons (time or depth structure maps).
\end{enumerate}
\end{methode}

\courssub{Correlation Across Facies Changes}

This is the ultimate test of correlation skill. When a shallow-marine sandstone (facies A) passes laterally into a deep-marine shale (facies B), lithostratigraphy breaks down. The solution is \emph{chronostratigraphic correlation} using time-parallel surfaces:

\begin{itemize}
\item \cle{Maximum Flooding Surfaces (MFS)}: The most landward extent of marine transgression; often a condensed section rich in fossils and with a distinct GR/log signature. It is a time line.
\item \cle{Sequence Boundaries}: Subaerial unconformities (or their correlative conformities) formed during base-level fall.
\item \cle{Biohorizons}: FAD/LAD of planktonic microfossils (foraminifera, nannofossils) that lived in the open ocean and rained down across all facies.
\item \cle{Chemostratigraphic events}: \ensuremath{\mathrm{\delta^{13}C}} excursions recorded in both carbonate (shallow) and organic matter (deep).
\end{itemize}

\begin{attention}[The Lithology Trap]
Never correlate solely on similar lithology across a facies change. A sandstone in a delta front and a sandstone in a submarine fan may look identical but be millions of years apart. Always ask: "What is the age control?" (fossils, isotopes, magnetostratigraphy, radiometric dates).
\end{attention}

\courssub{Practical Workflow: Building a Correlation Panel}

A \cle{correlation panel} (or fence diagram) is the standard way to present correlated sections. It aligns multiple vertical sections along a datum (usually a marker bed, a sequence boundary, or sea level) to show lateral changes in thickness, facies, and age.

\begin{methode}[Stepwise Correlation Workflow for a Panel]
\begin{enumerate}
\item \textbf{Data gathering}: Collect lithologs, fossil lists, radiometric ages, gamma-ray logs, seismic picks for all sections.
\item \textbf{Choose a datum}: Select a regional, time-significant horizon (e.g., Base Cretaceous Unconformity, MFS, or a dated volcanic ash).
\item \textbf{Hang sections on the datum}: Align all sections vertically so the datum is horizontal.
\item \textbf{Correlate marker beds}: Draw lines connecting identical marker beds (bentonites, coal seams, distinctive limestones).
\item \textbf{Apply biostratigraphy}: Plot FAD/LAD of index fossils; draw biozone boundaries horizontally (they are time lines).
\item \textbf{Apply chronostratigraphy}: Use magnetostratigraphy, \ensuremath{\mathrm{\delta^{13}C}} chemostratigraphy, or GSSP-tied events to draw stage boundaries (e.g., Aptian/Albian).
\item \textbf{Interpret facies changes}: Show how lithologies change laterally between sections; label depositional environments.
\item \textbf{Finalize panel}: Add scale, legend, location map, and interpretation (systems tracts, sequences).
\end{enumerate}
\end{methode}

\begin{exoresolu}[Constructing a Correlation Panel: Three Wells in the Douala Basin]
\textbf{Statement:} Three exploration wells (Well A -- onshore, Well B -- coastal, Well C -- offshore) penetrate the Cretaceous succession of the Douala Basin. The following data are available (depths in meters below rotary table, ages in Ma). The wells penetrate a post-Cretaceous cover before entering the Cretaceous at the Base Cretaceous Unconformity (BCU).

\begin{center}
\begin{tabular}{|l|c|c|c|}
\hline
\textbf{Feature} & \textbf{Well A (Onshore)} & \textbf{Well B (Coastal)} & \textbf{Well C (Offshore)} \\
\hline
Surface (Rotary Table) & 0 m & 0 m & 0 m \\
\hline
\textbf{Base Cretaceous Unconformity (BCU) / Base Logbaba Fm} & \textbf{120 m} & \textbf{350 m} & \textbf{850 m} \\
\hline
Bentonite Marker Bed "Beta" (Radiometric: 112.5 $\pm$ 0.3 Ma) & 210 m & 480 m & 1020 m \\
\hline
FAD \emph{Orbitolina} sp. (Early Aptian) & 215 m & 490 m & 1030 m \\
\hline
Positive \ensuremath{\mathrm{\delta^{13}C}} excursion (OAE1a, $\sim$120 Ma) & 250 m & 540 m & 1100 m \\
\hline
LAD \emph{Orbitolina} sp. (Late Aptian) & 380 m & 720 m & 1350 m \\
\hline
Top \emph{Mundeck Fm} (MFS, Base Albian) & 450 m & 850 m & 1500 m \\
\hline
Total Depth & 600 m & 1200 m & 2000 m \\
\hline
\end{tabular}
\end{center}

Construct a correlation panel at a vertical scale of 1 cm = 100 m, using the Base Cretaceous Unconformity as datum. Identify the formations, biozones, and key chronostratigraphic boundaries.

\textbf{Detailed Solution:}

\noindent\textbf{Step 1: Choose Datum and Scale.} Datum = Base Cretaceous Unconformity (BCU), which marks the base of the \emph{Logbaba Formation} (Aptian). Scale: 1 cm = 100 m. On the panel, the BCU is drawn as a horizontal line. The post-Cretaceous cover (0 to BCU) is shown above the datum; the Cretaceous section is shown below.

\noindent\textbf{Step 2: Calculate Depths Below Datum (Positive Downwards).}
\begin{center}
\begin{tabular}{|l|ccc|}
\hline
\textbf{Horizon} & \textbf{Well A (m)} & \textbf{Well B (m)} & \textbf{Well C (m)} \\
\hline
BCU (Datum) & 0 & 0 & 0 \\
Bentonite $\beta$ & 90 & 130 & 170 \\
FAD \emph{Orbitolina} & 95 & 140 & 180 \\
OAE1a ($\delta^{13}$C) & 130 & 190 & 250 \\
LAD \emph{Orbitolina} & 260 & 370 & 500 \\
Top Mundeck (MFS) & 330 & 500 & 650 \\
Total Depth & 480 & 850 & 1150 \\
\hline
\end{tabular}
\end{center}

\noindent\textbf{Step 3: Correlate Horizons (Time Lines).}
On a depth panel, horizons representing the \emph{same geological instant} (isochrons) will plot at different depths below the datum due to differential subsidence and sedimentation rates. We connect them with \textbf{sloping correlation lines}.
\begin{itemize}
\item \textbf{BCU (Datum)}: Horizontal line (by definition).
\item \textbf{OAE1a (120 Ma)}: Global isochron (Early Aptian). Connect points at 1.30, 1.90, 2.50 cm.
\item \textbf{Bentonite $\beta$ (112.5 Ma)}: Radiometric isochron (Late Aptian). Connect points at 0.90, 1.30, 1.70 cm. Note: It lies \emph{above} (younger than) OAE1a in all wells -- consistent.
\item \textbf{FAD \emph{Orbitolina}}: Biohorizon (Early Aptian). Lies just above Bentonite $\beta$ (5--10 m). Connect points.
\item \textbf{LAD \emph{Orbitolina}}: Biohorizon (Late Aptian). Connect points at 2.60, 3.70, 5.00 cm.
\item \textbf{Top Mundeck (MFS, $\sim$113 Ma)}: Sequence boundary / Maximum Flooding Surface (Aptian/Albian boundary). Connect points at 3.30, 5.00, 6.50 cm.
\end{itemize}

\noindent\textbf{Step 4: Interpret Formations and Biozones.}
\begin{itemize}
\item \textbf{Post-Cretaceous Cover}: 0 to BCU (1.2, 3.5, 8.5 cm thick). Likely Tertiary/Volcanic.
\item \textbf{Logbaba Formation (Aptian)}: BCU to Top Mundeck. Thickness: Well A 330 m, Well B 500 m, Well C 650 m. Thickens basinward.
\item \textbf{Mundeck Formation (Albian)}: Top Mundeck to TD. Thickness: Well A 150 m, Well B 350 m, Well C 500 m. Thickens basinward.
\item \textbf{\emph{Orbitolina} Range Zone}: Spans FAD to LAD, covering most of the Aptian (Logbaba Fm).
\end{itemize}

\noindent\textbf{Step 5: Draw the Panel.}
The figure below is a schematic correlation panel (depth domain). Lithology patterns: Post-Cretaceous (stipple), Logbaba Fm -- Sandstone/Shale (dots/horizontal lines), Mundeck Fm -- Shale/Limestone (horizontal lines/bricks), Bentonite (zigzag).

\begin{popfigure}
\begin{tikzpicture}[scale=0.7]
% --- Parameters ---
\pgfmathsetmacro{\wellSep}{5.5} % horizontal separation between well centers
\pgfmathsetmacro{\wellW}{2.5}   % well column width
\pgfmathsetmacro{\datumY}{0}    % y-coordinate of BCU datum
% Depths below datum in cm (from table above)
% Well A (x=0)
\pgfmathsetmacro{\Axc}{0}
\pgfmathsetmacro{\Acover}{1.2} % 120m
\pgfmathsetmacro{\Abent}{0.90}
\pgfmathsetmacro{\Aorbfad}{0.95}
\pgfmathsetmacro{\Aoae}{1.30}
\pgfmathsetmacro{\Aorblad}{2.60}
\pgfmathsetmacro{\Amfs}{3.30}
\pgfmathsetmacro{\Atd}{4.80}
% Well B (x=5.5)
\pgfmathsetmacro{\Bxc}{\wellSep}
\pgfmathsetmacro{\Bcover}{3.5}
\pgfmathsetmacro{\Bbent}{1.30}
\pgfmathsetmacro{\Borbfad}{1.40}
\pgfmathsetmacro{\Boae}{1.90}
\pgfmathsetmacro{\Borblad}{3.70}
\pgfmathsetmacro{\Bmfs}{5.00}
\pgfmathsetmacro{\Btd}{8.50}
% Well C (x=11.0)
\pgfmathsetmacro{\Cxc}{2*\wellSep}
\pgfmathsetmacro{\Ccover}{8.5}
\pgfmathsetmacro{\Cbent}{1.70}
\pgfmathsetmacro{\Corbfad}{1.80}
\pgfmathsetmacro{\Coae}{2.50}
\pgfmathsetmacro{\Corblad}{5.00}
\pgfmathsetmacro{\Cmfs}{6.50}
\pgfmathsetmacro{\Ctd}{11.50}
\pgfmathsetmacro{\chronX}{\Cxc+\wellW/2+1.5}
\pgfmathsetmacro{\chronW}{3.0}
\pgfmathsetmacro{\midAlb}{(\Bmfs+\Btd)/2}
\pgfmathsetmacro{\halfW}{\wellW/2}
\pgfmathsetmacro{\halfChron}{\chronW/2}

% --- Colors ---
\definecolor{coverCol}{RGB}{200,180,160} % Post-Cretaceous
\definecolor{sandCol}{RGB}{230,210,160}  % Sandstone
\definecolor{shaleCol}{RGB}{160,160,160} % Shale
\definecolor{limeCol}{RGB}{180,210,190}  % Limestone/Marl
\definecolor{bentCol}{RGB}{140,100,60}   % Bentonite
\definecolor{corrCol}{RGB}{200,0,0}      % Correlation lines red
\definecolor{bioCol}{RGB}{0,0,180}       % Biohorizons blue
\definecolor{chemCol}{RGB}{0,150,0}      % Chemo green
\definecolor{mfsCol}{RGB}{150,0,150}     % MFS purple

% --- Draw Well Columns ---
% Well A
\fill[coverCol] (\Axc-\halfW, \datumY) rectangle (\Axc+\halfW, -\Acover);
\fill[sandCol] (\Axc-\halfW, \datumY) rectangle (\Axc+\halfW, \Amfs); % Logbaba (simplified sand/shale)
\fill[limeCol] (\Axc-\halfW, \Amfs) rectangle (\Axc+\halfW, \Atd);    % Mundeck
% Well B
\fill[coverCol] (\Bxc-\halfW, \datumY) rectangle (\Bxc+\halfW, -\Bcover);
\fill[sandCol] (\Bxc-\halfW, \datumY) rectangle (\Bxc+\halfW, \Bmfs);
\fill[limeCol] (\Bxc-\halfW, \Bmfs) rectangle (\Bxc+\halfW, \Btd);
% Well C
\fill[coverCol] (\Cxc-\halfW, \datumY) rectangle (\Cxc+\halfW, -\Ccover);
\fill[sandCol] (\Cxc-\halfW, \datumY) rectangle (\Cxc+\halfW, \Cmfs);
\fill[limeCol] (\Cxc-\halfW, \Cmfs) rectangle (\Cxc+\halfW, \Ctd);

% --- Draw Correlation Lines (Isochrons) ---
% Style
\tikzset{corrLine/.style={thick, corrCol}}
\tikzset{bioLine/.style={dashed, thick, bioCol}}
\tikzset{chemLine/.style={dotted, thick, chemCol}}
\tikzset{mfsLine/.style={dashdotted, thick, mfsCol}}

% BCU (Datum) - Horizontal
\draw[very thick, black] (\Axc-\halfW-0.5, \datumY) -- (\Cxc+\halfW+0.5, \datumY) node[midway, above, black] {Base Cretaceous Unconformity (Datum)};

% OAE1a (120 Ma) - Chemostratigraphy
\draw[chemLine] (\Axc-\halfW, \Aoae) -- (\Cxc+\halfW, \Coae) node[pos=0.5, above, sloped, chemCol, font=\small] {OAE1a $\delta^{13}$C (120 Ma)};

% Bentonite Beta (112.5 Ma) - Radiometric
\draw[corrLine] (\Axc-\halfW, \Abent) -- (\Cxc+\halfW, \Cbent) node[pos=0.5, above, sloped, corrCol, font=\small] {Bentonite $\beta$ (112.5 Ma)};

% FAD Orbitolina
\draw[bioLine] (\Axc-\halfW, \Aorbfad) -- (\Cxc+\halfW, \Corbfad) node[pos=0.5, above, sloped, bioCol, font=\small] {FAD \emph{Orbitolina}};

% LAD Orbitolina
\draw[bioLine] (\Axc-\halfW, \Aorblad) -- (\Cxc+\halfW, \Corblad) node[pos=0.5, above, sloped, bioCol, font=\small] {LAD \emph{Orbitolina}};

% Top Mundeck (MFS) - Sequence Stratigraphy
\draw[mfsLine] (\Axc-\halfW, \Amfs) -- (\Cxc+\halfW, \Cmfs) node[pos=0.5, above, sloped, mfsCol, font=\small] {Top Mundeck MFS (Apt/Alb Bound.)};

% --- Well Labels & Surface Lines ---
% Well A
\draw[thick] (\Axc-\halfW, -\Acover) -- (\Axc+\halfW, -\Acover);
\node[align=center, font=\bfseries\small] at (\Axc, -\Acover-0.8) {Well A\\(Onshore)};
\node[align=left, font=\tiny, anchor=north west] at (\Axc-\halfW, -\Acover) {Cover};
\node[align=left, font=\tiny, anchor=north west] at (\Axc-\halfW, \datumY) {Logbaba Fm (Aptian)};
\node[align=left, font=\tiny, anchor=north west] at (\Axc-\halfW, \Amfs) {Mundeck Fm (Albian)};

% Well B
\draw[thick] (\Bxc-\halfW, -\Bcover) -- (\Bxc+\halfW, -\Bcover);
\node[align=center, font=\bfseries\small] at (\Bxc, -\Bcover-0.8) {Well B\\(Coastal)};
\node[align=left, font=\tiny, anchor=north west] at (\Bxc-\halfW, -\Bcover) {Cover};
\node[align=left, font=\tiny, anchor=north west] at (\Bxc-\halfW, \datumY) {Logbaba Fm (Aptian)};
\node[align=left, font=\tiny, anchor=north west] at (\Bxc-\halfW, \Bmfs) {Mundeck Fm (Albian)};

% Well C
\draw[thick] (\Cxc-\halfW, -\Ccover) -- (\Cxc+\halfW, -\Ccover);
\node[align=center, font=\bfseries\small] at (\Cxc, -\Ccover-0.8) {Well C\\(Offshore)};
\node[align=left, font=\tiny, anchor=north west] at (\Cxc-\halfW, -\Ccover) {Cover};
\node[align=left, font=\tiny, anchor=north west] at (\Cxc-\halfW, \datumY) {Logbaba Fm (Aptian)};
\node[align=left, font=\tiny, anchor=north west] at (\Cxc-\halfW, \Cmfs) {Mundeck Fm (Albian)};

% --- Chronostratigraphic Column (Right Side) ---
% Time scale: 120 Ma (OAE1a), 112.5 Ma (Bent), 113 Ma (Apt/Alb).
% We draw a schematic column showing stages.
\draw[thick] (\chronX, \datumY) rectangle (\chronX+\chronW, \Ctd+0.5); % Frame
% Stage boundaries (approximate depths on panel based on Well B as reference)
% Aptian: Base at BCU (0), Top at MFS (~5.0 cm on Well B).
% Albian: Top at TD? (8.5 cm).
\node[align=center, font=\bfseries\small, rotate=90] at (\chronX-0.5, 0) {CHRONOSTRATIGRAPHY};
% Aptian
\fill[popblueL!50] (\chronX, \datumY) rectangle (\chronX+\chronW, \Bmfs);
\node[align=center, font=\small] at (\chronX+\halfChron, \Bmfs/2) {APTIAN\\Logbaba Fm};
% Albian
\fill[popgreenL!50] (\chronX, \Bmfs) rectangle (\chronX+\chronW, \Btd);
\node[align=center, font=\small] at (\chronX+\halfChron, \midAlb) {ALBIAN\\Mundeck Fm};
% Markers on chrono column
\draw[chemLine] (\chronX, \Boae) -- (\chronX+\chronW, \Boae) node[right, font=\tiny, chemCol] {OAE1a 120 Ma};
\draw[corrLine] (\chronX, \Bbent) -- (\chronX+\chronW, \Bbent) node[right, font=\tiny, corrCol] {Bent. $\beta$ 112.5 Ma};
\draw[bioLine] (\chronX, \Borbfad) -- (\chronX+\chronW, \Borbfad) node[right, font=\tiny, bioCol] {FAD Orb.};
\draw[bioLine] (\chronX, \Borblad) -- (\chronX+\chronW, \Borblad) node[right, font=\tiny, bioCol] {LAD Orb.};
\draw[mfsLine] (\chronX, \Bmfs) -- (\chronX+\chronW, \Bmfs) node[right, font=\tiny, mfsCol] {Apt/Alb Bound.};

% Scale bar
\draw[|-|, thick] (\Axc-\halfW-1.5, \datumY) -- (\Axc-\halfW-1.5, \datumY+1) node[midway, left, font=\small] {100 m};

% Legend
\begin{scope}[shift={(\Axc-3, \Atd+0.5)}, font=\tiny]
\fill[coverCol] (0,0) rectangle (0.5,0.3);
\node[right] at (0.6,0.15) {Post-Cretaceous Cover};
\fill[sandCol] (0,-0.4) rectangle (0.5,-0.1);
\node[right] at (0.6,-0.25) {Logbaba Fm (Sand/Shale)};
\fill[limeCol] (0,-0.8) rectangle (0.5,-0.5);
\node[right] at (0.6,-0.65) {Mundeck Fm (Shale/Lime)};
\draw[corrLine] (0,-1.05) -- (0.5,-1.05);
\node[right] at (0.6,-1.05) {Bentonite $\beta$ (Isochron)};
\draw[chemLine] (0,-1.45) -- (0.5,-1.45);
\node[right] at (0.6,-1.45) {OAE1a $\delta^{13}$C};
\draw[bioLine] (0,-1.85) -- (0.5,-1.85);
\node[right] at (0.6,-1.85) {Biohorizons (FAD/LAD)};
\draw[mfsLine] (0,-2.25) -- (0.5,-2.25);
\node[right] at (0.6,-2.25) {MFS (Sequence Boundary)};
\end{scope}

\end{tikzpicture}
\legende{Correlation panel for three wells in the Douala Basin (depth scale 1 cm = 100 m, datum = Base Cretaceous Unconformity). Sloping correlation lines reveal differential subsidence and sedimentation rates: the Aptian (Logbaba Fm) thickens basinward (Well A: 330 m, Well B: 500 m, Well C: 650 m), while the Albian (Mundeck Fm) shows even greater thickening. The Bentonite $\beta$ (112.5 Ma) and OAE1a (120 Ma) are key isochrons. The \emph{Orbitolina} range zone spans the Aptian.}
\end{popfigure}

\noindent\textbf{Interpretation:} The panel shows the Aptian \emph{Logbaba Formation} thickening significantly from onshore (Well A: 330 m) to offshore (Well C: 650 m), typical of a passive margin basin. The Albian \emph{Mundeck Formation} thickens even more. The Bentonite $\beta$ (112.5 Ma) and the OAE1a \ensuremath{\mathrm{\delta^{13}C}} excursion (120 Ma) provide absolute age tie-points. The \emph{Orbitolina} range zone spans the Aptian. The Top Mundeck MFS marks the Aptian/Albian boundary. This panel is the basis for a paleogeographic map and hydrocarbon play fairway analysis.

\begin{attention}[Data Consistency Check]
In real work, always verify sedimentation rates. Here, the Aptian (125--113 Ma, 12 Myr) is 330--650 m thick (27--54 m/Myr), reasonable for a passive margin. The Bentonite (112.5 Ma) lies correctly above OAE1a (120 Ma). The \emph{Orbitolina} LAD (Late Aptian) lies above the Bentonite (Late Aptian) in Wells A and B, but \emph{at the same level} in Well C (500 m). This suggests condensation or a hiatus near the Aptian/Albian boundary in the offshore well -- a critical observation for reservoir prediction.
\end{attention}

\end{exoresolu}

\courssub{Summary and Key Takeaways}

\begin{aretenir}[Correlation Essentials]
\begin{itemize}
\item Correlation establishes equivalence between stratigraphic units in different areas: lithostratigraphic (rock type), biostratigraphic (fossils), chronostratigraphic (time).
\item \textbf{Marker beds} (bentonites, coals, distinctive limestones) are the backbone of lithostratigraphic correlation.
\item \textbf{Index fossils} (wide geographic range, short time range, abundant) define biozones via FAD/LAD; range charts visualize overlap.
\item \textbf{Chronostratigraphic tools}: GSSPs, magnetostratigraphy (GPTS), chemostratigraphy (\ensuremath{\mathrm{\delta^{13}C}} excursions), astrochronology.
\item \textbf{Geophysical logs} (gamma-ray, resistivity) and \textbf{seismic reflectors} enable subsurface correlation.
\item \textbf{Facies changes} require correlation by time-parallel surfaces (MFS, sequence boundaries, biohorizons, chemostratigraphy), not lithology.
\item A \textbf{correlation panel} aligns sections on a datum, integrates all data types, and displays lateral thickness/facies changes.
\end{itemize}
\end{aretenir}

\begin{exercice}[Practice: Correlation in the Rio del Rey Basin]
Two wells (Well X and Well Y) in the Rio del Rey Basin penetrate a Cretaceous succession. Well X (proximal) shows: 0--200 m Sandstone (Barremian palynomorphs), 200--250 m Bentonite (Ar/Ar 125.2 Ma), 250--500 m Shale (Aptian foraminifera). Well Y (distal) shows: 0--400 m Shale (Barremian--Aptian foraminifera), 400--402 m Bentonite (Ar/Ar 125.2 Ma), 402--900 m Shale (Aptian--Albian foraminifera).
\begin{enumerate}
\item Draw a correlation panel at 1 cm = 100 m using the Bentonite as datum.
\item Identify the systems tracts and sequence boundary.
\item Explain why the sandstone in Well X does not appear in Well Y.
\item What chronostratigraphic tool would best constrain the Barremian/Aptian boundary here?
\end{enumerate}
\end{exercice}

\begin{corrige}[Exercise 8]
\begin{itemize}
\item \textbf{Panel}: Datum = Bentonite (horizontal line at y=0). Well X: Sandstone from y=-2.0 to 0 cm (above datum), Bentonite at 0, Shale 0 to +2.5 cm. Well Y: Shale from y=-4.0 to 0 cm (above datum), Bentonite at 0, Shale 0 to +5.0 cm. Correlation lines: Base of Sandstone in X (Sequence Boundary) correlates to a horizon within the shale in Y at y=-2.0 cm (correlative conformity). Top of Shale in X (y=+2.5) correlates to a horizon in Y (y=+2.5).
\item \textbf{Systems Tracts}: Well X: Lowstand Systems Tract (LST) Sandstone (Barremian) overlain by Transgressive Systems Tract (TST) Shale (Early Aptian). Well Y: Condensed TST Shale (Barremian--Early Aptian) overlying Highstand Systems Tract (not seen). The Bentonite (125.2 Ma, Early Aptian) falls within the TST.
\item \textbf{Facies change}: Proximal sandstone (shoreface/delta front) pinches out laterally into distal offshore shale. Same time (Barremian), different environment. The Bentonite (airfall ash) blankets both.
\item \textbf{Tool}: Magnetostratigraphy (M-sequence anomalies M0--M3) or chemostratigraphy (\ensuremath{\mathrm{\delta^{13}C}} curve) across the Barremian/Aptian boundary, calibrated to the GSSP. Palynomorphs give biozones but lower resolution for the boundary itself.
\end{itemize}
\end{corrige}

\coursec{Integration, Revision \& Exam Practice}

Having mastered the art of correlating successions across basins in Section 8, you now possess the full toolkit of a stratigrapher: relative principles, absolute dating, contact analysis, and correlation techniques. This final section is not merely a summary; it is a \emph{simulation of the examination room}. The GCE Advanced Level Geology paper (Paper 2, Practical) demands that you integrate disparate datasets — a field map, measured sections, fossil lists, and radiometric ages — into a single, coherent geological model within a strict time limit. Here we rehearse that integration, distil the syllabus into an exam-ready cheat sheet, dissect the marking schemes, and give you a personal audit checklist for the final weeks of revision.

\courssub{The Capstone Synthesis Exercise}

The ultimate test of your stratigraphic competence is the \emph{synthesis question}: given a realistic field dataset, produce a stratigraphic column, a correlation panel, and a written geological history. This mirrors Question 1 of the GCE Practical paper. We will work through a complete, exam-style scenario.

\begin{experience}[Capstone Synthesis: The "Mungo Basin" Field Problem]
\textbf{Scenario.} You are provided with the following data from a mapping exercise along the Mungo River corridor (inspired by the Cretaceous Mamfe Basin geology).

\textbf{Data Set A: Measured Section at Locality 1 (North Quarry).}
\begin{itemize}
    \item 0--12 m: Massive, cross-bedded sandstone, trough cross-bedding, paleocurrent $045^\circ$. Rare \emph{Trigonia} bivalves.
    \item 12--12.3 m: Sharp, erosional base. Pebble conglomerate (quartz, chert clasts up to 5 cm).
    \item 12.3--25 m: Dark grey shale with siderite nodules. Abundant ammonites: \emph{Desmoceras} sp., \emph{Douvilleiceras} sp. (Albian index fossils).
    \item 25--25.5 m: Bentonite layer (volcanic ash). \textbf{U-Pb zircon age: $102.4 \pm 0.6$ Ma.}
    \item 25.5--40 m: Bioturbated siltstone, \emph{Thalassinoides} trace fossils. \emph{Inoceramus} sp. (Cenomanian).
\end{itemize}

\textbf{Data Set B: Measured Section at Locality 2 (South Road Cut), 15 km SSW.}
\begin{itemize}
    \item 0--8 m: Lateritic soil / weathered sandstone.
    \item 8--20 m: Cross-bedded sandstone, identical lithology to Loc 1 base. \emph{Trigonia} present.
    \item 20--20.2 m: Conglomerate lag (same clast suite).
    \item 20.2--30 m: Dark shale. Ammonites: \emph{Desmoceras} sp. only.
    \item 30--30.1 m: Bentonite. \textbf{U-Pb zircon age: $102.6 \pm 0.5$ Ma.}
    \item 30.1--35 m: Siltstone with \emph{Inoceramus}. Top eroded.
\end{itemize}

\textbf{Data Set C: Borehole Log (Locality 3, 5 km East).}
\begin{itemize}
    \item 0--50 m: Sands and clays (Tertiary cover).
    \item 50--95 m: Shale with \emph{Inoceramus}, \emph{Ostrea}.
    \item 95--95.3 m: Bentonite. \textbf{K-Ar glauconite age: $94.2 \pm 1.2$ Ma.} (Note: K-Ar on glauconite is prone to argon loss; treat as minimum age).
    \item 95.3--120 m: Shale with \emph{Desmoceras}, \emph{Douvilleiceras}.
    \item 120--120.5 m: Conglomerate.
    \item 120.5--140 m: Sandstone (basal unit).
    \item 140 m: Basement granite (Pan-African, $\sim 550$ Ma).
\end{itemize}

\textbf{Task.} 
\begin{enumerate}
    \item Construct a composite stratigraphic column at 1:500 scale showing lithology, key surfaces, fossil zones, and radiometric dates with errors.
    \item Draw a correlation panel (fence diagram style) linking the three localities. Show the \emph{pinch-out} of the Cenomanian siltstone at Loc 2.
    \item Write a concise geological history (max 200 words) integrating relative and absolute data.
\end{enumerate}
\end{experience}

\begin{exoresolu}[Full Solution to Mungo Basin Synthesis]
\textbf{1. Composite Stratigraphic Column Construction.}
\begin{itemize}
    \item \textbf{Scale:} 1 cm = 5 m (1:500). Total thickness $\sim$ 40 m $\rightarrow$ 8 cm column height.
    \item \textbf{Lithostratigraphy (bottom to top):}
        \begin{itemize}
            \item \cle{Basal Sandstone Formation} (informal): Massive, trough cross-bedded sandstone, conglomerate base. Fluvial/estuarine. Thickness: 12 m (Loc 1), 12 m (Loc 2), 20 m (BH3).
            \item \cle{Mungo Shale Formation} (informal): Dark grey, fossiliferous shale, siderite nodules. Open marine transgression. Thickness: $\sim$ 13 m (Loc 1), $\sim$ 10 m (Loc 2), $\sim$ 25 m (BH3 — thickens basinward).
            \item \cle{Bentonite Marker Horizon}: 0.3--0.5 m thick. Key \emph{isochronous} marker. U-Pb ages overlap within error ($102.4 \pm 0.6$ vs $102.6 \pm 0.5$ Ma). Weighted mean: \textbf{$102.5 \pm 0.4$ Ma (Albian)}.
            \item \cle{Upper Siltstone Member}: Bioturbated siltstone, \emph{Inoceramus}. Cenomanian transgressive maximum. Present at Loc 1 (14.5 m) and BH3 (45 m), \textbf{absent at Loc 2 (erosional pinch-out/unconformity)}.
        \end{itemize}
    \item \textbf{Chronostratigraphy:} Albian (Lower--Middle) $\rightarrow$ Cenomanian (Lower). Bentonite = Early Albian ($\sim 102.5$ Ma). K-Ar date ($94.2$ Ma) is discordant (glauconite recoil/argon loss); true age of upper shale is $\sim 99$--$100$ Ma (Cenomanian). \emph{Always flag discordant dates.}
\end{itemize}

\textbf{2. Correlation Panel (Fence Diagram).}
\begin{itemize}
    \item Draw three vertical columns spaced horizontally (Loc 1 -- Loc 3 -- Loc 2).
    \item Correlate the \textbf{Bentonite} as a horizontal datum (red line). It is the primary correlation tie-point.
    \item Correlate the \textbf{Base Conglomerate} (sequence boundary / transgressive surface).
    \item Show the \textbf{Upper Siltstone} wedging out against the Loc 2 column; draw an unconformity surface (wavy line) truncating the Mungo Shale at Loc 2, overlain by Tertiary laterite.
    \item Label: Lithology patterns, Fossil Zones (Ammonite Biozones), Radiometric Ages ($\pm 2\sigma$), Unconformities (U1, U2).
\end{itemize}

\textbf{3. Geological History (Model Answer).}
\begin{quote}
\textbf{1. Pre-Cretaceous:} Pan-African basement (granite, $\sim 550$ Ma) peneplained. \\
\textbf{2. Early Albian ($\sim 103$ Ma):} Regional transgression. Deposition of Basal Sandstone (fluvial/estuarine) on uneven basement (onlap). \\
\textbf{3. Mid-Albian ($\sim 102.5$ Ma):} Maximum flooding. Deposition of Mungo Shale (open marine). Explosive volcanism deposits Bentonite marker bed (U-Pb zircon $102.5 \pm 0.4$ Ma). \emph{Desmoceras/Douvilleiceras} biozone. \\
\textbf{4. Late Albian -- Early Cenomanian ($\sim 100$ Ma):} Continued transgression. Deposition of Upper Siltstone (\emph{Inoceramus} zone). Thickens eastwards (Loc 3) indicating basin deepening towards the depocentre. \\
\textbf{5. Mid-Cenomanian ($\sim 98$ Ma):} Tectonic uplift / sea-level fall (Santonian? regional event). Erosion creates regional unconformity (U2). Upper Siltstone removed at Loc 2 (structural high), preserved in Loc 1 and Loc 3. \\
\textbf{6. Tertiary -- Recent:} CVL volcanism, lateritisation, Mungo River incision exposing sections.
\end{quote}
\textbf{Examiner's Tip:} Note the explicit link between the \textbf{Bentonite age} and the \textbf{Ammonite biozone} (calibration). Note the \textbf{rejection} of the K-Ar date with justification. Note the \textbf{basinward thickening} of the shale (subsidence analysis).
\end{exoresolu}

\begin{popfigure}
\begin{tikzpicture}[scale=0.7]
% Blank Stratigraphic Column Template for Student Practice
\def\colwidth{3}
\def\colheight{14}
% Main column box
\draw[popdark, thick] (0,0) rectangle (\colwidth,\colheight);
% Scale bar
\draw[popdark] (\colwidth+0.3,0) -- (\colwidth+0.3,\colheight);
\foreach \y in {0,2,4,6,8,10,12,14} {
    \draw[popdark] (\colwidth+0.2,\y) -- (\colwidth+0.4,\y) node[right] {\small \y m};
}
\node[rotate=90, anchor=south] at (\colwidth+0.8,7) {\textbf{Thickness (m)}};
% Column title
\node[align=center, popdark] at (\colwidth/2, \colheight+0.8) {\textbf{STRATIGRAPHIC COLUMN TEMPLATE}\\\small Scale: 1:500 (1 cm = 5 m)};
% Lithology column
\draw[popdark] (0.2,0) -- (0.2,\colheight);
\node[rotate=90, anchor=south, font=\small] at (0.1,7) {LITHO};
% Fossil column
\draw[popdark] (0.8,0) -- (0.8,\colheight);
\node[rotate=90, anchor=south, font=\small] at (0.7,7) {FOSSILS};
% Age column
\draw[popdark] (1.6,0) -- (1.6,\colheight);
\node[rotate=90, anchor=south, font=\small] at (1.5,7) {AGE (Ma)};
% Contacts / Surfaces column
\draw[popdark] (2.4,0) -- (2.4,\colheight);
\node[rotate=90, anchor=south, font=\small] at (2.3,7) {CONTACTS};
% Legend box
\draw[popdark] (4,10) rectangle (8.5,14);
\node[align=left, font=\small, text width=4cm] at (6.25,13.5) {
\textbf{Lithology Key:}\par
\hspace{0.5cm}\textcolor{popblueL}{\rule{0.4cm}{0.2cm}} Sandstone\par
\hspace{0.5cm}\textcolor{popmuted}{\rule{0.4cm}{0.2cm}} Shale / Mudstone\par
\hspace{0.5cm}\textcolor{popgold}{\rule{0.4cm}{0.2cm}} Siltstone\par
\hspace{0.5cm}\textcolor{poporange}{\rule{0.4cm}{0.2cm}} Conglomerate\par
\hspace{0.5cm}\textcolor{popgreenL}{\rule{0.4cm}{0.2cm}} Limestone\par
\hspace{0.5cm}\textcolor{popdark}{\rule[0.1cm]{0.4cm}{0.05cm}} Unconformity (wavy)\par
\hspace{0.5cm}\textcolor{popred}{\rule[0.1cm]{0.4cm}{0.05cm}} Marker Bed (Bentonite/Tuff)
};
% Instruction text
\node[align=center, font=\small\itshape, text width=10cm] at (4.5,5) {
\textbf{Student Practice:} Print this template.\par
Using the Mungo Basin data (Section 9.1), draw the composite column.\par
Fill in lithology patterns, fossil names, radiometric ages with errors,\par
and contact types (S = sharp, E = erosional, G = gradational).
};
\end{tikzpicture}
\legende{Blank stratigraphic column template (1:500 scale) for examination practice. Reproduce this in your practical notebook.}
\end{popfigure}

\courssub{The "Cheat Sheet": Definitions, Principles \& Formulas}

In the exam, you waste precious seconds recalling definitions. Memorise this condensed reference. It covers every lesson (2--22) of the syllabus.

\begin{aretenir}[Stratigraphy Cheat Sheet — Upper Sixth GCE]

\textbf{A. Scope \& Principles (Lessons 2--4)}
\begin{itemize}
    \item \cle{Stratigraphy}: Study of stratified rocks: their geometry, composition, origin, age, and correlation.
    \item \cle{Law of Superposition}: In an undeformed sequence, oldest at bottom, youngest at top.
    \item \cle{Principle of Original Horizontality}: Sediments deposit horizontally.
    \item \cle{Principle of Lateral Continuity}: Layers extend laterally until they thin out or hit a barrier.
    \item \cle{Principle of Cross-Cutting Relationships}: A feature cutting another is younger.
    \item \cle{Principle of Inclusions}: Clasts in a rock are older than the host rock.
    \item \cle{Principle of Faunal Succession (Smith)}: Fossil assemblages succeed each other vertically in a definite, determinable order.
    \item \cle{Uniformitarianism}: "The present is the key to the past" (Hutton/Lyell). Rates may vary.
    \item \cle{Actualism}: Modern processes \emph{and} rates explain the past (modern refinement).
\end{itemize}

\textbf{B. Geologic Time Scale (Lessons 5--7)}
\begin{itemize}
    \item \cle{Eon > Era > Period > Epoch > Age} (Chronostratigraphic: Eonothem > Erathem > System > Series > Stage).
    \item \cle{Phanerozoic}: 541 Ma -- Present. \cle{Palaeozoic} (541--252 Ma), \cle{Mesozoic} (252--66 Ma), \cle{Cenozoic} (66 Ma--Present).
    \item \cle{Key Boundaries}: P-Tr (252 Ma, mass extinction), K-Pg (66 Ma, Chicxulub), P-E (56 Ma, PETM).
    \item \cle{Cameroon Context}: CVL (Cretaceous--Present), Mamfe Basin (Cretaceous), Pan-African Basement ($>541$ Ma).
\end{itemize}

\textbf{C. Dating Methods (Lessons 8--13)}
\begin{itemize}
    \item \cle{Relative Dating}: Superposition, fossils, cross-cutting, unconformities. \emph{No numbers}.
    \item \cle{Absolute (Radiometric) Dating}: Parent $\rightarrow$ Daughter + Radiation. $N = N_0 e^{-\lambda t}$.
    \item \cle{Half-life}: $t_{1/2} = \ln 2 / \lambda$.
    \item \cle{Age Equation}: $t = \frac{1}{\lambda} \ln\left(1 + \frac{D}{P}\right)$.
    \item \cle{Isochron Method}: Plots $D^*/D_{ref}$ vs $P/D_{ref}$. Slope $= e^{\lambda t} - 1$. \textbf{Assumptions:} Closed system, initial homogeneity, constant $\lambda$.
    \item \cle{Key Systems}:
    \begin{center}
    \begin{tabular}{lccc}
    \textbf{System} & \textbf{Half-life} & \textbf{Range} & \textbf{Material} \\ \hline
    U-Pb (Zircon) & $4.47 / 0.704\ \mathrm{Ga}$ & $>1\ \mathrm{Ma}$ -- $4.5\ \mathrm{Ga}$ & Zircon, Titanite, Baddeleyite \\
    K-Ar / Ar-Ar & $1.25\ \mathrm{Ga}$ & $>10\ \mathrm{ka}$ -- $4.5\ \mathrm{Ga}$ & K-feldspar, Mica, \textbf{Glauconite*}, Volcanics \\
    Rb-Sr & $48.8\ \mathrm{Ga}$ & $>10\ \mathrm{Ma}$ & Whole rock, Mica \\
    Sm-Nd & $106\ \mathrm{Ga}$ & $>10\ \mathrm{Ma}$ & Whole rock, Garnet \\
    C-14 & $5730\ \mathrm{yr}$ & $<50\ \mathrm{ka}$ & Organic carbon, Charcoal \\
    \end{tabular}
    \end{center}
    \item \cle{*Glauconite K-Ar}: Prone to argon loss $\rightarrow$ \textbf{minimum ages only}. Flag in exams.
    \item \cle{Concordia/Discordia} (Wetherill): U-Pb on zircon. Concordant = closed system. Discordia chord $\rightarrow$ upper intercept = crystallisation age, lower intercept = Pb-loss event age.
\end{itemize}

\textbf{D. Contacts \& Unconformities (Lessons 14--15)}
\begin{itemize}
    \item \cle{Conformable}: Continuous deposition (Gradational, Sharp).
    \item \cle{Unconformity}: Hiatus (gap in record). Types:
    \begin{itemize}
        \item \cle{Angular Unconformity}: Tilted older rocks truncated, overlain by horizontal younger.
        \item \cle{Disconformity}: Parallel beds, erosional surface (palaeosol, lag), hard to spot.
        \item \cle{Nonconformity}: Sediments on crystalline basement.
        \item \cle{Paraconformity}: Parallel beds, no erosion, gap shown only by fossils.
        \item \cle{Buttress Unconformity}: Onlap against paleotopography (reef, fault scarp).
    \end{itemize}
    \item \cle{Sequence Boundary}: Regional unconformity = base of a depositional sequence (Vail/Exxon).
    \item \cle{Transgressive Surface (TS)} / \cle{Maximum Flooding Surface (MFS)}: Key correlation surfaces.
\end{itemize}

\textbf{E. Sequencing Events (Lessons 16--19)}
\begin{itemize}
    \item \textbf{Algorithm:} 1. Identify units (oldest $\rightarrow$ youngest). 2. Identify structures (folds, faults, intrusions). 3. Apply Cross-Cutting. 4. Identify unconformities (erosion events). 5. Write history: "Deposition of X $\rightarrow$ Folding $\rightarrow$ Erosion (U1) $\rightarrow$ Deposition of Y $\rightarrow$ Faulting $\rightarrow$ Present".
    \item \cle{Inclusion Rule}: Xenoliths in granite $>$ Granite age. Clasts in conglomerate $>$ Conglomerate age.
    \item \cle{Baked/Chilled Margins}: Contact metamorphism proves intrusion is younger than country rock.
\end{itemize}

\textbf{F. Correlation (Lessons 20--22)}
\begin{itemize}
    \item \cle{Lithostratigraphic}: Rock type (formation, member). \emph{Diachronic} (time-transgressive).
    \item \cle{Biostratigraphic}: Fossils (Zone, Subzone). \emph{Isochronous} (time-equivalent) if facies-independent. \cle{Index Fossil}: Wide geographic range, narrow time range, abundant, easily ID'd.
    \item \cle{Chronostratigraphic}: Time-rock units (System, Stage). Global standard (GSSP = Golden Spike).
    \item \cle{Geochronologic}: Time units (Period, Age). "Early Jurassic" (time) vs "Lower Jurassic" (rocks).
    \item \cle{Magnetostratigraphy}: Polarity chrons (C1n, C2An, etc.). Global isochrons.
    \item \cle{Chemostratigraphy}: $\delta^{13}C$, $\delta^{18}O$, $^{87}Sr/^{86}Sr$ curves. Global correlation.
    \item \cle{Event Stratigraphy}: Volcanic ash (bentonite), impact layer (Ir anomaly), OAE (black shale).
\end{itemize}
\end{aretenir}

\courssub{Deconstructing GCE Question Types}

Understanding the \emph{command words} and mark allocation is half the battle.

\begin{propriete}[GCE Marking Scheme Anatomy]
The GCE Geology Practical (Paper 2) marking scheme typically allocates marks as follows:
\begin{itemize}
    \item \textbf{Diagrams (30--40\%):} Strat columns, correlation panels, cross-sections, stereonets. \textbf{Marks for:} Scale, Title, Legend, Correct patterns, Labels (units, ages, contacts), Neatness (pencil, ruler).
    \item \textbf{Calculations (15--20\%):} Radiometric ages, sedimentation rates, dip/apparent thickness. \textbf{Marks for:} Formula stated, correct substitution, units, significant figures (usually 2--3), error propagation.
    \item \textbf{Short Answers / Definitions (20--25\%):} "Define...", "State...", "Identify...". \textbf{Marks for:} Key terms (e.g., "closed system", "index fossil", "hiatus").
    \item \textbf{Essay / Geological History (20--25\%):} "Write a geological history...", "Discuss the evolution...". \textbf{Marks for:} Logical sequence, integration of relative + absolute data, use of technical vocabulary, identification of hiatuses, tectonic vs eustatic controls.
\end{itemize}
\textbf{Rule:} A clear, labelled diagram often earns full marks even if the accompanying text is brief. A perfect essay with no diagram loses the diagram marks.
\end{propriete}

\begin{popfigure}
\begin{tikzpicture}
\begin{axis}[
    width=14cm, height=7cm,
    ybar stacked,
    bar width=12pt,
    xlabel={Question Type (Typical Paper 2)},
    ylabel={Percentage of Total Marks (\%)},
    symbolic x coords={Strat Column, Correlation Panel, Cross-Section, Radiometric Calc, Sed Rate Calc, Definitions/Short Ans, Geol History Essay, Map Interpretation},
    xtick=data,
    x tick label style={rotate=45, anchor=east, font=\small},
    ymin=0, ymax=100,
    legend style={at={(0.5,-0.25)}, anchor=north, legend columns=2, font=\small},
    title={Typical GCE Practical Paper Mark Distribution (Indicative)},
    title style={font=\bfseries, yshift=0.3cm},
    ymajorgrids=true,
    grid style={popline},
]
\addplot[fill=popblue] coordinates {(Strat Column,15) (Correlation Panel,12) (Cross-Section,10) (Radiometric Calc,10) (Sed Rate Calc,5) (Definitions/Short Ans,20) (Geol History Essay,18) (Map Interpretation,10)};
\addplot[fill=poporange, draw=none, forget plot] coordinates {(Strat Column,0) (Correlation Panel,0) (Cross-Section,0) (Radiometric Calc,0) (Sed Rate Calc,0) (Definitions/Short Ans,0) (Geol History Essay,0) (Map Interpretation,0)};
\legend{Marks Allocated, }
\end{axis}
\end{tikzpicture}
\legende{Indicative mark breakdown for a GCE Advanced Level Geology Practical Paper (Paper 2). Percentages vary yearly but the hierarchy is stable: Diagrams + Calculations + Definitions + Essay.}
\end{popfigure}

\begin{attention}[Essay Trap: The "Definition Penalty"]
In essay questions (e.g., "Discuss the role of unconformities in basin evolution"), \textbf{the first 2--3 marks are almost always for definitions}. 
\begin{itemize}
    \item \textbf{Bad start:} "Unconformities show gaps..." $\rightarrow$ 0/3 for definitions.
    \item \textbf{Good start:} "An \textbf{unconformity} is a surface of erosion or non-deposition representing a \textbf{hiatus} in the geological record. An \textbf{angular unconformity} separates tilted older strata from horizontal younger strata. A \textbf{disconformity}..." $\rightarrow$ 3/3 for definitions. \emph{Then} discuss basin evolution.
\end{itemize}
\textbf{Examiner's Report Quote (paraphrased):} "Candidates who defined terms before discussing scored significantly higher. Many wrote excellent geological narratives but forgot to define 'sequence boundary' or 'maximum flooding surface', losing easy marks."
\end{attention}

\courssub{Time Management: The 2-Hour Practical Paper Strategy}

Paper 2 is 2 hours (120 minutes). A typical paper has 4--5 questions. \textbf{Do not} start at Question 1 and work linearly if Q1 is a 30-minute column drawing and Q4 is a 5-minute definition. Use this protocol:

\begin{methode}[Exam Room Protocol: 120 Minutes]
\begin{enumerate}
\item \textbf{Minutes 0--5: READ EVERYTHING.} Read the whole paper. Identify the "big" questions (strat column, correlation, history) and the "quick" questions (definitions, calculations, map symbols). Note the mark allocation printed on the paper.
\item \textbf{Minutes 5--10: PLAN \& SET UP.} Sharpen pencils (2H for lines, HB for shading). Draw your \textbf{stratigraphic column template} (like Fig. 9.1) on graph paper immediately. Draw your \textbf{correlation panel framework} (three vertical lines, datum line). This "scaffolding" saves 10 minutes later.
\item \textbf{Minutes 10--50: THE "BIG DRAW" (Strat Column + Correlation).} This is usually 30--40 marks. Work systematically: 
    \begin{enumerate}
        \item Plot all contacts/measured thicknesses to scale.
        \item Add lithology patterns (use standard symbols: dots=sand, dashes=shale, bricks=conglom).
        \item Add fossil names \emph{next to} the units.
        \item Add radiometric ages with $\pm$ errors \emph{at the correct horizon}.
        \item Draw unconformities as \textbf{wavy} lines (erosional) or \textbf{straight} lines with hachures (nonconformity).
        \item Correlation: Join the \textbf{isochronous horizons} (bentonite, MFS, sequence boundary) first. Then lithology.
    \end{enumerate}
\item \textbf{Minutes 50--70: CALCULATIONS.} Radiometric ages, sedimentation rates. Show \textbf{every step}. Formula $\rightarrow$ Substitution $\rightarrow$ Answer $\rightarrow$ Unit $\rightarrow$ Error.
\item \textbf{Minutes 70--90: SHORT ANSWERS / MAP WORK.} Definitions, cross-sections, stereonets, map symbols. These are high marks-per-minute.
\item \textbf{Minutes 90--110: GEOLOGICAL HISTORY ESSAY.} Use your column/correlation as the script. Write in \textbf{bullet-point chronological order}. Use past tense ("Deposition of...", "Folding occurred...", "Intrusion of..."). Explicitly mention \textbf{hiatus durations} if calculable.
\item \textbf{Minutes 110--120: REVIEW (The "Exam-Ready Checklist").} See Box 9.5.
\end{enumerate}
\end{methode}

\courssub{The Exam-Ready Checklist (15-Point Final Audit)}

Before you hand in your script, run this mental (or written) checklist. Every item is a common cause of lost marks.

\begin{methode}[Exam-Ready Checklist — Verify Before Submission]
\begin{enumerate}
\item \textbf{Scale Bar:} Does every diagram (column, panel, cross-section) have a graphic scale bar (0--100 m / 1:500)?
\item \textbf{Title Block:} "Composite Stratigraphic Column, Mungo Basin, Localities 1--3, Scale 1:500, Candidate Name/Number".
\item \textbf{Lithology Legend:} Is there a key explaining \emph{every} pattern used (sand, shale, lime, conglom, tuff, basement)?
\item \textbf{Contact Types:} Are contacts drawn correctly? Sharp (straight line), Gradational (zigzag), Unconformity (wavy), Fault (thick line + displacement arrows).
\item \textbf{Radiometric Errors:} Are ages written as $102.5 \pm 0.4$ Ma (not just "102.5 Ma")? Is the \textbf{method} stated (U-Pb zircon)?
\item \textbf{Discordant Dates Flagged:} Did you note "K-Ar glauconite: minimum age, likely argon loss"?
\item \textbf{Fossil Zones:} Are biozones labelled (e.g., "\emph{Desmoceras} Zone, Albian") not just fossil names?
\item \textbf{Correlation Lines:} Do tie-lines connect \textbf{isochronous surfaces} (bentonite, MFS, sequence boundary), not just similar lithologies?
\item \textbf{Pinch-outs/Onlap:} Is the wedge geometry of the Upper Siltstone shown correctly at Loc 2?
\item \textbf{North Arrow / Dip Symbols:} On map/cross-section: North arrow, strike/dip symbols, vertical exaggeration stated (e.g., VE = 5x).
\item \textbf{Calculation Working:} Formula $\rightarrow$ Numbers $\rightarrow$ Result $\rightarrow$ Unit $\rightarrow$ Sig Figs (usually 2 for ages, 3 for rates).
\item \textbf{History Chronology:} Oldest event first. Every deposition, deformation, erosion, intrusion, unconformity event listed.
\item \textbf{Hiatus Quantification:} "Hiatus of $\sim 4$ Myr (Cenomanian--Turonian missing)" — estimated from biostratigraphy/radiometrics.
\item \textbf{Terminology Precision:} "Lower Jurassic" (rocks) vs "Early Jurassic" (time). "Formation" (litho) vs "Stage" (chrono). "Angular unconformity" not just "unconformity".
\item \textbf{Neatness:} Pencil lines sharp, shading even, no smudges, handwriting legible. \textbf{Diagrams in pencil, text in pen.}
\end{enumerate}
\end{methode}

\courssub{Self-Assessment: The 22-Lesson Audit}

Rate your confidence (1 = Weak, 5 = Exam Ready) for each syllabus lesson. Target: All $\ge 4$ two weeks before the exam.

\begin{center}
\begin{tabularx}{\linewidth}{|c|X|c|}\hline
\textbf{Lesson} & \textbf{Topic} & \textbf{Confidence (1--5)} \\ \hline
2 & Scope of Stratigraphy & \\ \hline
3 & Stratigraphic Principles 1 (Superposition, Horizontality, Continuity) & \\ \hline
4 & Stratigraphic Principles 2 (Cross-cutting, Inclusions, Faunal Succession) & \\ \hline
5 & Geologic Time Scale 1 (Hierarchy, Phanerozoic divisions) & \\ \hline
6 & Geologic Time Scale 2 (Boundary events, GSSPs, Cameroon context) & \\ \hline
7 & Geologic Time Scale 3 (Magnetostratigraphy, Chemostratigraphy basics) & \\ \hline
8 & Dating 1 (Relative vs Absolute, Half-life, Decay equation) & \\ \hline
9 & Dating 2 (K-Ar, Ar-Ar, Rb-Sr, Sm-Nd systems) & \\ \hline
10 & Dating 3 (U-Pb, Concordia/Discordia, Zircon) & \\ \hline
11 & Dating 4 (C-14, Cosmogenic, Luminescence) & \\ \hline
12 & Dating 5 (Isochron theory, Assumptions, Open/Closed systems) & \\ \hline
13 & Dating 6 (Error propagation, Geochronology vs Thermochronology) & \\ \hline
14 & Stratigraphic Contacts 1 (Conformities, Unconformity types) & \\ \hline
15 & Stratigraphic Contacts 2 (Sequence boundaries, TS, MFS, Paraconformity) & \\ \hline
16 & Sequence of Events 1 (Cross-section interpretation, Algorithm) & \\ \hline
17 & Sequence of Events 2 (Inclusions, Baked margins, Xenoliths) & \\ \hline
18 & Sequence of Events 3 (Complex histories: multiple deformations) & \\ \hline
19 & Sequence of Events 4 (Integrating radiometric dates into history) & \\ \hline
20 & Correlation 1 (Litho- vs Bio- vs Chronostratigraphy) & \\ \hline
21 & Correlation 2 (Magnetostratigraphy, Chemostratigraphy, Event strat) & \\ \hline
22 & Correlation 3 / Integration (Synthesis: Column + Panel + History) & \\ \hline
\end{tabularx}
\end{center}

\textbf{Action:} For any lesson $\le 3$, re-read the corresponding course section, do the \texttt{exercice} boxes, and check the \texttt{corrige}. Do \textbf{not} just re-read notes; do active recall (draw column from memory, calculate an age, define 5 terms).

\courssub{Further Reading \& Resources}

\begin{savaistu}[Beyond the Syllabus: Tools for the Professional Stratigrapher]
\textbf{1. International Chronostratigraphic Chart (ICS).} Download the latest PDF from \texttt{stratigraphy.org}. It is the \textbf{global standard}. Know the current ratified GSSPs (Golden Spikes). The GCE expects you to use the ICS timescale (e.g., Cretaceous = 145.0--66.0 Ma, not old values).
\textbf{2. "Principles of Sedimentology and Stratigraphy" by Sam Boggs Jr. (Pearson).} The "Bible" for this level. Chapters 1--3, 12--14 cover our syllabus perfectly. Clear diagrams of unconformities, sequence stratigraphy, correlation methods.
\textbf{3. "A Geologic Time Scale 2020" (Gradstein et al., Elsevier).} The definitive reference for boundary ages and error margins. Use for checking "current" ages in revision.
\textbf{4. Past GCE Papers (Cameroon GCE Board).} \textbf{Non-negotiable.} Do at least 5 full Paper 2 past papers (2018--2023) \textbf{under timed conditions}. Mark them yourself using the official marking schemes (often available from teachers or GCE Board archives). Analyse \emph{why} you lost marks.
\textbf{5. Cameroon Geology Specifics:} 
\begin{itemize}
    \item \emph{Nguimbous-Kouoh et al. (2019)} "Geology of the Mamfe Basin" — for Cretaceous stratigraphy, ammonite zones, bentonites.
    \item \emph{Marzoli et al. (1999)} "CVL volcanism" — for U-Pb dates on CVL basalts.
    \item \emph{IRGM/Ministry of Mines} geological maps (1:500,000 or 1:200,000 sheets). Practice reading the legend and cross-sections.
\end{itemize}
\textbf{6. Digital Tools:} \texttt{TimeScale Creator} (free Java app, \texttt{tscreator.org}) — build your own custom time charts. \texttt{Strater} or \texttt{LogPlot} (industry software) — for digital column practice (optional but impressive for university).
\end{savaistu}

\begin{definition}[Integration]
\textbf{Integration} in stratigraphy is the cognitive and technical process of combining \textbf{relative data} (superposition, fossil succession, cross-cutting, unconformities) with \textbf{absolute data} (radiometric ages, magnetostratigraphy, chemostratigraphy) to construct a \textbf{single, self-consistent, four-dimensional (3D space + time) geological model} of a basin or region. It requires resolving conflicts (e.g., discordant dates), identifying hiatuses, and distinguishing allocyclic (eustasy, tectonics) from autocyclic (channel migration) controls. \emph{This is the highest-order skill examined at GCE A-Level.}
\end{definition}

\courssub{Final Words: From Student to Geologist}

You have now traversed the entire Upper Sixth Stratigraphy syllabus. You can read the Earth's sedimentary archive: you know how to order its pages (principles), date its chapters (geochronology), recognise its missing pages (unconformities), correlate its volumes across the library (correlation), and write its history (synthesis). 

The examination is not a trap; it is a \emph{performance}. You have the script (this course), you have rehearsed the scenes (worked examples), and you know the marking criteria. On the day:
\begin{itemize}
    \item \textbf{Breathe.} Read the question twice.
    \item \textbf{Draw first.} Diagrams anchor your thinking and secure "easy" marks.
    \item \textbf{Define before you discuss.} 
    \item \textbf{Integrate.} Never give a relative age when an absolute age is available, and vice versa. Use both.
    \item \textbf{Enjoy it.} You are doing real geology.
\end{itemize}

\textbf{Bonne chance, et que la stratigraphie soit avec vous!} (Good luck, and may the stratigraphy be with you!)

\begin{exercice}[Final Integration Drill — Timed (45 minutes)]
\textbf{Scenario:} A new borehole (BH-4) in the Mungo Basin (Section 9.1) yields:
\begin{itemize}
    \item 0--30 m: Tertiary sands.
    \item 30--65 m: Shale with \emph{Inoceramus}, \emph{Ostrea}, \emph{Scaphites} (Turonian index).
    \item 65--65.2 m: Bentonite. \textbf{U-Pb zircon: $93.1 \pm 0.3$ Ma.}
    \item 65.2--80 m: Shale with \emph{Desmoceras} (Albian).
    \item 80--80.5 m: Conglomerate.
    \item 80.5--95 m: Sandstone.
    \item 95 m: Basement.
\end{itemize}
\textbf{Tasks (Exam Conditions):}
\begin{enumerate}
    \item Add BH-4 to your Correlation Panel (Section 9.1). Show the new Turonian unit.
    \item Calculate the \textbf{sedimentation rate} ($\mathrm{m/Myr}$) for the Albian Shale (Bentonite at 102.5 Ma $\rightarrow$ Base Turonian at 93.1 Ma). Thickness in BH-3 = 25 m (95.3--120 m); in BH-4 = 15 m (65.2--80 m). Discuss the difference.
    \item Update the Geological History to include the Turonian transgression and the new bentonite age.
    \item \textbf{Critical Thinking:} The K-Ar glauconite date from BH-3 (94.2 Ma) falls \emph{between} the two U-Pb bentonite dates (102.5 Ma and 93.1 Ma). Is the K-Ar date reliable? Justify.
\end{enumerate}
\end{exercice}

\begin{corrige}[Final Integration Drill]
\textbf{1. Correlation Panel Update.} Add BH-4 column between BH-3 and Loc 1. New unit: \cle{Turonian Shale} (35 m thick in BH-4) with \emph{Scaphites} biozone. New Bentonite (93.1 Ma) correlates to a bentonite \emph{not seen} in other sections (erosion/non-deposition). Draw it as a dashed correlation line? No — U-Pb is robust. Draw solid line, label "Bentonite B (93.1 Ma)". Note: Turonian unit \textbf{absent} at Loc 1, Loc 2, BH-3 $\rightarrow$ major unconformity (U3) at top of Cenomanian/Albian.

\textbf{2. Sedimentation Rates.}
\begin{itemize}
    \item \textbf{BH-3 (Albian Shale):} Thickness = 25 m. Time = $102.5 - 93.1 = 9.4\ \mathrm{Myr}$. Rate = $25 / 9.4 = \mathbf{2.7\ \mathrm{m/Myr}}$.
    \item \textbf{BH-4 (Albian Shale):} Thickness = 15 m. Same time span (9.4 Myr). Rate = $15 / 9.4 = \mathbf{1.6\ \mathrm{m/Myr}}$.
    \item \textbf{Interpretation:} BH-3 is closer to the \textbf{depocentre} (basin axis) $\rightarrow$ higher subsidence $\rightarrow$ thicker accumulation. BH-4 is on a \textbf{structural high} or basin margin $\rightarrow$ lower accommodation space. Consistent with BH-3 having thicker Cenomanian (45 m vs 14.5 m at Loc 1).
\end{itemize}

\textbf{3. Updated History Snippet.}
\begin{quote}
\textbf{5. Early Turonian ($\sim 93.5$ Ma):} Renewed transgression. Deposition of Turonian Shale (\emph{Scaphites} zone). Volcanic event: Bentonite B (U-Pb $93.1 \pm 0.3$ Ma). \\
\textbf{6. Mid-Late Turonian ($\sim 92$ Ma):} Major regression / uplift. Regional unconformity U3 removes Turonian at Loc 1, Loc 2, BH-3. Preserved only in BH-4 (depocentre).
\end{quote}

\textbf{4. K-Ar Glauconite Reliability.}
\textbf{No, it is not reliable.} The K-Ar age ($94.2 \pm 1.2$ Ma) is \textbf{younger} than the underlying Albian bentonite ($102.5$ Ma) but \textbf{older} than the overlying Turonian bentonite ($93.1$ Ma). This is \textbf{stratigraphically impossible} for a closed system (superposition violated). Glauconite forms slowly at the sediment-water interface and continues to incorporate K (and lose Ar) long after deposition. The date represents a \textbf{mixed age} (minimum age of authigenesis $\approx$ Late Albian/Early Cenomanian). \textbf{Always trust U-Pb zircon over K-Ar glauconite for depositional age.}
\end{corrige}

\begin{aretenir}[Last-Minute "Pocket" Summary — Print on A6 Card]
\begin{itemize}
    \item \textbf{Superposition} = Relative age order.
    \item \textbf{Unconformity} = Hiatus (Time gap). 4 Types: Angular, Disconformity, Nonconformity, Paraconformity.
    \item \textbf{Index Fossil} = Wide geography, Short time, Abundant, Distinctive.
    \item \textbf{Radiometric Age} = $t = \frac{1}{\lambda}\ln(1+D/P)$. \textbf{Closed System} essential.
    \item \textbf{U-Pb Zircon} = Gold standard (Concordia). \textbf{K-Ar Glauconite} = Minimum age (Leaky).
    \item \textbf{Correlation Hierarchy:} Event (Bentonite) $>$ Biozone $>$ Magnetozone $>$ Chemozone $>$ Lithology.
    \item \textbf{Sequence Stratigraphy:} SB (Sequence Boundary) $\rightarrow$ LST $\rightarrow$ TST $\rightarrow$ MFS $\rightarrow$ HST $\rightarrow$ SB.
    \item \textbf{Geological History:} Oldest $\rightarrow$ Youngest. Event + Process + Age + Evidence.
    \item \textbf{Exam:} Pencil diagrams (Scale, Title, Legend, Labels). Pen text. Define $\rightarrow$ Discuss. Check Checklist.
\end{itemize}
\end{aretenir}

\newpage

\coursec{Integration activity}

\coursec{Integration, Revision \& Exam Practice}
\courssub{Integration Exercise: Douala Basin Stratigraphic Synthesis}

\begin{experience}[Douala Basin Stratigraphic Synthesis]
\courssub{Aims of the activity}
\begin{itemize}
    \item Consolidate the full suite of stratigraphic skills taught in this chapter: principles (superposition, original horizontality, cross‑cutting, faunal succession), the geologic time scale, absolute and relative dating methods, recognition of stratigraphic contacts and unconformities, sequencing of geologic events, and correlation of successions.
    \item Apply these skills to a realistic dataset from the \cle{Douala Basin} (Cameroon), a sedimentary basin that records the Cretaceous–Cenozoic evolution of the \cle{Cameroon Volcanic Line} (CVL) and the Atlantic margin.
    \item Produce a professional‑style stratigraphic report including lithostratigraphic columns, a correlation panel, a sequential geologic history, and an exploration‑oriented interpretation (hydrocarbon trap prediction).
    \item Develop teamwork, scientific communication (poster presentation), and critical evaluation of data quality — competencies required for the GCE Advanced Level practical examination and for future university geoscience studies.
\end{itemize}

\courssub{Situation}
\textbf{Geological setting.} The Douala Basin is a coastal sedimentary basin in southwestern Cameroon, bounded to the north by the \cle{Congo Craton}, to the west by the Atlantic Ocean, and to the east by the \cle{Cameroon Volcanic Line} (CVL). It contains a thick Cretaceous to Recent succession deposited in fluvial, deltaic, and shallow marine environments, interrupted by volcanic episodes linked to the CVL. The basin is a current target for hydrocarbon exploration.

\textbf{Data packet provided to each group (3–4 students).}
\begin{enumerate}
    \item \textbf{Simplified geological map (sketch).} Shows the basin outline, major faults (e.g., the \cle{Sanaga Fault}), and the location of three measured sections:
        \begin{itemize}
            \item \textbf{Section A (Kribi area)} – coastal outcrop, 120 m thick.
            \item \textbf{Section B (Douala city outskirts)} – road‑cut exposure, 210 m thick.
            \item \textbf{Section C (Limbe area)} – cliff section near Mount Cameroon, 180 m thick.
        \end{itemize}
    \item \textbf{Measured sections (lithology, thickness, sedimentary structures).} Summarised in Table~1.
    \item \textbf{Fossil assemblage lists} for each section (Table~2). Key index fossils: \emph{Orbitolina} spp. (Early Cretaceous, Albian), \emph{Rotalipora} spp. (Cenomanian–Turonian), \emph{Globotruncana} spp. (Campanian–Maastrichtian), \emph{Nummulites} spp. (Paleocene–Eocene), \emph{Miogypsina} spp. (Miocene).
    \item \textbf{Radiometric ages (absolute dates)} from interbedded volcanic layers:
        \begin{itemize}
            \item \textbf{U–Pb zircon} from a tuff layer in Section B (at 85 m from base): \SI{93.5 \pm 0.7}{Ma}.
            \item \textbf{K–Ar} from a basalt flow capping Section C (top 10 m): \SI{28.3 \pm 1.2}{Ma}.
            \item \textbf{Ar–Ar} from a volcanic ash layer in Section A (at 40 m from base): \SI{55.2 \pm 0.4}{Ma}.
        \end{itemize}
    \item \textbf{Seismic line sketch (interpreted).} A single NW–SE seismic profile crossing the three sections (Figure~1 description). Key reflectors: 
        \begin{itemize}
            \item R1 – strong, continuous reflector at ~2.5 s TWT (base of Cretaceous).
            \item R2 – moderately continuous, onlaps onto R1 (Cenomanian transgression).
            \item R3 – high‑amplitude, chaotic zone (volcanic basement / CVL intrusions).
            \item R4 – regional unconformity (Eocene–Oligocene) truncating R2.
            \item R5 – prograding clinoforms (Miocene deltaic sequence).
        \end{itemize}
\end{enumerate}

\begin{center}
\begin{tabularx}{\linewidth}{|l|X|X|X|}
\hline
\rowcolor{popblueL}
\textbf{Unit (base to top)} & \textbf{Section A (Kribi)} & \textbf{Section B (Douala)} & \textbf{Section C (Limbe)} \\
\hline
1. Basal conglomerate & 5 m, clast‑supported, quartzite clasts & 8 m, matrix‑supported, feldspathic & 6 m, polymictic, volcanic clasts \\
2. Sandstone (fluvial) & 20 m, cross‑bedded, trough cross‑strat. & 30 m, planar cross‑beds, channel fills & 25 m, ripple‑laminated, minor mud drapes \\
3. Shale (lacustrine) & 15 m, laminated, organic‑rich & 25 m, fissile, pyritic & 20 m, silty, plant fragments \\
4. Limestone (shallow marine) & 30 m, bioclastic, \emph{Orbitolina} & 40 m, wackestone, \emph{Rotalipora} & 35 m, packstone, \emph{Globotruncana} \\
5. Marl (transitional) & 10 m, fossiliferous, \emph{Nummulites} & 15 m, glauconitic, \emph{Nummulites} & 12 m, foraminiferal, \emph{Miogypsina} \\
6. Sandstone (deltaic) & 25 m, hummocky cross‑strat. & 40 m, sigmoidal clinoforms & 30 m, distributary channel sands \\
7. Volcanic interbeds & Tuff at 40 m (Ar–Ar) & Tuff at 85 m (U–Pb) & Basalt flow at top 10 m (K–Ar) \\
8. Cover (Recent) & 15 m, soil, laterite & 12 m, alluvium & 10 m, volcanic ash soil \\
\hline
\end{tabularx}
\end{center}
\textbf{Table 1:} Simplified lithostratigraphic logs (thickness in metres). All sections rest on a Precambrian basement (not exposed).

\begin{center}
\begin{tabularx}{\linewidth}{|l|X|X|X|}
\hline
\rowcolor{popblueL}
\textbf{Fossil group} & \textbf{Section A} & \textbf{Section B} & \textbf{Section C} \\
\hline
Benthic foraminifera & \emph{Orbitolina} sp. (Unit 4) & \emph{Rotalipora} sp. (Unit 4) & \emph{Globotruncana} sp. (Unit 4) \\
Planktonic foraminifera & – & \emph{Rotalipora cushmani} (Unit 4) & \emph{Globotruncana arca} (Unit 4) \\
Larger foraminifera & \emph{Nummulites} sp. (Unit 5) & \emph{Nummulites} sp. (Unit 5) & \emph{Miogypsina} sp. (Unit 5) \\
Palynomorphs & \emph{Classopollis} (Unit 3) & \emph{Afropollis} (Unit 3) & \emph{Retimonocolpites} (Unit 3) \\
Macrofauna & Oyster fragments (Unit 4) & Ammonite aptychi (Unit 4) & Echinoid spines (Unit 4) \\
\hline
\end{tabularx}
\end{center}
\textbf{Table 2:} Key fossil assemblages per section (units refer to Table~1).

\courssub{Tasks}
\begin{enumerate}
    \item \textbf{Lithostratigraphic columns.} Using the data in Table~1, draw three scaled lithostratigraphic columns (1 cm = 10 m) showing lithology, thickness, sedimentary structures, and the position of volcanic interbeds. Label each unit with a code (e.g., A1–A8).
    \item \textbf{Contacts and unconformities.} On each column, identify and name every contact (conformable, disconformity, angular unconformity, nonconformity). Justify each interpretation using the principles of stratigraphy and the fossil/absolute age data.
    \item \textbf{Correlation panel.} Construct a correlation panel aligning the three sections. Use:
        \begin{itemize}
            \item Lithological similarity (facies matching).
            \item Biostratigraphic markers (index fossils from Table~2).
            \item Absolute ages (radiometric dates) to constrain the age of key horizons.
        \end{itemize}
        Indicate the chronostratigraphic stages (e.g., Albian, Cenomanian, Paleocene, Miocene) on the panel.
    \item \textbf{Sequential geologic history.} Write a concise (≈300 words) narrative that sequences all events from basement formation to the present, incorporating: depositional environments, tectonic phases (rifting, thermal subsidence, CVL volcanism), sea‑level changes, and the identified unconformities. Reference the seismic reflectors (R1–R5) where appropriate.
    \item \textbf{Hydrocarbon trap proposal.} Based on the correlation panel and seismic sketch, propose a potential hydrocarbon trap linked to an identified unconformity. Specify:
        \begin{itemize}
            \item Trap type (structural, stratigraphic, or combination).
            \item Reservoir unit(s) and seal unit(s).
            \item Source rock candidate and maturation timing relative to trap formation.
            \item Risk factors (e.g., fault seal, volcanic intrusion).
        \end{itemize}
    \item \textbf{Poster preparation.} Design an A1 poster (portrait) that presents: (i) the three columns, (ii) the correlation panel, (iii) the geologic history as a timeline, (iv) the trap model with a cross‑section sketch, and (v) a conclusion box. The poster will be presented in a 5‑minute oral defence.
\end{enumerate}
\end{experience}

\begin{exoresolu}[Task 1: Lithostratigraphic Columns]
\textbf{Statement:} Draw three scaled lithostratigraphic columns (1 cm = 10 m) for Sections A, B, and C using Table 1 data. Show lithology, thickness, sedimentary structures, volcanic interbeds, and label units A1–A8, B1–B8, C1–C8.

\tcblower
\textbf{Detailed solution:}
\begin{popfigure}
\centering
\begin{tikzpicture}[scale=0.8]
% Column for Section A (Kribi) - 120m total = 12cm
% x=0 to 2cm width, y=0 to 12cm height (base at y=0)
% Lithology patterns (renamed to avoid conflicts with predefined patterns)
\pgfdeclarepatternformonly{mynwlines}{\pgfpointorigin}{\pgfpoint{4pt}{4pt}}{\pgfpoint{4pt}{4pt}}{
  \pgfpathmoveto{\pgfpoint{0pt}{4pt}}
  \pgfpathlineto{\pgfpoint{4pt}{0pt}}
  \pgfusepath{stroke}
}
\pgfdeclarepatternformonly{mycrosshatch}{\pgfpointorigin}{\pgfpoint{4pt}{4pt}}{\pgfpoint{4pt}{4pt}}{
  \pgfpathmoveto{\pgfpoint{0pt}{0pt}}
  \pgfpathlineto{\pgfpoint{4pt}{4pt}}
  \pgfpathmoveto{\pgfpoint{0pt}{4pt}}
  \pgfpathlineto{\pgfpoint{4pt}{0pt}}
  \pgfusepath{stroke}
}
\pgfdeclarepatternformonly{mybrick}{\pgfpointorigin}{\pgfpoint{6pt}{6pt}}{\pgfpoint{6pt}{6pt}}{
  \pgfpathrectangle{\pgfpoint{0pt}{0pt}}{\pgfpoint{6pt}{3pt}}
  \pgfpathrectangle{\pgfpoint{3pt}{3pt}}{\pgfpoint{6pt}{3pt}}
  \pgfusepath{stroke}
}
\pgfdeclarepatternformonly{mydots}{\pgfpointorigin}{\pgfpoint{4pt}{4pt}}{\pgfpoint{4pt}{4pt}}{
  \pgfpathcircle{\pgfpoint{2pt}{2pt}}{1pt}
  \pgfusepath{fill}
}

% Basement
\fill[pattern=mynwlines, pattern color=popdark] (0,0) rectangle (2,0.5);
\node[align=center, font=\tiny] at (1,0.25) {Basement};

% Unit A1: Basal conglomerate 5m = 0.5cm
\fill[pattern=mybrick, pattern color=poporange] (0,0.5) rectangle (2,1.0);
\node[align=center, font=\tiny] at (1,0.75) {A1: Conglomerate\\5m, clast-supp., quartzite};

% Unit A2: Sandstone 20m = 2cm
\fill[pattern=mycrosshatch, pattern color=popgold] (0,1.0) rectangle (2,3.0);
\node[align=center, font=\tiny] at (1,2.0) {A2: Sandstone (fluvial)\\20m, trough cross-strat.};

% Unit A3: Shale 15m = 1.5cm
\fill[pattern=mydots, pattern color=popmuted] (0,3.0) rectangle (2,4.5);
\node[align=center, font=\tiny] at (1,3.75) {A3: Shale (lacustrine)\\15m, laminated, organic-rich};

% Unit A4: Limestone 30m = 3cm
\fill[pattern=mynwlines, pattern color=popblue] (0,4.5) rectangle (2,7.5);
\node[align=center, font=\tiny] at (1,6.0) {A4: Limestone\\30m, bioclastic, Orbitolina};

% Tuff at 40m from base = 4cm (at top of A3 / base A4)
\draw[popred, thick] (0,4.5) -- (2,4.5);
\node[popred, font=\tiny, right] at (2.1,4.5) {Tuff: Ar-Ar 55.2 Ma};

% Unit A5: Marl 10m = 1cm
\fill[pattern=mycrosshatch, pattern color=popgreen] (0,7.5) rectangle (2,8.5);
\node[align=center, font=\tiny] at (1,8.0) {A5: Marl\\10m, Nummulites};

% Unit A6: Sandstone 25m = 2.5cm
\fill[pattern=mybrick, pattern color=popgold] (0,8.5) rectangle (2,11.0);
\node[align=center, font=\tiny] at (1,9.75) {A6: Sandstone (deltaic)\\25m, hummocky cross-strat.};

% Unit A7: Volcanic interbeds (none in middle, tuff already shown)
% Unit A8: Cover 15m = 1.5cm
\fill[gray!30] (0,11.0) rectangle (2,12.5);
\node[align=center, font=\tiny] at (1,11.75) {A8: Cover\\15m, soil, laterite};

% Scale bar
\draw[|-|, thick] (2.2,0) -- (2.2,10) node[midway, right, font=\small] {100 m};
\draw[|-|, thick] (2.2,10) -- (2.2,12) node[midway, right, font=\small] {20 m};

% Labels
\node[font=\bfseries\small, align=center] at (1,-0.8) {Section A (Kribi)\\120 m};
\end{tikzpicture}
\legende{Example lithostratigraphic column for Section A (Kribi) at 1 cm = 10 m. Students should draw similar columns for Sections B and C. Lithology patterns: conglomerate (brick), sandstone (crosshatch), shale (dots), limestone (NW lines), marl (crosshatch green), cover (grey). Volcanic interbeds shown as red lines with age labels.}
\end{popfigure}

\textbf{Key points for all columns:}
\begin{itemize}
    \item Use standard lithological patterns (see figure).
    \item Mark the exact position of volcanic interbeds: Section A tuff at 40 m (base of Unit 4), Section B tuff at 85 m (within Unit 4), Section C basalt at top 10 m (top of Unit 6 / base of Unit 8).
    \item Annotate sedimentary structures: trough cross-stratification (A2, B2), planar cross-beds (B2), ripple lamination (C2), lamination (A3, B3, C3), hummocky cross-stratification (A6), sigmoidal clinoforms (B6), distributary channels (C6).
    \item Label the basement contact at the base as a nonconformity (jagged line).
\end{itemize}
\end{exoresolu}

\begin{exoresolu}[Task 2: Contacts and Unconformities]
\textbf{Statement:} Identify and name every contact in each column. Justify using stratigraphic principles and fossil/absolute age data.

\tcblower
\textbf{Detailed solution:}
\begin{enumerate}
    \item \textbf{Base of each section (Units 1 on basement): Nonconformity.} 
    \emph{Justification:} Sedimentary rocks (conglomerates) overlie Precambrian crystalline basement. The contact represents a profound erosion surface with a significant time gap (Precambrian to Cretaceous). Principle of \cle{nonconformity}: sedimentary rocks on igneous/metamorphic basement.
    
    \item \textbf{Unit 1/2 contact (Conglomerate/Sandstone): Conformable (gradational).} 
    \emph{Justification:} Upward fining from conglomerate to sandstone indicates continuous deposition in an alluvial-fan to fluvial transition. No fossil break or erosional surface.
    
    \item \textbf{Unit 2/3 contact (Sandstone/Shale): Conformable.} 
    \emph{Justification:} Gradational change from fluvial channel sands to lacustrine laminated shales. Palynomorphs (\emph{Classopollis}, \emph{Afropollis}, \emph{Retimonocolpites}) in Unit 3 indicate continuous non-marine deposition.
    
    \item \textbf{Unit 3/4 contact (Shale/Limestone): Conformable (transgressive surface).} 
    \emph{Justification:} Sharp but conformable transition from lacustrine shale to shallow marine limestone with marine fossils (\emph{Orbitolina} in A, \emph{Rotalipora} in B, \emph{Globotruncana} in C). This records a marine transgression (seismic reflector R2). No erosion; deepening upward.
    
    \item \textbf{Unit 4/5 contact (Limestone/Marl): Conformable.} 
    \emph{Justification:} Gradational change from carbonate to marl with continuous fossil record (appearance of \emph{Nummulites} in A and B). Represents continued marine deposition with increased terrigenous input.
    
    \item \textbf{Unit 5/6 contact (Marl/Deltaic Sandstone):}
    \begin{itemize}
        \item \textbf{Sections A \& B: Disconformity.} 
        \emph{Justification:} Subtle erosional surface with lag gravel (implied by change from marine marl to deltaic sandstone), and a fossil gap: \emph{Nummulites} (Paleocene–Eocene) in Unit 5, but no fossils in Unit 6. The seismic reflector R4 (Eocene–Oligocene unconformity) correlates here. Missing Oligocene.
        \item \textbf{Section C: Paraconformity (a type of disconformity without obvious erosion).} 
        \emph{Justification:} Unit 5 contains \emph{Miogypsina} (Miocene) directly overlying Unit 4 with \emph{Globotruncana} (Maastrichtian). The contact is sharp but shows no erosional relief. A hiatus of ~30 Myr (Paleocene–Oligocene) is indicated by fossil mismatch. This is a \cle{paraconformity}.
    \end{itemize}
    
    \item \textbf{Volcanic interbeds:}
    \begin{itemize}
        \item \textbf{Section A tuff (40 m, Ar–Ar 55.2 Ma): Isochronous marker horizon.} It lies at the base of Unit 4 (limestone). The age (Early Eocene) conflicts with \emph{Orbitolina} (Albian) in the same unit. \textbf{Critical evaluation:} Either the tuff is a sill intruded later (cross-cutting), the Ar–Ar date reflects alteration, or the \emph{Orbitolina} is reworked. Given the principle of cross-cutting relationships, if the tuff is a true interbed, the limestone must be Eocene, implying \emph{Orbitolina} ranges higher here or is misidentified. Most likely: the tuff is a sill, and the 55.2 Ma date gives a minimum age for the host limestone.
        \item \textbf{Section B tuff (85 m, U–Pb 93.5 Ma): Isochronous marker horizon.} It lies within Unit 4 (limestone with \emph{Rotalipora}, Cenomanian–Turonian). The U–Pb age (93.5 Ma, Cenomanian) is consistent with the biostratigraphy. This is a reliable tie-point.
        \item \textbf{Section C basalt (top 10 m, K–Ar 28.3 Ma): Angular unconformity at its base.} The basalt flows over Unit 6 (deltaic sandstone) which is tilted 15° (field observation). The contact is an \cle{angular unconformity}: younger horizontal basalt on older tilted strata. The K–Ar age (28.3 Ma, Late Oligocene) dates the volcanism. \textbf{Note:} Unit 5 below contains \emph{Miogypsina} (Miocene, <23 Ma), creating an age inversion (Oligocene basalt over Miocene sediment). This suggests either: (a) the basalt is a sill, not a flow; (b) excess argon gives an erroneously old K–Ar age; (c) \emph{Miogypsina} is reworked. Field relations (flow topology vs. sill) are crucial.
    \end{itemize}
    
    \item \textbf{Top of sections (Unit 8): Nonconformity / modern erosion surface.} Soil/laterite/alluvium on bedrock.
\end{enumerate}

\begin{attention}[Common Pitfalls in Contact Identification]
\begin{itemize}
\item Do not confuse a \textbf{paraconformity} (gap without erosion) with a \textbf{conformable} contact. Always check fossil assemblages for time gaps.
\item A \textbf{nonconformity} specifically requires sedimentary rocks overlying igneous/metamorphic basement. An \textbf{angular unconformity} requires tilted/folded older rocks overlain by younger, horizontal strata.
\item \textbf{Radiometric dates on volcanic layers:} A date on a tuff \emph{within} a sedimentary unit gives a maximum age for overlying strata and minimum for underlying (if primary). A date on a basalt \emph{overlying} tilted strata dates the unconformity (post-tilting). Always verify field relationships (flow vs. sill) before accepting the date as depositional.
\end{itemize}
\end{attention}
\end{exoresolu}

\begin{exoresolu}[Task 3: Correlation Panel]
\textbf{Statement:} Construct a correlation panel aligning the three sections using lithology, biostratigraphy, and absolute ages. Indicate chronostratigraphic stages.

\tcblower
\textbf{Detailed solution:}
The correlation panel (Figure 2) aligns the sections by \cle{chronostratigraphic} age, not just lithological thickness. The key is the \cle{diachronous} nature of the marine transgression (Unit 4) and the major Eocene–Oligocene unconformity (R4).

\begin{popfigure}
\centering
\begin{tikzpicture}[scale=0.7]
% Correlation panel: three columns aligned by age (Ma)
% Age axis on left: 0 to 145 Ma (but we focus on 0-100)
% Columns at x=1 (A), x=5 (B), x=9 (C), width=3 each
% y = age in Ma, from 0 (top) to 100 (bottom). 1 cm = 10 Ma.

% Age scale
\draw[->, thick] (0,0) -- (0,10) node[above, font=\small] {Age (Ma)};
\foreach \age in {0,10,20,30,40,50,60,70,80,90,100} {
  \draw (0,\age/10) -- (-0.1,\age/10) node[left, font=\tiny] {\age};
}

% Section A column (x=1 to 4)
\draw[thick] (1,0) rectangle (4,12);

% Time lines (horizontal)
\foreach \age/\lbl in {0/Recent, 5/Pliocene, 23/Miocene, 28/Oligocene, 34/Eocene, 56/Paleocene, 66/K-Pg, 93/Cenomanian, 100/Albian, 145/Early Cret.} {
  \draw[popline, dashed] (0.5,\age/10) -- (11.5,\age/10) node[right, font=\tiny] {\lbl};
}

% Section A units with age ranges
% A1-A3: Early Cretaceous (Berriasian-Albian) ~145-100 Ma
\fill[poporange!40] (1,10) rectangle (4,14.5); % 145-100
\node[align=center, font=\tiny] at (2.5,12.25) {A1-A3: Rift\\Fluvial/Lacustrine};

% A4: Limestone with Orbitolina -> Albian ~113-100 Ma
\fill[popblue!40] (1,10) rectangle (4,11.3); % Albian 113-100
\node[align=center, font=\tiny] at (2.5,10.65) {A4: Limestone\\Orbitolina (Albian)};
\draw[poporange, thick] (1,10) -- (4,10) node[right, poporange, font=\tiny] {Tuff 55.2 Ma?};

% A5: Marl with Nummulites -> Paleocene-Eocene 66-34 Ma
\fill[popgreen!40] (1,3.4) rectangle (4,6.6);
\node[align=center, font=\tiny] at (2.5,5.0) {A5: Marl\\Nummulites (Paleo-Eoc)};

% A6: Deltaic sandstone -> Miocene 23-5 Ma (seismic R5)
\fill[popgold!50] (1,0.5) rectangle (4,3.4);
\node[align=center, font=\tiny] at (2.5,1.95) {A6: Deltaic Sand\\Miocene (R5)};

% A8: Cover
\fill[gray!30] (1,0) rectangle (4,0.5);
\node[font=\tiny] at (2.5,0.25) {Cover};

% Section B (x=5 to 8)
% B1-B3: Early Cretaceous
\fill[poporange!40] (5,10) rectangle (8,14.5);
\node[align=center, font=\tiny] at (6.5,12.25) {B1-B3: Rift\\Fluvial/Lacustrine};

% B4: Limestone with Rotalipora -> Cenomanian-Turonian ~100-90 Ma
\fill[popblue!40] (5,9) rectangle (8,11);
\node[align=center, font=\tiny] at (6.5,10) {B4: Limestone\\Rotalipora (Cenom-Tur)};
\draw[poporange, thick] (5,9.35) -- (8,9.35) node[right, poporange, font=\tiny] {Tuff 93.5 Ma (U-Pb)};

% B5: Marl with Nummulites -> Paleocene-Eocene
\fill[popgreen!40] (5,3.4) rectangle (8,6.6);
\node[align=center, font=\tiny] at (6.5,5.0) {B5: Marl\\Nummulites (Paleo-Eoc)};

% B6: Deltaic sandstone -> Miocene
\fill[popgold!50] (5,0.5) rectangle (8,3.4);
\node[align=center, font=\tiny] at (6.5,1.95) {B6: Deltaic Sand\\Miocene (R5 clinoforms)};

% B8: Cover
\fill[gray!30] (5,0) rectangle (8,0.5);

% Section C (x=9 to 12)
% C1-C3: Early Cretaceous
\fill[poporange!40] (9,10) rectangle (12,14.5);
\node[align=center, font=\tiny] at (10.5,12.25) {C1-C3: Rift\\Fluvial/Lacustrine};

% C4: Limestone with Globotruncana -> Campanian-Maastrichtian ~83-66 Ma
\fill[popblue!40] (9,6.6) rectangle (12,9);
\node[align=center, font=\tiny] at (10.5,7.8) {C4: Limestone\\Globotruncana (Camp-Maas)};

% C5: Marl with Miogypsina -> Miocene? But below basalt 28 Ma. Conflict.
\fill[popgreen!40] (9,2.3) rectangle (12,5.3); % Miocene 23-5
\node[align=center, font=\tiny] at (10.5,3.8) {C5: Marl\\Miogypsina (Miocene?)};
\draw[poporange, thick] (9,2.8) -- (12,2.8) node[right, poporange, font=\tiny] {Basalt 28.3 Ma (K-Ar)};

% C6: Deltaic sandstone -> Miocene? But below basalt.
\fill[popgold!50] (9,0.5) rectangle (12,2.3);
\node[align=center, font=\tiny] at (10.5,1.4) {C6: Deltaic Sand\\Pre-basalt?};

% C8: Cover
\fill[gray!30] (9,0) rectangle (12,0.5);

% Unconformity R4 (Eocene-Oligocene ~34 Ma)
\draw[poppurple, very thick, decorate, decoration={zigzag, segment length=4, amplitude=1}] (0.5,3.4) -- (11.5,3.4) node[above, poppurple, font=\small] {R4: Eocene-Oligocene Unconformity};

% Labels
\node[font=\bfseries\small] at (2.5,15) {Section A (Kribi)};
\node[font=\bfseries\small] at (6.5,15) {Section B (Douala)};
\node[font=\bfseries\small] at (10.5,15) {Section C (Limbe)};

% Correlation lines
\draw[popblue, thick, dashed] (4,10.65) -- (5,10) -- (8,10) -- (9,7.8); % Transgression
\draw[popgreen, thick, dashed] (4,5) -- (5,5) -- (8,5); % Nummulites
\draw[popgold, thick, dashed] (4,2) -- (5,2) -- (8,2) -- (9,1.4); % Miocene delta
\end{tikzpicture}
\legende{Correlation panel for the Douala Basin sections. Time increases upward (0 Ma at top). Lithological patterns as in Task 1. Red lines: volcanic horizons with radiometric ages. Purple zigzag: regional unconformity R4. Dashed lines: correlation of key horizons. Note the diachronous marine transgression (Unit 4: Albian in A, Cenomanian in B, Campanian in C) and the age conflict in Section C (Miocene fossils below Oligocene basalt).}
\end{popfigure}

\textbf{Correlation Table (Chronostratigraphic Alignment):}
\begin{center}
\begin{tabularx}{\linewidth}{|l|X|X|X|}
\hline
\rowcolor{popblueL}
\textbf{Chronostratigraphic Stage} & \textbf{Section A (Kribi)} & \textbf{Section B (Douala)} & \textbf{Section C (Limbe)} \\
\hline
Early Cretaceous (Berriasian–Albian) & Units A1–A3 (rift fill); A4 base (\emph{Orbitolina}) & Units B1–B3 (rift fill); B4 base (\emph{Rotalipora}) & Units C1–C3 (rift fill); C4 base (\emph{Globotruncana}) \\
Cenomanian (93.5 Ma) & \textbf{Hiatus / condensed?} & Unit B4 continues; \textbf{Tuff 93.5 Ma (U–Pb)} & – \\
Turonian–Santonian & – & Unit B4 continues & – \\
Campanian–Maastrichtian & – & – & Unit C4 (\emph{Globotruncana}) \\
Paleocene–Eocene & Unit A5 (\emph{Nummulites}); \textbf{Tuff 55.2 Ma (Ar–Ar) at base A4?} & Unit B5 (\emph{Nummulites}) & \textbf{Hiatus (Paraconformity)} \\
Eocene–Oligocene (~34 Ma) & \textbf{Unconformity R4 (Disconformity)} & \textbf{Unconformity R4 (Disconformity)} & \textbf{Angular Unconformity (tilting of Unit C6)} \\
Oligocene (28.3 Ma) & – & – & \textbf{Basalt flow 28.3 Ma (K–Ar)} on tilted Unit C6 \\
Miocene & Unit A6 (deltaic, hummocky) & Unit B6 (deltaic, sigmoidal clinoforms = R5) & Unit C5 (\emph{Miogypsina}), Unit C6 (deltaic) \\
Pliocene–Recent & Unit A8 (cover) & Unit B8 (cover) & Unit C8 (cover) \\
\hline
\end{tabularx}
\end{center}

\textbf{Key Correlation Conclusions:}
\begin{itemize}
    \item The marine transgression (Unit 4) is strongly \cle{diachronous}: Albian in the south (Kribi), Cenomanian in the centre (Douala), Campanian in the north (Limbe). This reflects the progressive northward flooding of the rift basin (seismic R2 onlap).
    \item The 93.5 Ma tuff in Section B provides a robust \cle{isochron} for the Cenomanian.
    \item The 55.2 Ma tuff in Section A is problematic: it occurs at the base of the Albian limestone. It likely represents a sill or a remobilized ash, giving a minimum age for the limestone. The \emph{Nummulites} in Unit A5 confirm Paleocene–Eocene deposition above.
    \item The regional unconformity R4 (Eocene–Oligocene) is a major sequence boundary. In A and B it is a disconformity at the Unit 5/6 boundary. In C it is an angular unconformity at the base of the basalt (Unit 7), with Unit 6 tilted beneath.
    \item The Miocene deltaic progradation (Unit 6) correlates basin-wide with seismic clinoforms R5. \emph{Miogypsina} in Section C Unit 5 confirms Miocene age for the marl, but its position below the 28.3 Ma basalt requires the basalt to be a sill or the K–Ar date to be erroneous (excess argon).
\end{itemize}
\end{exoresolu}

\begin{exoresolu}[Task 4: Sequential Geologic History]
\textbf{Statement:} Write a concise (≈300 words) narrative sequencing events from basement to present.

\tcblower
\textbf{Detailed solution:}
\begin{enumerate}
    \item \textbf{Precambrian (>540 Ma):} Formation and stabilization of the \cle{Congo Craton} basement. Pan-African orogeny (ca. 600–550 Ma) imprints a structural grain (NW–SE faults) later reactivated during rifting.
    \item \textbf{Early Cretaceous (Berriasian–Albian, ~145–100 Ma):} \cle{Rifting} of the South Atlantic initiates the Douala Basin. Alluvial fans (Unit 1) and braided fluvial systems (Unit 2) deposit on the basement \cle{nonconformity}. Lacustrine shales (Unit 3) accumulate in isolated rift lakes (organic-rich, potential source rocks). Seismic reflector R1 marks the base of this syn-rift sequence.
    \item \textbf{Albian–Cenomanian (~113–93 Ma):} Marine transgression floods the basin from the south (seismic R2 onlap). Shallow marine carbonates (Unit 4) with \emph{Orbitolina} (Section A) indicate early Albian seaway. In Section B, \emph{Rotalipora} and a \textbf{U–Pb zircon age of 93.5 ± 0.7 Ma} from an interbedded tuff anchor the Cenomanian transgression. Section C remains non-marine or receives deeper water facies later.
    \item \textbf{Turonian–Maastrichtian (~93–66 Ma):} Thermal subsidence dominates. Carbonate deposition continues (Unit 4 in Section C yields \emph{Globotruncana}, Campanian–Maastrichtian). The basin is a passive margin.
    \item \textbf{Paleocene–Eocene (~66–34 Ma):} Sea-level highstand deposits marls with \emph{Nummulites} (Unit 5) in Sections A and B. An ash fall in Section A yields \textbf{Ar–Ar 55.2 ± 0.4 Ma} (Early Eocene), but its position at the base of Unit 4 suggests it may be a sill; the \emph{Nummulites} in Unit 5 confirm Paleocene–Eocene age for the marl. Section C shows a \cle{paraconformity} (hiatus of ~30 Myr), likely due to uplift near the nascent Mount Cameroon.
    \item \textbf{Eocene–Oligocene (~34–28 Ma):} Major tectonic uplift and erosion linked to the onset of intense \cle{Cameroon Volcanic Line} activity. The \cle{regional unconformity R4} truncates the Paleogene sequence. In Section C, deltaic sands (Unit 6) are tilted 15° and truncated.
    \item \textbf{Late Oligocene (~28 Ma):} A basalt flow (\textbf{K–Ar 28.3 ± 1.2 Ma}) erupts over the eroded surface in Section C, forming an \cle{angular unconformity}. This marks the main CVL volcanic phase. (Note: The K–Ar age is older than the overlying Miocene \emph{Miogypsina} in Unit 5; the basalt is likely a sill or the date affected by excess argon.)
    \item \textbf{Miocene (~23–5 Ma):} Renewed subsidence and high sediment supply from the rising CVL. Deltaic sandstones (Unit 6) prograde basinward (seismic clinoforms R5). \emph{Miogypsina} in Section C Unit 5 confirms Miocene age for the transgressive marl.
    \item \textbf{Pliocene–Present:} Continued deltaic deposition, lateritic soil formation (Unit 8), and modern coastal processes.
\end{enumerate}
\end{exoresolu}

\begin{exoresolu}[Task 5: Hydrocarbon Trap Proposal]
\textbf{Statement:} Propose a potential hydrocarbon trap linked to an identified unconformity.

\tcblower
\textbf{Detailed solution:}
\begin{itemize}
    \item \textbf{Trap type:} Combined \cle{structural–stratigraphic trap}: a \textbf{truncation (onlap) trap} at the Eocene–Oligocene unconformity (R4). Miocene deltaic sandstones (Unit 6, reservoir) onlap the unconformity surface and are sealed by overlying Miocene shales/marls (top Unit 6 / Unit 8) and, locally, by the basalt flow (Section C).
    \item \textbf{Reservoir unit(s):} Unit 6 deltaic sandstones (hummocky cross-stratification in A, sigmoidal clinoforms in B, distributary channels in C). Good primary porosity/permeability expected. Laterally extensive on seismic (R5 clinoforms).
    \item \textbf{Seal unit(s):} 
    \begin{itemize}
        \item Vertical: Overlying Miocene marine shales/marls (top of Unit 6) and the regional basalt flow (where present, Section C) provide effective vertical seals.
        \item Lateral: Pinch-out of deltaic sandstones against the R4 unconformity surface (stratigraphic seal).
    \end{itemize}
    \item \textbf{Source rock candidate:} 
    \begin{itemize}
        \item \textbf{Primary:} Unit 3 lacustrine shales (organic-rich, Type I/II kerogen, TOC >2\% inferred from "organic-rich" description). Deposited in Early Cretaceous rift lakes.
        \item \textbf{Secondary:} Unit 5 marls (Type II kerogen, marine).
    \end{itemize}
    \item \textbf{Maturation timing relative to trap formation:} 
    \begin{itemize}
        \item Basin modelling suggests Unit 3 entered the oil window in Late Cretaceous (~80 Ma) and gas window in Miocene (~15 Ma).
        \item Trap formation (R4 unconformity + Miocene onlap) occurred at ~34 Ma (unconformity) and 23–5 Ma (reservoir deposition/seal). 
        \item \textbf{Critical timing:} Oil generation (Late Cretaceous) \emph{pre-dates} trap formation (Oligocene–Miocene). However, gas generation (Miocene) is \emph{synchronous} with trap formation. Migration pathways: carrier beds along the Sanaga Fault and along the R4 unconformity surface.
    \end{itemize}
    \item \textbf{Risk factors:}
    \begin{itemize}
        \item \textbf{Fault seal integrity:} The Sanaga Fault may act as a leak pathway if not sealed by shale gouge.
        \item \textbf{Volcanic intrusions:} Basalt flows/sills (Sections A, B, C) may have fractured seals or caused local thermal overmaturity (gas cracking).
        \item \textbf{Reservoir continuity:} Uncertainty in lateral extent of Unit 6 sandstones away from the three control points; seismic R5 clinoforms suggest good continuity but need 3D calibration.
        \item \textbf{Source rock maturity:} Risk that Unit 3 is overmature (gas only) in the basin centre; oil potential may be limited to flanks.
    \end{itemize}
    \item \textbf{Recommendation:} Acquire 3D seismic to map the onlap geometry of Unit 6 onto R4, image fault seals on the Sanaga Fault, and delineate basalt/sill distribution. Drill a test well at the crest of the onlap pinch-out (mid-basin, between Sections A and B) to test the truncation trap.
\end{itemize}
\end{exoresolu}

\begin{exoresolu}[Task 6: Poster Preparation \& Oral Defence]
\textbf{Statement:} Design an A1 poster (portrait) presenting the synthesis. Prepare a 5-minute oral defence.

\tcblower
\textbf{Detailed solution:}
\textbf{Poster Layout (A1 Portrait, 594 $\times$ 841 mm):}
\begin{enumerate}
    \item \textbf{Title Block (Top 10\%):} ``Stratigraphic Synthesis \& Hydrocarbon Potential of the Douala Basin, Cameroon'' --- Names, Date, GCE Advanced Level Geology.
    \item \textbf{Left Column (30\% width): Three Lithostratigraphic Columns.} Scaled 1 cm = 10 m, aligned at the basement nonconformity. Colour-coded lithology patterns (legend). Volcanic horizons highlighted in red with ages. Contacts labelled (NC, DIS, ANG, PARA).
    \item \textbf{Centre Column (40\% width): Correlation Panel.} As constructed in Task 3 (Figure 2). Chronostratigraphic scale on left. Correlation lines (dashed) for key horizons: Base Cretaceous (R1), Cenomanian Transgression (R2 + 93.5 Ma tuff), Paleogene Marl (Nummulites), R4 Unconformity, Miocene Delta (R5). Colour-code: Yellow = Reservoir (Unit 6), Grey = Seal (Unit 5, 8), Red = Volcanic, Blue = Source (Unit 3).
    \item \textbf{Right Column (30\% width): Geologic History Timeline.} Vertical timeline (0--145 Ma) with coloured bars for each tectono-stratigraphic phase: Rift (red), Transgression (blue), Thermal Subsidence (green), Uplift/Unconformity (purple), CVL Volcanism (orange), Deltaic Progradation (yellow). Key events annotated with ages and seismic reflectors (R1--R5).
    \item \textbf{Bottom Section (20\% height): Trap Model Cross-Section.} NW--SE sketch across the basin showing: Basement, Rift fill, Carbonate platform, R4 unconformity (zigzag), Miocene deltaic onlap (yellow), Seal (grey), Basalt sill/flow (red), Sanaga Fault. Hydrocarbon migration arrows from Unit 3 (source) up fault and along R4 into Unit 6 trap. Label: ``Truncation Trap at R4''.
    \item \textbf{Conclusion Box (Bottom right):} Bullet points: (1) Diachronous Albian--Campanian transgression. (2) R4 unconformity (34 Ma) creates regional trap. (3) Miocene deltaic sandstones (Unit 6) = reservoir. (4) Unit 3 lacustrine shales = source (gas prone). (5) Critical risks: fault seal, volcanic overmaturity. (6) Next step: 3D seismic + exploration well.
\end{enumerate}

\textbf{Oral Defence Key Messages (5 minutes):}
\begin{itemize}
    \item \textbf{Integration:} ``We integrated biostratigraphy (\emph{Orbitolina} $\rightarrow$ \emph{Rotalipora} $\rightarrow$ \emph{Globotruncana} $\rightarrow$ \emph{Nummulites} $\rightarrow$ \emph{Miogypsina}) with three radiometric dates (U--Pb, Ar--Ar, K--Ar) to build a robust chronostratigraphic framework.''
    \item \textbf{Critical Data Evaluation:} ``We identified two age conflicts: (a) the 55.2 Ma tuff at the base of Albian limestone in Section A --- interpreted as a sill giving a minimum age; (b) the 28.3 Ma basalt overlying Miocene \emph{Miogypsina} in Section C --- interpreted as a sill or excess argon. This critical evaluation is essential in real exploration.''
    \item \textbf{Trap Genesis:} ``The Eocene--Oligocene unconformity (R4), driven by CVL uplift, created the structural template. Miocene deltaic progradation (R5) provided the reservoir and seal in a classic onlap truncation trap.''
    \item \textbf{Cameroon Context:} ``This model applies the stratigraphic principles taught in this chapter to the Douala Basin, Cameroon's premier hydrocarbon province, demonstrating the link between CVL tectonics and petroleum systems.''
\end{itemize}

\begin{aretenir}[Examiner's Tips for GCE Advanced Level]
\begin{itemize}
\item In the practical exam, you may be given a similar dataset (logs, fossils, dates, seismic sketch).
\item \textbf{Always:} Draw columns to scale. Label every contact with its \emph{type} and \emph{justification} (principle + evidence).
\item \textbf{Correlation:} Align by \emph{time} (chronostratigraphy), not just lithology. Use fossils for relative age, radiometric dates for absolute calibration. Note diachroneity.
\item \textbf{History:} Use the sequence: Basement $\rightarrow$ Rift $\rightarrow$ Transgression $\rightarrow$ Subsidence $\rightarrow$ Uplift/Unconformity $\rightarrow$ Volcanism $\rightarrow$ Deltaic Progradation $\rightarrow$ Present. Link each to a seismic reflector if provided.
\item \textbf{Trap Proposal:} Be specific: Trap type, Reservoir, Seal, Source, Timing, Risks. Mention ``migration pathway''.
\item \textbf{Poster:} Clear visual hierarchy. Columns aligned. Correlation panel is the centrepiece. Cross-section must show trap geometry. Conclusion box = ``take-home message''.
\end{itemize}
\end{aretenir}
\end{exoresolu}

\begin{savaistu}[Cameroon Context: The Douala Basin \& CVL]
The Douala Basin (also called the Rio del Rey–Douala Basin) is the onshore/offshore extension of the Niger Delta petroleum system into Cameroon. It hosts the \cle{Logbaba} gas field (discovered 1980s, producing since 2013) and the \cle{Sanje} gas discovery. The basin's evolution is intimately tied to the \cle{Cameroon Volcanic Line} (CVL), a 1600 km Y-shaped chain of volcanoes from Pagalu (Annobón) to Lake Chad. The CVL's origin is debated: mantle plume vs. lithospheric fracture. Its activity since the Cretaceous has provided: (1) heat flow for source rock maturation, (2) volcaniclastics as reservoirs/seals, (3) structural traps (faults, folds, unconformities), and (4) risks (seal breach, overmaturity). The 1986 \cle{Lake Nyos} disaster (CO₂ release from a CVL maar) underscores the ongoing volcanic hazard. Understanding the stratigraphy of the Douala Basin is thus not only academic but of direct economic and societal relevance for Cameroon.
\end{savaistu}

\newpage

\coursec{Methods and worked examples}

\courssub{Methods and Worked Examples}

This section presents the essential methodological tools for the Stratigraphy chapter. Each method is broken down into clear, examinable steps, followed by a GCE-style worked example showing the expected layout, reasoning and precision. Master these techniques: they are the backbone of Paper 2 (Structured/Essay) and of the practical mapwork and data-handling questions.

\begin{methode}[Applying Stratigraphic Principles to Decipher a Cross-Section]
\textbf{Goal:} Determine the relative chronological order of geological events (deposition, folding, faulting, intrusion, erosion, unconformity) from a vertical cross-section.

\textbf{Step-by-step procedure:}
\begin{enumerate}
    \item \textbf{Identify all distinct rock units.} Label each sedimentary layer, igneous body (dyke, sill, pluton) and metamorphic basement with a letter (A, B, C\ldots). Note the lithology (shale, sandstone, limestone, granite, basalt).
    \item \textbf{Apply the \cle{Principle of Superposition}.} In an undeformed sequence, the oldest layer is at the bottom and the youngest at the top. \emph{Check for overturning:} look for younging indicators (graded bedding, cross-bedding, ripple marks, mudcracks). If the beds are overturned, the \emph{apparent} superposition is reversed.
    \item \textbf{Apply the \cle{Principle of Cross-Cutting Relationships}.} Any feature (fault, dyke, vein, unconformity surface) that cuts across a rock unit is \emph{younger} than that unit. A fault displacing a dyke is younger than the dyke.
    \item \textbf{Apply the \cle{Principle of Inclusions}.} Fragments (xenoliths) of country rock inside an igneous intrusion are older than the intrusion. A conglomerate containing granite pebbles implies the granite is older than the conglomerate.
    \item \textbf{Identify \cle{Unconformities}.} Look for a surface separating two sequences with different dips (\cle{angular unconformity}), a gap in the fossil record between parallel beds (\cle{disconformity}), or sedimentary rocks resting on eroded crystalline basement (\cle{nonconformity}). The rocks \emph{below} are older; the surface represents missing time (\cle{hiatus}); the rocks \emph{above} are younger.
    \item \textbf{Use \cle{Faunal Succession} (if fossils are given).} Index fossils constrain the relative age of specific layers (e.g., trilobites $\rightarrow$ Palaeozoic, typically Cambrian--Permian; ammonites $\rightarrow$ Mesozoic, especially Jurassic--Cretaceous).
    \item \textbf{Construct the \cle{Geological History}.} List events from \textbf{oldest (1)} to \textbf{youngest (last)} using standard phrasing: ``Deposition of Unit X'', ``Folding/tilting'', ``Erosion and planation'', ``Intrusion of Dyke Y'', ``Faulting (normal/reverse)'', ``Deposition of Unit Z''.
\end{enumerate}

\textbf{Examiner's reflexes:}
\begin{itemize}
    \item Always state the \textbf{principle used} for each deduction (e.g., ``By superposition, Shale A is older than Sandstone B'').
    \item Distinguish \textbf{relative age} (order of events) from \textbf{absolute age} (numerical age in years). This method gives relative age only.
    \item On maps: apply the \cle{rule of V's} for dipping beds crossing valleys, and remember that outcrop width increases as dip decreases (true thickness $=$ outcrop width $\times \sin(\text{dip})$ for horizontal ground).
\end{itemize}
\end{methode}

\begin{exoresolu}[Relative Dating from a Complex Cross-Section (GCE Style)]
\textbf{Statement:} A geological cross-section through a hypothetical area shows the following units:
\begin{itemize}
    \item \textbf{A:} Massive granite (contains xenoliths of D).
    \item \textbf{B:} Horizontal sandstone with \emph{Productus} (brachiopod) fossils (Carboniferous).
    \item \textbf{C:} Tilted shale/sandstone sequence (dipping $40^\circ$ SE), no fossils.
    \item \textbf{D:} Metamorphic gneiss (basement).
    \item \textbf{E:} Dolerite dyke (vertical, cuts A, C and D, but does \emph{not} cut B).
    \item \textbf{F:} Normal fault (displaces B, C, D and E; downthrow to the SE).
    \item \textbf{G:} Alluvial deposits (unconsolidated, flat-lying, covering all other units).
\end{itemize}
An angular unconformity separates C (below) from B (above). The dyke E is truncated by the unconformity surface beneath B.

\textbf{Task:} Establish the complete chronological sequence of geological events from oldest to youngest, justifying each step with the appropriate stratigraphic principle.

\tcblower
\textbf{Detailed Solution:}

\textbf{Step 1: Basement and intrusion.}
\begin{itemize}
    \item Unit \textbf{D (gneiss)} is the country rock; unit \textbf{A (granite)} contains \textbf{xenoliths of D}.
    \item \textbf{Principle of Inclusions:} the included rock (D) is older than the host (A).
    \item \textbf{Event 1:} Formation and metamorphism of the \textbf{gneiss D} (oldest basement).
    \item \textbf{Event 2:} Intrusion of the \textbf{granite A}.
\end{itemize}

\textbf{Step 2: Sedimentation, deformation and dyke emplacement.}
\begin{itemize}
    \item Unit \textbf{C} rests on the eroded basement (A, D) and is now tilted at $40^\circ$.
    \item \textbf{Principle of Original Horizontality:} C was deposited horizontally, so the tilting \emph{post-dates} the deposition of C.
    \item \textbf{Event 3:} Deposition of the \textbf{sediments C} on the eroded basement.
    \item \textbf{Event 4:} \textbf{Tectonic deformation (folding/tilting)} of C (and of A, D) to a $40^\circ$ dip.
    \item \textbf{Event 5:} \textbf{Erosion and planation} truncating the tilted edges of C, A and D, producing a sub-horizontal erosion surface (the future unconformity).
    \item The dyke \textbf{E} cuts A, C and D but is truncated by the unconformity and does \emph{not} cut B.
    \item \textbf{Cross-Cutting Relationships:} E is younger than A, C, D and younger than the tilting (it is itself vertical and undeformed); because it is truncated by the unconformity surface and sealed by B, it is \emph{older than B}.
    \item \textbf{Event 6:} Intrusion of the \textbf{dolerite dyke E}, after the tilting and before the deposition of B.
\end{itemize}

\textbf{Step 3: Unconformity and cover sequence.}
\begin{itemize}
    \item \textbf{B (horizontal sandstone)} lies on the eroded surface cutting C, A, D and E. B is horizontal while C dips at $40^\circ$: this is an \textbf{angular unconformity}.
    \item \textbf{Principle of Unconformities:} the rocks below (C, A, D, E) were deformed, intruded and eroded \emph{before} the deposition of B.
    \item Fossils in B: \emph{Productus} $\rightarrow$ Carboniferous (faunal succession).
    \item \textbf{Event 7:} Deposition of the \textbf{sandstone B} (Carboniferous).
\end{itemize}

\textbf{Step 4: Faulting and recent deposits.}
\begin{itemize}
    \item \textbf{F (normal fault)} displaces B, C, D and E. \textbf{Cross-cutting:} F is younger than all the units it offsets.
    \item \textbf{G (alluvium)} covers everything, including the fault scarp. \textbf{Superposition:} G is the youngest unit.
    \item \textbf{Event 8:} \textbf{Normal faulting F} (post-Carboniferous extension).
    \item \textbf{Event 9:} \textbf{Erosion} of the fault scarp.
    \item \textbf{Event 10:} Deposition of the \textbf{alluvial deposits G} (youngest).
\end{itemize}

\textbf{Final chronological table (oldest $\rightarrow$ youngest):}

\begin{center}
\begin{tabularx}{\linewidth}{|c|X|}\hline
\rowcolor{popblueL}\textbf{Order} & \textbf{Geological event} \\ \hline
1 & Formation of the metamorphic basement \textbf{gneiss (D)}. \\ \hline
2 & Intrusion of the \textbf{granite (A)} (xenoliths of D). \\ \hline
3 & Deposition of the \textbf{shale/sandstone sequence (C)} on the eroded basement. \\ \hline
4 & \textbf{Tectonic deformation (folding/tilting)} of C, A, D to a $40^\circ$ dip. \\ \hline
5 & \textbf{Erosion and planation} forming the unconformity surface. \\ \hline
6 & Intrusion of the \textbf{dolerite dyke (E)} cutting the basement and the tilted cover. \\ \hline
7 & Deposition of the \textbf{horizontal sandstone (B)} with \emph{Productus} (Carboniferous) on the unconformity. \\ \hline
8 & \textbf{Normal faulting (F)} displacing B, C, D, E (downthrow SE). \\ \hline
9 & Sub-aerial erosion of the fault scarp. \\ \hline
10 & Deposition of the \textbf{alluvial deposits (G)} (youngest). \\ \hline
\end{tabularx}
\end{center}

\textbf{Summary of justifications for the exam script:}
\begin{itemize}
    \item D older than A: inclusions (xenoliths).
    \item A, D older than C: superposition (C rests on the basement).
    \item C older than tilting: original horizontality.
    \item Tilting older than erosion: angular unconformity (truncation of C).
    \item E younger than tilting, older than B: cross-cutting plus truncation at the unconformity.
    \item B older than F: cross-cutting (F displaces B).
    \item F older than G: superposition (G covers the fault scarp).
\end{itemize}
\end{exoresolu}

\begin{methode}[Calculating Absolute Ages by Radiometric Dating]
\textbf{Goal:} Calculate the absolute age of a rock or mineral from parent/daughter isotope measurements, using either the \cle{single-sample} age equation or the \cle{isochron} (graphical) method.

\textbf{Fundamental relations:}
\begin{enumerate}
    \item \textbf{Decay constant and half-life:} $\lambda = \dfrac{\ln 2}{T_{1/2}} = \dfrac{0.693}{T_{1/2}}$.
    \item \textbf{Single-sample age equation} (closed system, initial daughter zero or corrected):
    \[ t = \frac{1}{\lambda} \ln \left( 1 + \frac{D^*}{P} \right) \]
    where $D^*$ is the accumulated \emph{radiogenic} daughter and $P$ the remaining parent (both in moles or atoms).
    \item \textbf{Isochron equation} (e.g., Rb--Sr):
    \[ \left(\frac{D}{D_{\mathrm{ref}}}\right)_{\text{present}} = \left(\frac{D}{D_{\mathrm{ref}}}\right)_{\text{initial}} + \left(\frac{P}{D_{\mathrm{ref}}}\right)_{\text{present}} \left(e^{\lambda t} - 1\right) \]
    Plot $y = D/D_{\mathrm{ref}}$ against $x = P/D_{\mathrm{ref}}$ for several cogenetic samples or minerals:
    \begin{itemize}
        \item \textbf{slope} $m = e^{\lambda t} - 1 \Rightarrow t = \dfrac{1}{\lambda}\ln(m+1)$;
        \item \textbf{intercept} $= (D/D_{\mathrm{ref}})_{\text{initial}}$, which constrains the magma source.
    \end{itemize}
\end{enumerate}

\textbf{Single-sample procedure (e.g., K--Ar):}
\begin{enumerate}
    \item \textbf{Identify the system:} parent, daughter, half-life or decay constants given (e.g., $^{40}$K $\rightarrow$ $^{40}$Ar, $T_{1/2} = 1.25 \times 10^{9}$ yr).
    \item \textbf{Work in moles (atoms), never in mass.} Convert mass to moles with $n = \dfrac{\text{mass}}{\text{molar mass}}$.
    \item \textbf{Correct for atmospheric daughter} (K--Ar): $^{40}Ar^{*} = {}^{40}Ar_{\text{meas}} - 295.5 \times {}^{36}Ar_{\text{meas}}$.
    \item \textbf{Compute $\lambda$:} for branching decay of $^{40}$K, $\lambda_{\text{total}} = \lambda_{\beta} + \lambda_{ec} = 5.543 \times 10^{-10}\ \mathrm{yr^{-1}}$, with $\lambda_{ec} = 0.581 \times 10^{-10}\ \mathrm{yr^{-1}}$ (the branch producing $^{40}$Ar).
    \item \textbf{Apply the branching age equation:}
    \[ t = \frac{1}{\lambda_{\text{total}}} \ln \left( 1 + \frac{\lambda_{\text{total}}}{\lambda_{ec}} \cdot \frac{{}^{40}Ar^{*}}{{}^{40}K} \right) \]
    \item \textbf{Round} to the precision of the data (2--3 significant figures) and express the result in Ma or Ga.
\end{enumerate}

\textbf{Isochron procedure (Rb--Sr):}
\begin{enumerate}
    \item Tabulate $^{87}$Rb/$^{86}$Sr (x-axis) and $^{87}$Sr/$^{86}$Sr (y-axis) for each sample.
    \item Draw the best-fit straight line; measure the slope $m = \Delta y / \Delta x$.
    \item Age: $t = \dfrac{1}{\lambda_{Rb}} \ln(m+1)$, with $\lambda_{Rb} = 1.42 \times 10^{-11}\ \mathrm{yr^{-1}}$.
    \item Interpret the intercept (initial $^{87}$Sr/$^{86}$Sr): a high value suggests an old crustal source; a low value ($\approx 0.702$--$0.704$) a mantle source.
    \item If the points scatter widely, the system was \textbf{open} (alteration, metamorphism): the ``age'' is meaningless.
\end{enumerate}

\begin{attention}[Common Traps in Radiometric Calculations]
\begin{itemize}
    \item \textbf{Confusing $\lambda$ and $T_{1/2}$:} $\lambda$ is in $\mathrm{yr^{-1}}$; $T_{1/2}$ is in years. Convert before calculating.
    \item \textbf{Mass versus moles:} isotope ratios are \emph{atomic} (mole) ratios. Convert concentrations (wt\%, ppm) to moles first.
    \item \textbf{Branching decay of $^{40}$K:} only about $11\%$ of decays produce $^{40}$Ar; the rest produce $^{40}$Ca. Use $\lambda_{\text{total}}$ in the prefactor and the ratio $\lambda_{\text{total}}/\lambda_{ec}$ inside the logarithm.
    \item \textbf{Closed-system assumption:} heating (metamorphism) or weathering can reset or disturb the clock. A K--Ar date records \emph{cooling/last closure}, not necessarily crystallisation.
    \item \textbf{Units of the answer:} with $\lambda$ in $\mathrm{yr^{-1}}$, $t$ comes out in \emph{years}; divide by $10^{6}$ for Ma.
\end{itemize}
\end{attention}
\end{methode}

\begin{exoresolu}[K--Ar Dating with Atmospheric Correction (Cameroon Volcanic Line)]
\textbf{Statement:} A biotite separate from a trachyte of \textbf{Mount Manengouba} (Cameroon Volcanic Line) was analysed for K--Ar dating. Results per gram of sample:
\begin{itemize}
    \item Potassium concentration: $8.00\%$ K by weight (molar mass of K $= 39.10\ \mathrm{g\,mol^{-1}}$).
    \item Isotopic abundance of $^{40}$K: $0.0117\%$ of total K.
    \item Measured $^{40}$Ar: $5.69 \times 10^{-10}\ \mathrm{mol\,g^{-1}}$.
    \item Measured $^{36}$Ar: $5.00 \times 10^{-13}\ \mathrm{mol\,g^{-1}}$.
    \item Atmospheric ratio $({}^{40}Ar/{}^{36}Ar)_{\text{atm}} = 295.5$.
    \item Decay constants: $\lambda_{\beta} = 4.962 \times 10^{-10}\ \mathrm{yr^{-1}}$; $\lambda_{ec} = 0.581 \times 10^{-10}\ \mathrm{yr^{-1}}$.
\end{itemize}
\textbf{Calculate the K--Ar age of the trachyte in Ma.}

\tcblower
\textbf{Detailed Solution:}

\textbf{1. Total decay constant:}
\[ \lambda_{\text{total}} = \lambda_{\beta} + \lambda_{ec} = (4.962 + 0.581) \times 10^{-10} = 5.543 \times 10^{-10}\ \mathrm{yr^{-1}} \]

\textbf{2. Moles of $^{40}$K per gram (parent $P$):}
\begin{align*}
\text{mass of K per g} &= 0.0800\ \mathrm{g} \\
n(K) &= \frac{0.0800}{39.10} = 2.046 \times 10^{-3}\ \mathrm{mol\,g^{-1}} \\
n({}^{40}K) &= 2.046 \times 10^{-3} \times 1.17 \times 10^{-4} = 2.394 \times 10^{-7}\ \mathrm{mol\,g^{-1}}
\end{align*}

\textbf{3. Atmospheric correction (radiogenic daughter $D^{*}$):}
\begin{align*}
{}^{40}Ar_{\text{atm}} &= 295.5 \times 5.00 \times 10^{-13} = 1.478 \times 10^{-10}\ \mathrm{mol\,g^{-1}} \\
{}^{40}Ar^{*} &= 5.69 \times 10^{-10} - 1.478 \times 10^{-10} = 4.212 \times 10^{-10}\ \mathrm{mol\,g^{-1}}
\end{align*}

\textbf{4. Branching age equation:}
\begin{align*}
\frac{{}^{40}Ar^{*}}{{}^{40}K} &= \frac{4.212 \times 10^{-10}}{2.394 \times 10^{-7}} = 1.759 \times 10^{-3} \\
\frac{\lambda_{\text{total}}}{\lambda_{ec}} &= \frac{5.543}{0.581} = 9.540 \\
1 + 9.540 \times 1.759 \times 10^{-3} &= 1.01679 \\
t &= \frac{\ln(1.01679)}{5.543 \times 10^{-10}} = \frac{0.01665}{5.543 \times 10^{-10}} = 3.00 \times 10^{7}\ \mathrm{yr}
\end{align*}

\textbf{Final answer:}
\[ \boxed{t \approx 30.0\ \mathrm{Ma} \quad \text{(Oligocene)}} \]

\textbf{Examiner's comments:}
\begin{itemize}
    \item The result ($\approx 30$ Ma) is \textbf{geologically plausible}: the continental sector of the Cameroon Volcanic Line has been active since the Paleogene, and Oligocene ages are known from the CVL. Always \emph{check the plausibility} of a calculated age against the geological context.
    \item The atmospheric correction is essential: without it, the age would be overestimated (excess apparent $^{40}$Ar).
    \item Unit check: $\lambda$ in $\mathrm{yr^{-1}}$ gives $t$ in years; $3.00 \times 10^{7}$ yr $= 30.0$ Ma.
    \item Significant figures: the data carry 3 significant figures, so report \textbf{30.0 Ma}.
\end{itemize}
\end{exoresolu}

\begin{methode}[Constructing and Interpreting a Stratigraphic Column (Log)]
\textbf{Goal:} Synthesise lithological, palaeontological and structural observations from a measured section or borehole into a standard graphical log, then interpret the depositional environments and history.

\textbf{Standard symbols (GCE convention):}
\begin{itemize}
    \item \textbf{Lithology:} sandstone (dots), shale (horizontal dashes), limestone (brick pattern), conglomerate (circles), coal (solid black), basalt (v-symbol), granite (crosses).
    \item \textbf{Sedimentary structures:} cross-bedding, ripple marks, graded bedding (coarse $\rightarrow$ fine), mudcracks, burrows.
    \item \textbf{Fossils:} sketch diagnostic forms (ammonite, trilobite, brachiopod, coral, plant debris).
    \item \textbf{Contacts:} sharp line (conformable), wavy/erosional line (unconformity), fault (tick on the downthrow side).
    \item \textbf{Scale:} always state the vertical scale (e.g., 1 cm $=$ 10 m).
\end{itemize}

\textbf{Step-by-step construction:}
\begin{enumerate}
    \item Draw the vertical axis with the depth scale and a central lithology column.
    \item Fill in the lithology patterns between the correct depths; show grain-size trends (fining-up or coarsening-up).
    \item Add sedimentary structures and fossils at their depths on the margins.
    \item Mark the key surfaces: erosion surfaces/sequence boundaries (SB), flooding surfaces, and the level of deepest water (\cle{maximum flooding surface}, MFS).
    \item Group beds into \cle{facies associations} and interpret the environment of each (use \cle{Walther's Law}: a vertical succession records laterally adjacent environments):
    \begin{itemize}
        \item Coarsening-up + marine fossils $\rightarrow$ prograding shoreface/delta;
        \item Fining-up + cross-bedding + plant debris $\rightarrow$ fluvial channel;
        \item Limestone + corals $\rightarrow$ reef/shallow platform;
        \item Black shale + pyrite + no bottom fauna $\rightarrow$ anoxic basin.
    \end{itemize}
    \item Assign relative ages from index fossils (biostratigraphy) or superposition.
    \item Write a synthesis sentence, e.g., ``The log records a transgressive--regressive cycle: basal fluvial conglomerate $\rightarrow$ estuarine heterolithics $\rightarrow$ offshore shale (MFS) $\rightarrow$ prograding deltaic sandstone.''
\end{enumerate}

\textbf{Exam skills:}
\begin{itemize}
    \item \textbf{Correlation:} match the MFS (shale/fossil peak) or erosion surfaces between logs.
    \item \textbf{Rates:} sedimentation rate $=$ thickness $/$ duration (m/Ma).
    \item \textbf{Gaps:} identify missing time (hiatus) at unconformities using missing fossil zones.
\end{itemize}
\end{methode}

\begin{exoresolu}[Interpreting a Borehole Log: Transgression--Regression (Rio del Rey Basin)]
\textbf{Statement:} Borehole BH-01 (Rio del Rey Basin, offshore Cameroon) recovered the following succession (top $\rightarrow$ bottom). Construct a schematic log (scale 1 cm $=$ 10 m) and interpret the depositional history.

\begin{center}
\begin{tabular}{|c|l|l|c|}\hline
\rowcolor{popblueL}\textbf{Depth (m)} & \textbf{Lithology} & \textbf{Structures / Fossils} & \textbf{Gamma ray (API)} \\ \hline
0--50 & Sand, shelly & \emph{Turritella}, \emph{Ostrea}, bioturbation & 40 (low) \\ \hline
50--120 & Grey fissile shale & Planktonic forams (\emph{Globigerina}), pyrite & 120 (high) \\ \hline
120--122 & \textbf{Hardground} (Fe/Mn crust) & \emph{Thalassinoides} burrows, phosphatic nodules & 80 \\ \hline
122--200 & Heterolithic sand/shale & Flaser bedding, \emph{Ophiomorpha}, brackish forams & 60--80 \\ \hline
200--250 & Massive sandstone & Trough cross-bedding, plant fragments, coal at top & 30 (low) \\ \hline
250--255 & Conglomerate & Imbrication, erosional base & 25 \\ \hline
255+ & Granite basement & Weathered top (saprolite) & --- \\ \hline
\end{tabular}
\end{center}

\textbf{Tasks:} (1) Draw the log. (2) Identify facies associations and environments. (3) Identify the sequence-stratigraphic surfaces. (4) Reconstruct the relative sea-level history.

\tcblower
\textbf{Detailed Solution:}

\textbf{1. Schematic log (description for the drawing):} vertical column 25.5 cm tall (1 cm $=$ 10 m).
\begin{itemize}
    \item 0--50 m: dot pattern (sand); shells and burrows; low gamma-ray block.
    \item 50--120 m: dashed pattern (shale); forams, pyrite; \textbf{gamma-ray peak} (120 API).
    \item 120--122 m: thick wavy line (hardground); burrows, phosphatic nodules.
    \item 122--200 m: alternating dots/dashes (heterolithic); flaser bedding, \emph{Ophiomorpha}.
    \item 200--250 m: dots (sandstone); trough cross-beds; coal streak at top.
    \item 250--255 m: circles (conglomerate); imbrication; sharp erosional base.
    \item 255 m+: crosses (granite); stippled weathered top.
\end{itemize}

\textbf{2. Facies associations and environments:}

\begin{center}
\begin{tabularx}{\linewidth}{|l|X|X|}\hline
\rowcolor{popblueL}\textbf{Interval} & \textbf{Facies association} & \textbf{Environment} \\ \hline
0--50 m & Clean sand, marine shells, bioturbated & \textbf{Shoreface / inner shelf} (high energy, oxygenated). \\ \hline
50--120 m & Fissile shale, planktonic forams, pyrite, max.\ gamma ray & \textbf{Offshore / prodelta} (low energy, deepest water). \\ \hline
120--122 m & Hardground, bored, phosphatic & \textbf{Condensed section / hiatus} (sediment starvation). \\ \hline
122--200 m & Heterolithic, flaser bedding, \emph{Ophiomorpha}, brackish forams & \textbf{Tidal flat / estuary} (tide-dominated, mixed energy). \\ \hline
200--250 m & Massive sand, cross-beds, plant/coal, fining-up & \textbf{Fluvial channel / coastal plain}. \\ \hline
250--255 m & Pebble conglomerate, imbrication, erosional base & \textbf{Braided river / alluvial fan}. \\ \hline
255 m+ & Weathered granite & \textbf{Basement / palaeosol}. \\ \hline
\end{tabularx}
\end{center}

\textbf{3. Sequence-stratigraphic interpretation:}
\begin{itemize}
    \item \textbf{Sequence boundary (SB):} at \textbf{255 m} (conglomerate on weathered granite): subaerial exposure and erosion.
    \item \textbf{Lowstand systems tract (LST):} 250--255 m (conglomerate) and 200--250 m (fluvial channel fill).
    \item \textbf{Transgressive surface (TS):} at \textbf{200 m} (base of the heterolithics; appearance of brackish forams).
    \item \textbf{Transgressive systems tract (TST):} 122--200 m (estuarine/tidal deposits; landward shift of facies, retrogradation).
    \item \textbf{Maximum flooding surface (MFS):} at \textbf{120--122 m} (hardground/condensed section, gamma-ray peak, planktonic forams, pyrite). \textbf{Key correlation marker.}
    \item \textbf{Highstand systems tract (HST):} 0--120 m (offshore shale passing up to shoreface sand): \textbf{progradation}, shallowing-upward.
\end{itemize}

\textbf{4. Relative sea-level history:}
\begin{itemize}
    \item 255--200 m: lowstand then rapid rise (valley filling).
    \item 200--120 m: rapid rise (transgression); the shoreline moves landward.
    \item 120 m: \textbf{turnaround} (maximum flooding).
    \item 120--0 m: stillstand then relative fall (progradation); the shoreline moves seaward.
\end{itemize}

\textbf{Summary statement for the exam:} ``The log records a complete depositional sequence: a sequence boundary on basement (255 m), lowstand fluvial fill, transgressive estuarine deposits, a maximum flooding surface marked by a phosphatic hardground and a gamma-ray peak at 120 m, and a highstand prograding shoreface succession. The hardground at the MFS represents a significant hiatus (condensation).''
\end{exoresolu}

\begin{methode}[Correlating Stratigraphic Sections (Biostratigraphy and Lithostratigraphy)]
\textbf{Goal:} Demonstrate the time-equivalence (or non-equivalence) of rock units in different sections (boreholes, outcrops) in order to reconstruct basin geometry.

\textbf{Correlation tools (in order of reliability for \emph{time} correlation):}
\begin{enumerate}
    \item \textbf{Biostratigraphy (chronostratigraphy):} \cle{index fossils} (ammonites, graptolites, foraminifera, nannofossils, palynomorphs) defining \cle{biozones}. Highest resolution; the gold standard.
    \item \textbf{Event stratigraphy:} volcanic ash beds (bentonites, datable by radiometry), magnetic reversals (magnetostratigraphy), sequence boundaries and maximum flooding surfaces.
    \item \textbf{Lithostratigraphy:} formations defined by lithology. \textbf{Warning:} lithological boundaries are commonly \cle{diachronous} (time-transgressive): a sandstone formation migrates basinward during regression. \textbf{Never correlate on lithology alone.}
\end{enumerate}

\textbf{Step-by-step procedure:}
\begin{enumerate}
    \item \textbf{Establish the local zonation:} list the fossil assemblages of each section; identify index fossils (short range, wide distribution, abundant) and define biozones.
    \item \textbf{Identify marker horizons:} ash beds, distinctive limestones, MFS (gamma-ray peak, condensed fauna), sequence boundaries (erosion, palaeosol).
    \item \textbf{Construct the correlation panel:} draw the sections side by side at the same vertical scale and \textbf{hang them on a time-line datum} (an ash bed or MFS), not on the present-day surface.
    \item \textbf{Draw the correlation lines:} connect \textbf{time lines} (biozone boundaries, ash beds, MFS) horizontally; draw \textbf{lithostratigraphic boundaries} as dipping lines to show diachronism.
    \item \textbf{Analyse thickness trends:} thickening $\rightarrow$ depocentre (maximum subsidence); thinning/onlap $\rightarrow$ palaeohigh.
    \item \textbf{Identify hiatuses:} missing biozones at unconformities or condensed sections; estimate the duration of the gap from the biozone durations.
\end{enumerate}

\textbf{Cameroon context:}
\begin{itemize}
    \item In the \textbf{Douala, Rio del Rey and Logone Birni basins}, correlation relies heavily on \textbf{palynomorphs} (pollen, spores, dinoflagellates) and \textbf{foraminifera}, because macrofossils are rare in deltaic and shaly successions.
    \item \textbf{Bentonite/tuff layers} linked to Cameroon Volcanic Line activity can provide radiometric tie-points within sedimentary basins.
    \item In the \textbf{Benue Trough} (e.g., Mamfe Basin), the Cretaceous marine transgressions are correlated using ammonites and foraminifera.
\end{itemize}
\end{methode}

\begin{exoresolu}[Multi-Well Correlation Panel (Logone Birni Basin)]
\textbf{Statement:} Three wells (A, B, C) penetrate the Cretaceous of the Logone Birni Basin (North Cameroon). Palynology gives the following depths (m) of key bio-events (FO $=$ first occurrence, LO $=$ last occurrence).

\begin{center}
\begin{tabular}{|l|c|c|c|}\hline
\rowcolor{popblueL}\textbf{Event} & \textbf{Well A (m)} & \textbf{Well B (m)} & \textbf{Well C (m)} \\ \hline
FO \emph{Classopollis} (base Albian) & 2450 & 2100 & 2600 \\ \hline
LO \emph{Afropollis jardinus} (top Albian) & 2100 & 1850 & 2200 \\ \hline
FO \emph{Cicatricosisporites} (base Cenomanian) & 2050 & 1800 & 2150 \\ \hline
\textbf{Bentonite bed (dated 98.5 Ma)} & 1950 & 1720 & \textbf{absent} \\ \hline
LO \emph{Spinizonocolpites} (top Cenomanian) & 1800 & 1600 & 1950 \\ \hline
FO \emph{Gabonisporis} (base Turonian) & 1750 & 1550 & 1900 \\ \hline
\end{tabular}
\end{center}

\textbf{Tasks:} (1) Construct a correlation panel hung on the bentonite datum (98.5 Ma), scale 1 cm $=$ 100 m. (2) Calculate the average sedimentation rate for the Albian in each well (Albian duration $\approx 12$ Ma, $\approx$113--101 Ma). (3) Interpret the basin geometry.

\tcblower
\textbf{Detailed Solution:}

\textbf{1. Panel construction:}
\begin{itemize}
    \item Draw a horizontal \textbf{datum line at 98.5 Ma} (bentonite) through wells A (1950 m) and B (1720 m).
    \item Plot the bio-events of each well at their depths relative to the datum (e.g., Well A: base Albian 500 m below datum; top Albian 150 m below; base Cenomanian 100 m below; top Cenomanian 150 m above; base Turonian 200 m above).
    \item \textbf{Well C has no bentonite:} interpolate the 98.5 Ma level within its Cenomanian. Cenomanian thickness $= 2150 - 1950 = 200$ m over $\approx 6.5$ Ma $\Rightarrow$ rate $\approx 30.8$ m/Ma. At 98.5 Ma (2.0 Ma after the base Cenomanian): depth $\approx 2150 - 2 \times 30.8 \approx 2088$ m.
    \item Connect the \textbf{time lines} (biozone boundaries, bentonite) horizontally between wells.
\end{itemize}

\textbf{2. Albian sedimentation rates} (duration 12 Ma):

\begin{center}
\begin{tabular}{|c|c|c|c|c|}\hline
\rowcolor{popblueL}\textbf{Well} & \textbf{Base Albian (m)} & \textbf{Top Albian (m)} & \textbf{Thickness (m)} & \textbf{Rate (m/Ma)} \\ \hline
A & 2450 & 2100 & 350 & $350/12 = 29.2$ \\ \hline
B & 2100 & 1850 & 250 & $250/12 = 20.8$ \\ \hline
C & 2600 & 2200 & 400 & $400/12 = 33.3$ \\ \hline
\end{tabular}
\end{center}

\textbf{3. Basin geometry:}
\begin{itemize}
    \item Albian thickness: C (400 m) $>$ A (350 m) $>$ B (250 m).
    \item \textbf{Well C} occupied the Albian \textbf{depocentre} (maximum subsidence and sediment supply).
    \item \textbf{Well B} stood on a relative \textbf{palaeohigh} (thinnest Albian).
    \item Cenomanian check: A $= 250$ m ($\approx 38.5$ m/Ma), B $= 200$ m, C $= 200$ m ($\approx 30.8$ m/Ma): subsidence at A accelerated in the Cenomanian, while B remained a persistent high.
    \item The bentonite (98.5 Ma) provides an \textbf{absolute time tie} between A and B; its absence in C (the depocentre) is most likely a preservation/recognition problem, so the 98.5 Ma level in C is interpolated.
\end{itemize}

\textbf{Exam answer structure:} (1) neat panel with horizontal datum and time lines; (2) rate table with calculations shown; (3) interpretation paragraph: ``Well C is the Albian depocentre (400 m; 33 m/Ma); Well B sits on a persistent palaeohigh; Well A shows accelerated Cenomanian subsidence. The dated bentonite confirms the biostratigraphic correlation.''
\end{exoresolu}

\begin{methode}[Solving the Sequence of Events from a Geological Map (Paper 3 Style)]
\textbf{Goal:} Derive the geological history from a map (lithologies, dips, faults, folds, unconformities).

\textbf{Step-by-step protocol:}
\begin{enumerate}
    \item \textbf{Read the legend:} units are usually listed oldest $\rightarrow$ youngest. Note the structural symbols: dip/strike, fold axes, faults (with downthrow side), unconformities.
    \item \textbf{Determine dip and younging:} younging is perpendicular to strike in the dip direction (upright beds). Apply the \cle{rule of V's}: a dipping bed makes a V pointing in the dip direction where it crosses a valley (for moderate dips); horizontal beds follow the contours.
    \item \textbf{Identify structures:} anticline (oldest rocks in the core, limbs dip outward), syncline (youngest in the core, limbs dip inward); faults (repetition or omission of strata across the trace).
    \item \textbf{Identify unconformities:} \cle{truncation} of older units beneath the surface; \cle{overstep/onlap} of the younger cover onto progressively older units.
    \item \textbf{Recognise inliers and outliers:} inlier $=$ older rocks surrounded by younger; outlier $=$ younger rocks surrounded by older.
    \item \textbf{Write the history (oldest $\rightarrow$ youngest):} deposition of the oldest unit $\rightarrow$ tectonic event(s) $\rightarrow$ erosion (unconformity) $\rightarrow$ deposition of the cover $\rightarrow$ intrusion/faulting $\rightarrow$ recent erosion and superficial deposits.
\end{enumerate}

\begin{attention}[Map Traps]
\begin{itemize}
    \item A \textbf{fault} repeats or omits strata; a simple \textbf{lithological boundary} follows strike and dip consistently. Do not confuse them.
    \item A fault \textbf{sealed} (covered but not cut) by a younger unit is older than that unit: this brackets the age of faulting.
    \item Valleys often follow fault zones (crushed rock erodes easily); ridges follow resistant beds (quartzite, dolerite).
\end{itemize}
\end{attention}
\end{methode}

\begin{exoresolu}[Map Interpretation: Unconformity, Tilting and Faulting]
\textbf{Statement:} A geological map shows:
\begin{itemize}
    \item \textbf{Unit 1 (oldest):} Precambrian gneiss (basement), outcropping in the north and south.
    \item \textbf{Unit 2:} Cambrian quartzite, vertical, N--S strike, forming a narrow central N--S belt.
    \item \textbf{Unit 3:} Ordovician shale, dipping $30^\circ$ E, east of the quartzite.
    \item \textbf{Unit 4:} Silurian limestone, dipping $30^\circ$ E, east of the shale.
    \item \textbf{Unit 5:} Carboniferous sandstone, horizontal, in the east, \textbf{overlapping} Units 2, 3, 4 and resting directly on Unit 1 in the north and south.
    \item A \textbf{normal fault} (downthrow east, N--S trace) cuts Units 1--4 but \textbf{terminates at the base of Unit 5}.
    \item Topography: a N--S valley follows the fault; a ridge follows the quartzite.
\end{itemize}

\textbf{Tasks:} (1) Sketch an E--W cross-section. (2) Write the complete geological history. (3) Explain why Unit 5 rests on Unit 1 in the north/south but on Units 2--4 in the centre.

\tcblower
\textbf{1. Cross-section (E--W), described:}

\begin{popfigure}
\begin{tikzpicture}[scale=1.0]
% basement
\fill[gray!40] (0,0) rectangle (2.2,2.0);
\node at (1.1,1.0) {\small 1 Gneiss};
% quartzite vertical
\fill[popgold!60] (2.2,0) rectangle (2.8,2.0);
\node[rotate=90] at (2.5,1.0) {\small 2};
% tilted shale and limestone
\fill[popblueL] (2.8,0) -- (4.6,0) -- (5.4,2.0) -- (3.6,2.0) -- cycle;
\fill[popgreenL] (3.6,2.0) -- (5.4,2.0) -- (6.2,2.0) -- (4.4,2.0) -- cycle;
\draw (2.8,0) -- (3.6,2.0);
\draw (4.6,0) -- (5.4,2.0);
\node at (4.2,0.9) {\small 3};
\node at (5.2,1.6) {\small 4};
% fault
\draw[very thick,popdark] (6.2,0) -- (6.2,2.0);
\node[below] at (6.2,0) {\small F};
% unconformity
\draw[very thick,poporange] (0,2.0) -- (9,2.0);
% cover
\fill[popgold!25] (0,2.0) rectangle (9,3.0);
\node at (7.6,2.5) {\small 5 Sandstone (horizontal)};
\node[above,poporange] at (3.0,2.0) {\small unconformity};
\draw[->] (0,3.4) -- (1.2,3.4) node[right]{\small E};
\end{tikzpicture}
\legende{Schematic E--W cross-section: tilted Units 2--4 and the fault F are truncated by the unconformity and sealed by the horizontal Carboniferous sandstone (Unit 5).}
\end{popfigure}

\textbf{2. Geological history (oldest $\rightarrow$ youngest):}
\begin{enumerate}
    \item Formation of the \textbf{Precambrian gneissic basement (Unit 1)}.
    \item Deposition of the \textbf{Cambrian quartzite (Unit 2)} on the eroded basement (nonconformity).
    \item Conformable deposition of the \textbf{Ordovician shale (Unit 3)} then the \textbf{Silurian limestone (Unit 4)}.
    \item \textbf{Tectonism:} tilting of Units 2--4 (to $30^\circ$ E; quartzite vertical) and \textbf{normal faulting} (downthrow east) affecting Units 1--4. This is \emph{post-Silurian, pre-Carboniferous}.
    \item \textbf{Prolonged erosion and planation:} truncation of the tilted units and of the fault; in the north and south the Palaeozoic cover was completely removed, re-exposing the gneiss. An \textbf{angular unconformity} formed.
    \item \textbf{Carboniferous transgression:} deposition of the horizontal \textbf{sandstone (Unit 5)}, onlapping the unconformity.
    \item \textbf{Recent erosion} sculpting the present landscape (valley on the crushed fault zone; ridge on the resistant quartzite).
\end{enumerate}

\textbf{3. Explanation of the map pattern (overstep):}
\begin{itemize}
    \item During the long pre-Carboniferous erosion interval, the tilted Palaeozoic sequence was \textbf{stripped off the northern and southern areas} (relative highs), exposing the basement, while it was \textbf{preserved in the central area} (protected by downfaulting to the east).
    \item The Carboniferous sandstone then blanketed the whole palaeo-landscape: it rests on gneiss where the Palaeozoic rocks had been removed (north/south) and on the truncated Palaeozoic units where they survived (centre). This regional relationship is called \cle{overstep}.
    \item \textbf{Significance:} the fault is \emph{sealed} by Unit 5, proving that faulting and tilting are \textbf{pre-Carboniferous}; the unconformity represents a major hiatus (the Devonian is missing).
\end{itemize}

\textbf{Examiner's tips:} use ``overstep'' for the regional pattern and ``onlap'' for the termination of beds against the unconformity; always bracket the age of a fault (``post-Silurian, pre-Carboniferous'').
\end{exoresolu}

\begin{methode}[Graphical Correlation (Shaw Diagram) and Hiatus Quantification]
\textbf{Goal:} Detect and quantify missing time (hiatus) or condensation in a section by plotting the depths of \emph{isochronous} events (fossil first/last occurrences, dated ash beds) in two wells against each other.

\textbf{Principle:}
\begin{itemize}
    \item Plot depth in Well A (x-axis) against depth in Well B (y-axis) for each synchronous event.
    \item The \cle{Line of Correlation (LOC)} through the points expresses the ratio of sedimentation rates between the wells.
    \item An event that plots \textbf{off} the LOC (too shallow in one well) indicates a \textbf{hiatus or condensed section} there.
    \item \textbf{Missing thickness} $=$ (depth predicted on the LOC) $-$ (actual depth); \textbf{hiatus duration} $=$ missing thickness $/$ local sedimentation rate.
\end{itemize}

\textbf{Step-by-step:}
\begin{enumerate}
    \item Select only reliable, isochronous tie-points (index-fossil events, radiometric dates) --- never lithological boundaries.
    \item Plot the points (Depth$_A$, Depth$_B$).
    \item Fit the LOC through the reliable points (exclude the suspect point).
    \item Predict where the suspect event \emph{should} lie from its depth in the other well.
    \item Compute the missing thickness, then divide by the sedimentation rate of the nearest undisturbed interval in the same well to obtain the hiatus duration.
\end{enumerate}
\end{methode}

\begin{exoresolu}[Quantifying a Hiatus with a Shaw Diagram (Douala Basin)]
\textbf{Statement:} Wells Alpha and Beta (Douala Basin) penetrate the Cretaceous. Planktonic foraminifera give the following zone-boundary depths (m):

\begin{center}
\begin{tabular}{|l|c|c|c|}\hline
\rowcolor{popblueL}\textbf{Bio-event} & \textbf{Age (Ma)} & \textbf{Well Alpha (m)} & \textbf{Well Beta (m)} \\ \hline
Top Albian (base Cenomanian) & 100.5 & 2800 & 2500 \\ \hline
Mid-Cenomanian event (MCE) & 96.0 & 2450 & 2250 \\ \hline
Top Cenomanian (base Turonian) & 93.9 & 2350 & \textbf{2150} \\ \hline
Mid-Turonian event & 91.5 & 2170 & 2050 \\ \hline
Top Turonian (base Coniacian) & 89.8 & 2072 & 1980 \\ \hline
\end{tabular}
\end{center}

Well Beta is suspected to contain a \textbf{hiatus/condensed section} near the Cenomanian--Turonian boundary.

\textbf{Tasks:} (1) Plot the Shaw diagram for the four reliable events and determine the LOC. (2) Predict the depth of the top Cenomanian in Beta if no hiatus existed. (3) Calculate the missing thickness and the hiatus duration, using the Turonian sedimentation rate of Beta.

\tcblower
\textbf{Detailed Solution:}

\textbf{1. Data points (Alpha $x$, Beta $y$) and LOC:}
\begin{itemize}
    \item Top Albian: (2800, 2500); MCE: (2450, 2250); Mid-Turonian: (2170, 2050); Top Turonian: (2072, 1980).
    \item These four points are collinear. Using the end members:
    \[ R = \frac{\Delta \text{Depth}_{\alpha}}{\Delta \text{Depth}_{\beta}} = \frac{2800 - 2072}{2500 - 1980} = \frac{728}{520} = 1.40 \]
    \item LOC equation (through (2800, 2500)):
    \[ \text{Depth}_{\alpha} = 1.40\,\text{Depth}_{\beta} - 700 \]
    \item Check (MCE): $1.40 \times 2250 - 700 = 2450$ \checkmark; (Mid-Turonian): $1.40 \times 2050 - 700 = 2170$ \checkmark.
\end{itemize}

\textbf{2. Predicted top-Cenomanian depth in Beta:} the top Cenomanian lies at 2350 m in Alpha:
\[ 2350 = 1.40\,D_{\beta} - 700 \;\Rightarrow\; D_{\beta} = \frac{3050}{1.40} = 2178.6\ \mathrm{m} \]

\textbf{3. Hiatus calculation:}
\begin{itemize}
    \item Actual depth in Beta: \textbf{2150 m}; predicted: \textbf{2178.6 m}.
    \item \textbf{Missing thickness:} $\Delta h = 2178.6 - 2150 = 28.6$ m.
    \item Turonian sedimentation rate in Beta (Mid-Turonian $\rightarrow$ Top Turonian):
    \[ SR_{\beta} = \frac{2050 - 1980}{91.5 - 89.8} = \frac{70\ \mathrm{m}}{1.7\ \mathrm{Ma}} = 41.2\ \mathrm{m/Ma} \]
    \item \textbf{Hiatus duration:}
    \[ \Delta t = \frac{\Delta h}{SR_{\beta}} = \frac{28.6}{41.2} \approx 0.69\ \mathrm{Ma} \]
\end{itemize}

\textbf{Conclusion:} Well Beta contains a hiatus of approximately \textbf{0.7 Ma} ($\approx 28.6$ m of missing section) at the Cenomanian--Turonian boundary. This coincides with the well-known \textbf{Cenomanian--Turonian boundary event (Oceanic Anoxic Event 2)}, commonly recorded as a condensed or erosive surface in the Douala Basin.

\textbf{Examiner's note:} always use the sedimentation rate of the \emph{same well} and of the \emph{nearest undisturbed interval}; using the other well's rate mixes two different subsidence histories.
\end{exoresolu}

\begin{methode}[Integrated Stratigraphic Interpretation: the ``Big Picture'' Essay]
\textbf{Goal:} Synthesise lithostratigraphy, biostratigraphy, chronostratigraphy, sequence stratigraphy and geochronology into a coherent basin-evolution model (typical GCE long essay / integration question).

\textbf{Structure of a high-scoring essay:}
\begin{enumerate}
    \item \textbf{Introduction:} locate the basin, state the data available (outcrops, wells) and the objective (stratigraphic framework and evolution).
    \item \textbf{Chronostratigraphic framework:} the relevant part of the geologic time scale; the biozones used; the key marker surfaces (MFS, sequence boundaries, dated ash beds).
    \item \textbf{Lithostratigraphy and facies:} the formations (lithology, thickness, boundaries); facies associations and depositional environments; lateral facies distribution.
    \item \textbf{Sequence stratigraphy / sea-level history:} depositional sequences, systems tracts (LST, TST, HST); eustatic versus tectonic control.
    \item \textbf{Tectono-stratigraphic synthesis:} the successive phases (rifting $\rightarrow$ breakup/drift $\rightarrow$ thermal subsidence $\rightarrow$ inversion/recent events), each with its characteristic stratigraphy.
    \item \textbf{Applications:} economic significance (hydrocarbons, groundwater, building materials) and geohazards where relevant.
    \item \textbf{Conclusion:} one summarising sentence.
\end{enumerate}

\textbf{Key terms that impress examiners:} \cle{accommodation space}, \cle{progradation / retrogradation / aggradation}, \cle{condensed section}, \cle{hiatus}, \cle{diachronous}, \cle{onlap / overstep}, \cle{systems tract}.
\end{methode}

\begin{exoresolu}[Integrated Basin-Evolution Essay Plan (Douala / Rio del Rey)]
\textbf{Statement:} ``Using your knowledge of the stratigraphy of the Douala/Rio del Rey Basin (Cameroon), describe the tectono-stratigraphic evolution from the Early Cretaceous to the Present. Illustrate your answer with the main sequences and key stratigraphic surfaces, and discuss the petroleum-system elements.''

\tcblower
\textbf{Model answer structure:}

\textbf{1. Introduction.} The Douala and Rio del Rey basins form the Cameroonian segment of the West African margin, initiated during the \textbf{Early Cretaceous} opening of the South Atlantic (breakup of Gondwana). Their evolution reflects rift tectonics, thermal subsidence, eustatic sea-level changes and, later, magmatism of the \textbf{Cameroon Volcanic Line (CVL)}.

\textbf{2. Tectono-stratigraphic phases:}

\begin{center}
\begin{tabularx}{\linewidth}{|l|l|X|}\hline
\rowcolor{popblueL}\textbf{Phase} & \textbf{Time} & \textbf{Processes and stratigraphy} \\ \hline
1. Syn-rift & Aptian--Albian & Extensional faulting; half-grabens filled with \textbf{continental clastics} (alluvial, fluvial, lacustrine; e.g., Mundeck-type facies); local volcanics; high sedimentation rates in depocentres. \\ \hline
2. Breakup / drift onset & Late Albian--Cenomanian & \textbf{Major marine transgression}; marine shales and limestones; maximum flooding; organic-rich \textbf{black shales} near the Cenomanian--Turonian boundary (Oceanic Anoxic Event 2) $\rightarrow$ regional \textbf{source rock}. \\ \hline
3. Thermal subsidence & Turonian--Maastrichtian & Long-wavelength subsidence; \textbf{deltaic and shoreface progradation} (reservoir sands); third-order sequences driven by eustasy. \\ \hline
4. Inversion / uplift & Paleogene & Compression and uplift (linked to Benue Trough inversion); \textbf{erosion and a major unconformity}; onset of CVL continental volcanism. \\ \hline
5. Neogene--Present & Miocene--Recent & Massive \textbf{clastic progradation} (deltaic systems related to the Niger Delta dispersal); growth faults; \textbf{active CVL volcanism} (Mount Cameroon eruptions; Lake Nyos and Lake Monoun maar craters). \\ \hline
\end{tabularx}
\end{center}

\textbf{3. Key stratigraphic surfaces to cite:} breakup unconformity (near the Albian--Cenomanian transition); maximum flooding surfaces (mid-Cenomanian; Cenomanian--Turonian boundary); the Paleogene erosion unconformity; the condensed black-shale interval of OAE 2.

\textbf{4. Petroleum-system elements (applied stratigraphy):}
\begin{itemize}
    \item \textbf{Source rocks:} Cenomanian--Turonian black shales (OAE 2, oil- and gas-prone); subordinate Albian lacustrine shales.
    \item \textbf{Reservoirs:} Albian fluvio-deltaic sands; Cenomanian shoreface sands; Late Cretaceous channel and turbidite sands; Miocene deep-water turbidites.
    \item \textbf{Seals:} the thick marine shale intervals (Cenomanian--Turonian, Maastrichtian, Miocene).
    \item \textbf{Traps:} fault-bounded blocks (syn-rift), rollover anticlines above growth faults (Neogene), stratigraphic pinchouts and unconformity traps.
    \item \textbf{Timing:} trap formation must coincide with or post-date hydrocarbon generation and migration; CVL heat flow locally accelerates maturation.
\end{itemize}

\textbf{5. Conclusion.} The Douala/Rio del Rey Basin records a classic margin evolution: continental syn-rift fill (Aptian--Albian), transgressive drift-onset marine deposition with the OAE 2 source rock (Cenomanian--Turonian), progradational passive-margin fill (Late Cretaceous), Paleogene inversion and erosion, and Neogene clastic progradation accompanied by CVL volcanism. The Cenomanian--Turonian boundary interval is the cornerstone of the regional petroleum system.
\end{exoresolu}

\begin{aretenir}[Key Points of this Methods Section]
\begin{itemize}
    \item Relative dating rests on five principles: \cle{superposition}, \cle{original horizontality}, \cle{cross-cutting relationships}, \cle{inclusions}, \cle{faunal succession}; unconformities mark \cle{hiatuses}.
    \item Always \textbf{name the principle} behind every deduction in a cross-section or map exercise.
    \item Radiometric ages: work in \textbf{moles}, correct for initial/atmospheric daughter, use $\lambda_{\text{total}}$ and the branching ratio for K--Ar, and check the \textbf{geological plausibility} of the result.
    \item Logs: read grain-size trends and fossils $\rightarrow$ facies $\rightarrow$ environments; the \cle{MFS} (gamma-ray/fossil peak, condensed section) is the best correlation marker.
    \item Correlation: time lines (biozones, ash beds) are horizontal; lithological boundaries are commonly \cle{diachronous} --- never correlate on lithology alone.
    \item Shaw diagram: LOC through reliable tie-points; hiatus $=$ missing thickness $/$ local sedimentation rate.
    \item Essays: structure = framework $\rightarrow$ facies $\rightarrow$ sequences $\rightarrow$ tectonic phases $\rightarrow$ applications; use Cameroonian examples (Douala/Rio del Rey basins, CVL, Benue Trough).
\end{itemize}
\end{aretenir}

\newpage

\coursec{Exercises}

\courssub{I check my knowledge}

\begin{exercice}[MCQ: Stratigraphic Principles]
Which of the following statements best describes the \cle{Principle of Cross-Cutting Relationships}?
\begin{enumerate}
\item Sedimentary layers are deposited horizontally.
\item A fault or intrusion is younger than the rocks it cuts across.
\item Fossils succeed each other in a definite order.
\item In an undeformed sequence, the oldest layer is at the bottom.
\end{enumerate}
\end{exercice}

\begin{exercice}[True/False with Justification]
State whether each statement is True or False. Justify your answer in one sentence.
\begin{enumerate}
\item The half-life of a radioactive isotope changes with temperature and pressure.
\item An angular unconformity always indicates a period of deformation followed by erosion before deposition of the overlying strata.
\item \emph{Index fossils} must have a wide geographic distribution but a narrow time range.
\item Carbon-14 dating can be used to date a basaltic lava flow from Mount Cameroon that erupted in 1999.
\end{enumerate}
\end{exercice}

\begin{exercice}[Key Definitions]
Define the following terms precisely:
\begin{enumerate}
\item \cle{Lithostratigraphic unit} vs \cle{Chronostratigraphic unit}.
\item \cle{Concordia diagram} (in U-Pb geochronology).
\item \cle{Magnetostratigraphy}.
\item \cle{Hiatus}.
\end{enumerate}
\end{exercice}

\begin{exercice}[Direct Calculation: Radiometric Dating]
A zircon crystal from a granitic intrusion in the \cle{Cameroon Volcanic Line} contains \SI{12.5}{\%} of its original \ce{^{238}U} (half-life = \SI{4.468e9}{years}). Assuming a closed system, calculate the age of the crystallisation of the zircon. Show your working.
\end{exercice}

\courssub{I practise}

\begin{exercice}[Relative Dating from a Cross-Section]
The schematic cross-section below shows a sequence of sedimentary rocks (S1--S4), a granitic intrusion (G), a dolerite dyke (D), and an overlying flat-lying sandstone (S5). A fault (F) offsets all units except S5.
\begin{popfigure}
\begin{tikzpicture}[scale=0.8]
% Layers S1-S4 (folded)
\draw[fill=popblueL!30] (0,0) -- (10,0) -- (10,1) -- (0,1) -- cycle; % S1
\draw[fill=popgreenL!30] (0,1) -- (10,1) -- (10,2) -- (0,2) -- cycle; % S2
\draw[fill=poporange!30] (0,2) -- (10,2) -- (10,3) -- (0,3) -- cycle; % S3
\draw[fill=poppink!30] (0,3) -- (10,3) -- (10,4) -- (0,4) -- cycle; % S4
% Fold deformation (simple sinusoidal)
\foreach \x in {0,0.5,...,10} {
    \pgfmathsetmacro{\y}{0.2*sin(\x*180)+2.5}
    \draw[popdark, thick] (\x,2) -- (\x,\y);
}
% Granite G (cross-cuts folds)
\draw[fill=poppurple!40, pattern=north east lines] (3,0.5) -- (7,0.5) -- (7,3.5) -- (3,3.5) -- cycle;
\node[poppurple, font=\bfseries] at (5,2) {G};
% Dolerite Dyke D (vertical, cuts G and sediments)
\draw[fill=popgreen!50!black] (8,0) -- (8.5,0) -- (8.5,4) -- (8,4) -- cycle;
\node[popgreen!50!black, font=\bfseries, rotate=90] at (8.25,2) {D};
% Fault F (offsets left side down)
\draw[popred, ultra thick] (2,0) -- (2,4);
\draw[popred, ->] (1.5,1.5) -- (1.5,2.5);
\draw[popred, ->] (2.5,1.5) -- (2.5,0.5);
\node[popred] at (2,4.3) {F};
% Sandstone S5 (flat, unconformable on top)
\draw[fill=yellow!60!poporange] (0,4) -- (10,4) -- (10,5) -- (0,5) -- cycle;
\node[popdark] at (5,4.5) {S5};
% Labels
\node[left] at (0,0.5) {S1};
\node[left] at (0,1.5) {S2};
\node[left] at (0,2.5) {S3};
\node[left] at (0,3.5) {S4};
\end{tikzpicture}
\legende{Schematic cross-section for relative dating exercise.}
\end{popfigure}
List the geological events (deposition of S1--S4, folding, intrusion G, dyke D, faulting F, erosion, deposition of S5) in correct chronological order from oldest to youngest. State the stratigraphic principle used for each relative age constraint.
\end{exercice}

\begin{exercice}[Biostratigraphic Correlation]
You are given the stratigraphic ranges of five fossil species (A--E) in the table below (Ma = millions of years ago). A well in the \cle{Rio del Rey Basin} penetrates three fossil assemblages at different depths:
\begin{itemize}
\item Assemblage 1: Species A, C
\item Assemblage 2: Species B, D
\item Assemblage 3: Species C, E
\end{itemize}
\begin{center}
\begin{tabular}{|c|c|c|}\hline
Species & First Appearance (Ma) & Last Appearance (Ma) \\ \hline
A & 145 & 130 \\
B & 135 & 110 \\
C & 140 & 120 \\
D & 125 & 90 \\
E & 115 & 85 \\ \hline
\end{tabular}
\end{center}
\begin{enumerate}
\item Determine the possible age range (in Ma) for each assemblage.
\item Identify which species could serve as an \cle{index fossil} for the Early Cretaceous (approx. 145--100 Ma) in this basin. Justify.
\item If a volcanic ash bed within Assemblage 2 yields a U-Pb zircon age of \SI{112 \pm 2}{Ma}, how does this refine the correlation?
\end{enumerate}
\end{exercice}

\begin{exercice}[Concordia-Discordia Interpretation]
Three zircon analyses from a granitic gneiss in the \cle{Congo Craton} (Southern Cameroon) plot on a Concordia diagram as follows (ratios are \ce{^{207}Pb/^{235}U} vs \ce{^{206}Pb/^{238}U}):
\begin{itemize}
\item Analysis 1: (0.125, 0.0185)
\item Analysis 2: (0.210, 0.0310)
\item Analysis 3: (0.350, 0.0515)
\end{itemize}
The points define a straight line (Discordia) intersecting the Concordia curve at two points.
\begin{enumerate}
\item Sketch a labelled Concordia diagram showing the three analyses and the Discordia line. Indicate the upper and lower intercepts.
\item If the upper intercept corresponds to \SI{2050 \pm 20}{Ma} and the lower intercept to \SI{550 \pm 30}{Ma}, interpret the geological history of the sample.
\item Why do zircon grains often yield discordant ages? Mention two mechanisms.
\end{enumerate}
\end{exercice}

\begin{exercice}[Unconformities: Identification and Field Evidence]
Draw a single schematic stratigraphic column (approx. 15 cm tall) showing, from bottom to top:
\begin{enumerate}
\item A \cle{nonconformity} between Precambrian gneiss and Cambrian sandstone.
\item An \cle{angular unconformity} between folded Silurian shales and horizontal Devonian sandstone.
\item A \cle{disconformity} within the Carboniferous limestone (with a palaeosol).
\item A \cle{paraconformity} between two Jurassic clay units (identified only by fossil gap).
\end{enumerate}
Label each contact clearly. For each unconformity type, list \emph{two} distinct field criteria you would use to recognise it.
\end{exercice}

\courssub{I prepare for the GCE}

\begin{exercice}[Integrated Relative and Absolute Dating -- 20 marks]
The figure below represents a geological cross-section through part of the \cle{Benue Trough} (Northern Cameroon).
\begin{popfigure}
\begin{tikzpicture}[scale=0.7]
% Basement
\draw[fill=poppurple!30, pattern=north west lines] (0,0) rectangle (14,2);
\node at (7,1) {\textbf{Basement Gneiss (Bg)}};
% Unconformity (wavy)
\draw[thick, decorate, decoration={snake, amplitude=0.1cm, segment length=0.5cm}] (0,2) -- (14,2);
% Sedimentary sequence (tilted)
\draw[fill=popblueL!40] (0,2) -- (14,2) -- (14,5) -- (0,5) -- cycle;
\node at (7,3.5) {\textbf{Tilted Cretaceous Sandstones (Ks)}};
% Granitic Intrusion
\draw[fill=poppink!40, pattern=crosshatch dots] (4,1) -- (10,1) -- (10,4) -- (4,4) -- cycle;
\node[popdark, font=\bfseries] at (7,2.5) {\textbf{Granite (Gr)}};
% Dyke
\draw[fill=popgreen!50!black] (11,0) -- (11.5,0) -- (11.5,5) -- (11,5) -- cycle;
\node[popgreen!50!black, rotate=90, font=\bfseries] at (11.25,2.5) {\textbf{Dolerite Dyke (Dd)}};
% Fault
\draw[popred, ultra thick] (2,0) -- (2,5);
\draw[popred, ->] (1.5,2) -- (1.5,3);
\draw[popred, ->] (2.5,2) -- (2.5,1);
\node[popred] at (2,5.3) {\textbf{Fault (F)}};
% Flat-lying cover
\draw[fill=yellow!60!poporange] (0,5) -- (14,5) -- (14,7) -- (0,7) -- cycle;
\node at (7,6) {\textbf{Horizontal Tertiary Sands (Ts)}};
% U-Pb Age label
\draw[<-, popblue] (7,2.5) -- (7,8) node[above] {\textbf{U-Pb Zircon Age (Gr) = 88.5 \pm 0.5 Ma}};
\end{tikzpicture}
\legende{Cross-section for integrated dating exercise.}
\end{popfigure}
\begin{enumerate}
\item Establish the complete relative chronological sequence of \emph{all} events shown (deposition, tilting/folding, intrusion, dyke emplacement, faulting, erosion, deposition of Ts). [8 marks]
\item Using the U-Pb zircon age provided for the Granite (Gr), state the \emph{minimum absolute age} of the tilting/folding event affecting the Cretaceous Sandstones (Ks). Explain your reasoning. [4 marks]
\item The Dolerite Dyke (Dd) is not dated. However, a whole-rock K-Ar analysis on the dyke yields \SI{42 \pm 2}{Ma}. Discuss the reliability of this K-Ar age compared to the U-Pb zircon age for constraining the timing of dyke emplacement. Consider the closure temperatures of the systems. [4 marks]
\item The contact between the Basement Gneiss (Bg) and the Cretaceous Sandstones (Ks) is a \cle{nonconformity}. Describe the field evidence you would look for to confirm this interpretation. [4 marks]
\end{enumerate}
\end{exercice}

\begin{exercice}[Essay: Magnetostratigraphy and Chemostratigraphy in Cameroon -- 25 marks]
\textbf{``Discuss the role of magnetostratigraphy and chemostratigraphy in refining the correlation of Cretaceous sequences in the Atlantic margin of Cameroon.''}

In your answer, you should:
\begin{itemize}
\item Define both methods and the principle on which each relies.
\item Explain how the \cle{Geomagnetic Polarity Time Scale (GPTS)} is constructed and used for correlation.
\item Describe the main \cle{chemostratigraphic} proxies (e.g., $\delta^{13}$C, $\delta^{18}$O, $^{87}$Sr/$^{86}$Sr) applicable to Cretaceous marine sediments.
\item Evaluate the strengths and limitations of each method when applied to the \cle{Rio del Rey}, \cle{Douala}, and \cle{Kribi-Campo} basins.
\item Provide a specific example (real or hypothetical) of how the integration of magnetostratigraphy, chemostratigraphy, and biostratigraphy resolved a correlation problem in the Cretaceous of the Cameroon Atlantic margin.
\end{itemize}
\end{exercice}

\begin{exercice}[Multi-Section Correlation and Sequence Stratigraphy -- 20 marks]
Three boreholes (BH-1, BH-2, BH-3) are drilled along a 50 km transect perpendicular to the coast in the \cle{Douala Basin}. The simplified lithologs and key data are given below.

\begin{center}
\begin{tabular}{|c|c|c|c|}\hline
\textbf{Depth (m)} & \textbf{BH-1 (Onshore)} & \textbf{BH-2 (Coastal)} & \textbf{BH-3 (Offshore)} \\ \hline
0--200 & Yellow sands, lignite & Grey clays, forams & Grey clays, forams \\ \hline
200--450 & \textbf{Unconformity (U1)} & \textbf{Unconformity (U1)} & Grey marls, \emph{Orbitolina} \\ \hline
450--600 & Red beds, conglomerates & Yellow sands, lignite & Grey marls, \emph{Orbitolina} \\ \hline
600--800 & \textbf{Basement} & Red beds, conglomerates & Black shales, TOC 4\%, \emph{Inoceramus} \\ \hline
800--1000 & -- & \textbf{Basement} & Black shales, TOC 4\%, \emph{Inoceramus} \\ \hline
\end{tabular}
\end{center}

Additional data:
\begin{itemize}
\item Biostratigraphy: \emph{Orbitolina} (Aptian), \emph{Inoceramus} (Turonian-Coniacian). Planktonic forams in BH-2/3 top unit: \emph{Globotruncana} (Campanian-Maastrichtian).
\item Radiometric: Volcanic ash at 550 m in BH-3: U-Pb zircon = \SI{93.5 \pm 0.3}{Ma} (Cenomanian-Turonian boundary).
\item Seismic: U1 correlates with a major regional reflector (Base Tertiary Unconformity).
\end{itemize}

\begin{enumerate}
\item Construct a correlation panel (fence diagram) at a vertical scale of 1 cm = 100 m, showing the three boreholes, the correlation of lithostratigraphic units, and the unconformity U1. [6 marks]
\item Identify the \cle{systems tracts} (Lowstand, Transgressive, Highstand) represented in the offshore section (BH-3) between 450 m and 1000 m, assuming a eustatic sea-level curve with a major transgression in the Aptian-Albian and a highstand in the Late Cretaceous. Justify using the lithology and fossil data. [6 marks]
\item Calculate the average sedimentation rate (m/Ma) for the black shale unit in BH-3 (800--1000 m) using the radiometric age constraint and the biostratigraphic range of \emph{Inoceramus}. [4 marks]
\item Discuss the hydrocarbon potential (source rock, reservoir, seal, trap) of the correlated sequence, referencing the specific units in the three boreholes. [4 marks]
\end{enumerate}
\end{exercice}

\newpage

\coursec{Solutions to the exercises}

\begin{corrige}[Exercise 1 — MCQ: Stratigraphic Principles]
\textbf{Correct answer: (b)} — A fault or intrusion is younger than the rocks it cuts across.

\textbf{Justification:} The \cle{Principle of Cross-Cutting Relationships} (formulated by James Hutton and Charles Lyell) states that any geological feature (fault, dyke, intrusion, vein) that cuts across another rock body must be \emph{younger} than the rock it cuts, because the rock must already have existed in order to be cut.

\textbf{Why the other options are wrong:}
\begin{enumerate}
\item Statement (a) describes the \cle{Principle of Original Horizontality} (Steno).
\item Statement (c) describes the \cle{Principle of Faunal Succession} (William Smith).
\item Statement (d) describes the \cle{Principle of Superposition} (Steno).
\end{enumerate}

\begin{attention}[Examiner's tip]
In MCQs, read every option: three of them are \emph{true statements} but they describe \emph{other} principles. The question tests your ability to \emph{match each principle to its correct statement}, not merely to recognise a true sentence.
\end{attention}
\end{corrige}

\begin{corrige}[Exercise 2 — True/False with Justification]
\begin{enumerate}
\item \textbf{TRUE.} The half-life of a radioactive isotope is a nuclear property that is essentially constant and unaffected by temperature, pressure or chemical state; this is precisely why radiometric dating is reliable.

\item \textbf{TRUE.} An angular unconformity records a sequence of: deposition of the lower strata $\rightarrow$ tectonic deformation (tilting/folding) $\rightarrow$ uplift and erosion producing a truncated surface $\rightarrow$ renewed deposition of the overlying horizontal strata.

\item \textbf{TRUE.} A good index fossil combines a \emph{wide geographic distribution} (so it can correlate rocks over large areas) with a \emph{short stratigraphic range} (so it defines a narrow time interval), plus abundance and easy identification.

\item \textbf{FALSE.} Carbon-14 dating only works on \emph{organic} material younger than about \SI{50000}{years}; a basalt contains no organic carbon, and a 1999 eruption is far too recent — the method of choice for young lavas would be K-Ar or \ce{^{40}Ar}/\ce{^{39}Ar} dating.
\end{enumerate}
\end{corrige}

\begin{corrige}[Exercise 3 — Key Definitions]
\begin{enumerate}
\item \textbf{Lithostratigraphic unit vs Chronostratigraphic unit:} A \cle{lithostratigraphic unit} (e.g., formation, member, group) is a body of rock defined solely by its observable \emph{lithological characteristics} (rock type, texture, composition); its boundaries may be \emph{diachronous} (of different ages in different places). A \cle{chronostratigraphic unit} (e.g., system, series, stage) is a body of rock defined by the \emph{time interval} during which it was deposited; its boundaries are \emph{isochronous} (same age everywhere).

\item \textbf{Concordia diagram:} A graph plotting the ratio \ce{^{206}Pb}/\ce{^{238}U} (x-axis) against \ce{^{207}Pb}/\ce{^{235}U} (y-axis). The \emph{Concordia curve} is the locus of points for which both U-Pb decay systems give the \emph{same} age. A sample that has remained a closed system plots \emph{on} the curve (concordant age); samples that have lost lead plot below it along a straight line, the \emph{Discordia}, whose upper and lower intercepts date the crystallisation and the disturbance event respectively.

\item \textbf{Magnetostratigraphy:} The branch of stratigraphy that uses the record of \emph{reversals of the Earth's magnetic field} preserved in rocks (remanent magnetisation). Sequences of normal and reversed polarity zones are matched against the \cle{Geomagnetic Polarity Time Scale (GPTS)} to date and correlate strata.

\item \textbf{Hiatus:} The \emph{interval of geological time not represented by any rock} at an unconformity or depositional break — i.e., the time gap corresponding to non-deposition and/or erosion between two successive strata.
\end{enumerate}
\end{corrige}

\begin{corrige}[Exercise 4 — Direct Calculation: Radiometric Dating]
\textbf{Given:} $N/N_0 = 12.5\% = 0.125$; half-life $T_{1/2} = \SI{4.468e9}{years}$ (\ce{^{238}U} decay).

\textbf{Decay equation:}
\[
\frac{N}{N_0} = e^{-\lambda t}, \qquad \lambda = \frac{\ln 2}{T_{1/2}}.
\]

\textbf{Step 1 — recognise the number of half-lives:}
\[
0.125 = \frac{1}{8} = \left(\frac{1}{2}\right)^3
\]
so exactly \textbf{three half-lives} have elapsed.

\textbf{Step 2 — compute the age:}
\[
t = 3 \times T_{1/2} = 3 \times 4.468\times 10^{9} = 1.3404\times 10^{10}\ \text{years}.
\]

\textbf{Check with the logarithmic formula:}
\begin{align*}
t &= \frac{1}{\lambda}\ln\!\left(\frac{N_0}{N}\right) = \frac{T_{1/2}}{\ln 2}\ln\!\left(\frac{1}{0.125}\right)\\
&= \frac{4.468\times 10^{9}}{0.6931}\times \ln 8
= \frac{4.468\times 10^{9}}{0.6931}\times 2.0794\\
&\approx 1.34\times 10^{10}\ \text{years}.
\end{align*}

\[
\boxed{t \approx \SI{13.4}{Ga} \text{ (13.4 billion years)}}
\]

\begin{attention}[Critical discussion — do not skip this in the exam!]
This age ($\approx$ \SI{13.4}{Ga}) is \emph{older than the Solar System} ($\approx$ \SI{4.56}{Ga}) and nearly the age of the Universe — it is \textbf{geologically impossible} for a terrestrial zircon. The closed-system assumption must therefore have failed (e.g., the measured uranium content is wrong, uranium was lost, or the "12.5\%" figure is unrealistic for this isotope system). In an exam, \emph{always comment on the plausibility of a numerical result}: full marks require both the correct calculation \emph{and} the critical remark that such a zircon cannot exist, showing that you understand the limits of the method.
\end{attention}
\end{corrige}

\begin{corrige}[Exercise 5 — Relative Dating from a Cross-Section]
\textbf{Chronological sequence (oldest $\rightarrow$ youngest):}

\begin{center}
\begin{tabularx}{\linewidth}{|c|X|X|}\hline
\rowcolor{popblueL}
\textbf{Order} & \textbf{Event} & \textbf{Principle used} \\ \hline
1 & Deposition of S1, then S2, S3, S4 & Superposition (oldest at bottom) \\ \hline
2 & Folding of S1--S4 & Folds deform already-deposited layers; layers must pre-date deformation \\ \hline
3 & Intrusion of granite G & Cross-cutting relationships: G cuts the folded strata (and baked/contact-metamorphosed margin confirms intrusion after folding) \\ \hline
4 & Emplacement of dolerite dyke D & Cross-cutting relationships: D cuts across S1--S4 \emph{and} G, so D is younger than both \\ \hline
5 & Faulting F & Cross-cutting relationships: F offsets S1--S4, G and D \\ \hline
6 & Erosion (planation surface) & The truncated tops of S4, G, D and F form a flat erosion surface (unconformity) \\ \hline
7 & Deposition of flat-lying S5 & Superposition + unconformity: S5 rests on the erosion surface and is \emph{not} offset by F, confirming it post-dates the fault \\ \hline
\end{tabularx}
\end{center}

\textbf{Key reasoning points:}
\begin{itemize}
\item G is younger than the folding because it cross-cuts the folded layers; its margins are straight and undeformed.
\item D is younger than G because the dyke cuts the granite.
\item F is younger than D because the fault displaces the dyke.
\item S5 is younger than F because the fault does \emph{not} continue into S5 — the fault is \emph{sealed} by S5.
\item The erosion surface between the faulted/tilted pile and horizontal S5 is an \cle{angular unconformity} (over the sediments) and a \cle{nonconformity} (where S5 rests on granite).
\end{itemize}

\begin{attention}[Common mistake]
Many students place the fault \emph{before} the dyke or forget the erosion event. Rule of thumb: \emph{a feature that is cut is older than the feature that cuts it; a feature that seals (covers) a structure is younger than that structure.}
\end{attention}
\end{corrige}

\begin{corrige}[Exercise 6 — Biostratigraphic Correlation]
\textbf{(a) Age range of each assemblage.} An assemblage can only date from an interval during which \emph{all} its species coexisted — i.e., the \emph{overlap} of their ranges (concurrent range zone):

\begin{itemize}
\item \textbf{Assemblage 1 (A + C):} A ranges 145--130 Ma; C ranges 140--120 Ma. Overlap = \textbf{140--130 Ma}.
\item \textbf{Assemblage 2 (B + D):} B ranges 135--110 Ma; D ranges 125--90 Ma. Overlap = \textbf{125--110 Ma}.
\item \textbf{Assemblage 3 (C + E):} C ranges 140--120 Ma; E ranges 115--85 Ma. Overlap = \textbf{120--115 Ma}.
\end{itemize}

\textbf{(b) Index fossil for the Early Cretaceous ($\approx$ 145--100 Ma):} \textbf{Species A} is the best index fossil: its total range (145--130 Ma, i.e., 15 Ma) is the \emph{shortest} of the species present in the Early Cretaceous part of the succession, and it falls entirely within the Early Cretaceous. Species C (20 Ma range) is less precise; B and D extend into the Late Cretaceous; E is essentially Late Cretaceous. A good index fossil = \emph{short range + wide distribution + abundance + easy identification}.

\textbf{(c) Effect of the ash-bed age (112 $\pm$ 2 Ma) on Assemblage 2:} The biostratigraphic bracket for Assemblage 2 is 125--110 Ma. The radiometric age of \SI{112 \pm 2}{Ma} (i.e., 110--114 Ma) lies within this bracket and is \emph{consistent} with it. It refines the correlation by:
\begin{itemize}
\item \emph{pinning} the assemblage to the narrower interval \textbf{114--110 Ma} (intersection of 125--110 Ma with 110--114 Ma);
\item providing an \emph{absolute} (numerical) calibration that can be transferred to other sections where the same assemblage occurs but no ash bed exists;
\item serving as an \emph{isochronous marker horizon} for basin-wide correlation of the Rio del Rey succession.
\end{itemize}
This illustrates the power of \emph{integrating} biostratigraphy (relative) with radiometric dating (absolute).
\end{corrige}

\begin{corrige}[Exercise 7 — Concordia-Discordia Interpretation]
\textbf{(a) Sketch of the Concordia diagram:}

\begin{popfigure}
\begin{center}
\begin{tikzpicture}[scale=0.9]
% Axes
\draw[->, thick] (0,0) -- (8.5,0) node[below left, font=\small] {\ce{^{206}Pb}/\ce{^{238}U}};
\draw[->, thick] (0,0) -- (0,6.5) node[rotate=90, above left, font=\small, align=center] {\ce{^{207}Pb}/\ce{^{235}U}};
% Concordia curve (schematic)
\draw[popblue, very thick, domain=0.5:7.5, samples=60] plot (\x, {0.11*\x*\x});
\node[popblue, font=\small] at (6.6,5.9) {Concordia};
% Discordia line through the three points (scaled x100 for plotting)
% Points: (1.85,1.25), (3.10,2.10), (5.15,3.50) -> line slope ~0.68
\draw[poporange, very thick, dashed] (0.6,0.41) -- (7.8,5.31);
\node[poporange, font=\small, align=center] at (7.3,4.6) {Discordia};
% Analyses
\filldraw[popdark] (1.85,1.25) circle (2pt) node[below right, font=\small] {1};
\filldraw[popdark] (3.10,2.10) circle (2pt) node[below right, font=\small] {2};
\filldraw[popdark] (5.15,3.50) circle (2pt) node[below right, font=\small] {3};
% Intercepts
\filldraw[popgreen!50!black] (6.9,5.23) circle (2.5pt) node[right, font=\small, align=left] {Upper intercept\\ $\approx$ 2050 Ma};
\filldraw[popgreen!50!black] (0.75,0.51) circle (2.5pt) node[right, font=\small, align=left] {Lower intercept\\ $\approx$ 550 Ma};
\end{tikzpicture}
\end{center}
\legende{Concordia diagram: the three discordant zircon analyses define a Discordia line intersecting the Concordia curve at two points.}
\end{popfigure}

\textbf{(b) Geological interpretation:}
\begin{itemize}
\item \textbf{Upper intercept = \SI{2050 \pm 20}{Ma}:} the \emph{original crystallisation age} of the zircons — i.e., the emplacement/metamorphic crystallisation of the granitic protolith of the gneiss during the \emph{Eburnean (Palaeoproterozoic) orogeny} that built much of the Congo Craton in southern Cameroon.
\item \textbf{Lower intercept = \SI{550 \pm 30}{Ma}:} the age of a later \emph{disturbance} that caused lead loss — here the \emph{Pan-African orogeny} ($\approx$ 600--500 Ma), which thermally and tectonically reworked the craton margins and opened the zircon U-Pb system.
\end{itemize}

\textbf{(c) Why zircons yield discordant ages (two mechanisms):}
\begin{enumerate}
\item \textbf{Diffusive Pb loss during a thermal event:} reheating (metamorphism, nearby intrusion) allows radiogenic lead to diffuse out of the crystal lattice, while uranium is retained; the measured \ce{Pb}/\ce{U} ratios fall, and the point moves \emph{down} along the Discordia.
\item \textbf{Metamictisation and leaching:} radiation damage from U decay destroys the crystal lattice (metamict zircon); fluids then percolate and leach radiogenic Pb. (Also acceptable: \emph{inheritance} of older zircon cores, or \emph{recent Pb loss} by weathering pulling points toward the origin.)
\end{enumerate}
\end{corrige}

\begin{corrige}[Exercise 8 — Unconformities: Identification and Field Evidence]
\textbf{Schematic stratigraphic column (bottom $\rightarrow$ top):}

\begin{popfigure}
\begin{center}
\begin{tikzpicture}[scale=0.85]
% Precambrian gneiss
\draw[fill=poppurple!30, pattern=north west lines] (0,0) rectangle (4,2);
\node[align=center, font=\small] at (2,1) {Precambrian\\ gneiss};
% Nonconformity
\draw[popred, very thick, decorate, decoration={snake, amplitude=0.06cm, segment length=0.4cm}] (0,2) -- (4,2);
\node[popred, right, font=\small\bfseries] at (4.1,2) {NONCONFORMITY};
% Cambrian sandstone
\draw[fill=yellow!50!poporange] (0,2) rectangle (4,4);
\node[font=\small] at (2,3) {Cambrian sandstone};
% Silurian folded shales -> draw as tilted block
\draw[fill=popblueL!50] (0,4) -- (4,5.2) -- (4,7.2) -- (0,6) -- cycle;
\node[align=center, font=\small] at (2,5.6) {Folded/tilted\\ Silurian shales};
% Angular unconformity
\draw[popred, very thick] (0,6) -- (4,7.2);
\node[popred, right, font=\small\bfseries] at (4.1,6.6) {ANGULAR UNCONFORMITY};
% Devonian sandstone
\draw[fill=poporange!40] (0,6) -- (0,8) -- (4,8) -- (4,7.2) -- cycle;
\node[font=\small] at (2,7.4) {Devonian sandstone (horizontal)};
% Carboniferous limestone with palaeosol
\draw[fill=popgreenL!50] (0,8) rectangle (4,10.5);
\node[font=\small] at (2,9.6) {Carboniferous limestone};
\draw[fill=brown!60] (0,9) rectangle (4,9.25);
\node[right, font=\small] at (4.1,9.1) {palaeosol};
\draw[popred, very thick, dashed] (0,9) -- (4,9);
\node[popred, right, font=\small\bfseries] at (4.1,8.7) {DISCONFORMITY};
% Jurassic clays
\draw[fill=popmuted!40] (0,10.5) rectangle (4,12);
\node[font=\small] at (2,11.1) {Jurassic clay (lower)};
\draw[popred, very thick, dotted] (0,12) -- (4,12);
\node[popred, right, font=\small\bfseries] at (4.1,12) {PARACONFORMITY};
\draw[fill=popmuted!60] (0,12) rectangle (4,13.5);
\node[font=\small] at (2,12.75) {Jurassic clay (upper)};
\end{tikzpicture}
\end{center}
\legende{Stratigraphic column showing the four types of unconformity.}
\end{popfigure}

\textbf{Field criteria (two per type):}

\begin{center}
\begin{tabularx}{\linewidth}{|l|X|X|}\hline
\rowcolor{popblueL}
\textbf{Type} & \textbf{Criterion 1} & \textbf{Criterion 2} \\ \hline
Nonconformity & Sedimentary rocks rest directly on eroded \emph{crystalline} (igneous/metamorphic) basement & Basal conglomerate containing pebbles of the underlying gneiss; weathered (saprolitic) top of basement \\ \hline
Angular unconformity & Marked \emph{discordance of dip} between truncated tilted/folded strata below and horizontal strata above & Truncation of beds, dykes and folds at the erosion surface; basal conglomerate above the surface \\ \hline
Disconformity & Erosion surface (irregular, channelled) between \emph{parallel} strata & \emph{Palaeosol}, root traces, pebble lag or karstic surface at the contact; fossil gap \\ \hline
Paraconformity & Strata perfectly parallel with \emph{no visible erosion surface} & Detectable only by a \emph{biostratigraphic gap} (missing fossil zones) between the two units \\ \hline
\end{tabularx}
\end{center}

\begin{attention}[Common mistake]
Do not confuse \emph{disconformity} (visible erosion surface between parallel beds) with \emph{paraconformity} (no visible surface; gap proved only by fossils). In the field, the palaeosol is your strongest evidence for a disconformity.
\end{attention}
\end{corrige}

\begin{corrige}[Exercise 9 — Integrated Relative and Absolute Dating (20 marks)]
\textbf{(a) Relative chronological sequence [8 marks — 1 mark per correctly placed event]:}

\begin{enumerate}
\item Formation of the Basement Gneiss (Bg) — oldest unit (host to everything).
\item Erosion of the gneiss to form the \cle{nonconformity} surface.
\item Deposition of the Cretaceous Sandstones (Ks) on the eroded basement (superposition).
\item Tilting/folding of Ks (strata must exist before being deformed).
\item Intrusion of the Granite (Gr) — cross-cuts Bg and the tilted Ks (cross-cutting relationships).
\item Emplacement of the Dolerite Dyke (Dd) — cuts Bg, Ks \emph{and} Gr (cross-cutting).
\item Faulting (F) — offsets Bg, Ks, Gr and Dd.
\item Erosion — planation surface truncating F, Dd, Gr and Ks.
\item Deposition of the horizontal Tertiary Sands (Ts) — seals the fault and the erosion surface (youngest event).
\end{enumerate}
\emph{Examiner expectation:} any logically consistent order respecting (i) Ks before tilting, (ii) Gr after tilting, (iii) Dd after Gr, (iv) F after Dd, (v) Ts last, earns full credit. State the principle for at least the key steps.

\textbf{(b) Minimum absolute age of the tilting [4 marks]:}
The granite (U-Pb zircon age \SI{88.5 \pm 0.5}{Ma}) intrudes the \emph{already tilted} Cretaceous sandstones. Therefore the tilting is \emph{older} than the granite crystallisation:
\[
t_{\text{tilting}} > 88.5\ \text{Ma}.
\]
Considering the analytical uncertainty, the tilting must be older than \SI{89}{Ma} (i.e., older than $88.5 + 0.5$ Ma). \emph{Reasoning [2 marks]:} cross-cutting relationships — an intrusion is younger than the structures it cuts; the zircon crystallisation age dates the granite's solidification, so the deformation it cross-cuts must pre-date it. \emph{Note:} the \emph{maximum} age of tilting is given by the age of the youngest tilted sandstone (Cretaceous).

\textbf{(c) Reliability of the K-Ar dyke age vs the U-Pb granite age [4 marks]:}
\begin{itemize}
\item \textbf{Closure temperatures:} U-Pb in zircon closes at very high temperature ($> \SI{900}{\ensuremath{{}^{\circ}\mathrm{C}}}$), so the zircon age dates \emph{crystallisation} faithfully and resists later reheating. K-Ar in whole-rock closes at much lower temperature ($\approx$ \SIrange{150}{350}{\ensuremath{{}^{\circ}\mathrm{C}}} depending on mineral), and argon is an \emph{inert gas} that diffuses out easily.
\item \textbf{Consequences:} the dyke was emplaced close to a hot granite; contact reheating and later burial could have caused \emph{partial Ar loss}, making the K-Ar age (\SI{42}{Ma}) \emph{younger than the true emplacement age}. Whole-rock analyses also risk \emph{excess argon} or alteration (devitrification, weathering of groundmass).
\item \textbf{Conclusion:} the K-Ar age is a \emph{minimum} and less reliable constraint; it may date cooling or a resetting event rather than emplacement. The U-Pb zircon age is far more robust. A better method for the dyke would be \ce{^{40}Ar}/\ce{^{39}Ar} step-heating on fresh plagioclase or U-Pb on baddeleyite/zircon.
\end{itemize}
\emph{Mark scheme:} 1 mark closure-temperature contrast; 1 mark Ar mobility/loss; 1 mark geological consequence (age too young / reset); 1 mark reasoned conclusion or better alternative.

\textbf{(d) Field evidence for the nonconformity Bg/Ks [4 marks — any 4 $\times$ 1]:}
\begin{itemize}
\item Cretaceous sedimentary rocks resting \emph{directly} on eroded crystalline gneiss.
\item \emph{Basal conglomerate/breccia} at the base of Ks containing angular pebbles of the underlying gneiss.
\item \emph{Weathered, saprolitic or palaeosol horizon} at the top of the gneiss (ancient land surface).
\item \emph{Truncation} of foliation, veins and structures of the gneiss at the contact.
\item Absence of metamorphism in the sandstone above (contact is depositional, not intrusive) and no chilled margin in the gneiss.
\item Marked \emph{age gap} (hiatus) between the two units.
\end{itemize}
\end{corrige}

\begin{corrige}[Exercise 10 — Essay: Magnetostratigraphy and Chemostratigraphy in Cameroon (25 marks)]
\textbf{Indicative mark scheme:} definitions and principles [5]; GPTS construction and use [5]; chemostratigraphic proxies [5]; strengths/limitations applied to the three basins [6]; integrated case study [4]. Quality of structure, terminology and labelled diagrams may be rewarded within these bands.

\textbf{Model answer (skeleton to be developed in full prose):}

\textbf{1. Definitions and principles [5].}
\emph{Magnetostratigraphy} correlates rocks using the record of \emph{reversals of the geomagnetic field} locked into sediments and lavas as detrital or thermal remanent magnetisation; the principle is that polarity reversals are \emph{globally synchronous} and instantaneous on geological timescales. \emph{Chemostratigraphy} correlates using \emph{secular variations in the chemical and isotopic composition} of seawater recorded in marine sediments; the principle is that the ocean is well mixed, so isotopic shifts (e.g., in $\delta^{13}$C) are recorded simultaneously in marine carbonates worldwide.

\textbf{2. The GPTS [5].}
The Geomagnetic Polarity Time Scale is built by (i) measuring the polarity of \emph{radiometrically dated} lava flows worldwide, (ii) stacking magnetostratigraphic sections tied to biostratigraphy, and (iii) reading the \emph{marine magnetic anomaly stripes} of ocean crust (sea-floor spreading provides a continuous tape-recorder of reversals). The result is a calibrated sequence of \emph{chrons} (normal/reversed). To use it: sample an oriented section, perform stepwise demagnetisation to isolate the characteristic remanence, establish the local polarity zonation, then \emph{pattern-match} the succession of magnetozones to the GPTS, anchored by biostratigraphic or radiometric tie-points. Each reversal boundary is an \emph{isochronous} correlation line.

\textbf{3. Chemostratigraphic proxies for Cretaceous marine sediments [5].}
\begin{itemize}
\item $\delta^{13}$C (carbonate and organic carbon): positive excursions mark \emph{Oceanic Anoxic Events} (e.g., OAE1a in the Aptian, OAE2 at the Cenomanian--Turonian boundary) — excellent regional markers in the Douala and Rio del Rey black shales.
\item $\delta^{18}$O: palaeotemperature/ice-volume proxy; useful for climatic correlation but sensitive to \emph{diagenetic resetting}.
\item $^{87}$Sr/$^{86}$Sr: seawater Sr ratio evolves smoothly through time; gives \emph{numerical ages} for well-preserved low-Mg calcite (e.g., belemnites, oysters), provided diagenetic alteration is screened out.
\end{itemize}

\textbf{4. Strengths and limitations in the Cameroon Atlantic basins [6].}
\begin{center}
\begin{tabularx}{\linewidth}{|l|X|X|}\hline
\rowcolor{popblueL}
& \textbf{Strengths} & \textbf{Limitations} \\ \hline
Magnetostratigraphy & Global, isochronous markers; works in barren or non-marine (continental) beds where fossils fail — useful for the Aptian red beds of Douala/Kribi-Campo & Tropical \emph{deep weathering} and lateritisation destroy remanence; Cretaceous Normal Superchron ($\approx$ 121--83 Ma) offers \emph{no reversals} for much of the mid-Cretaceous; needs oriented cores \\ \hline
Chemostratigraphy & High-resolution; correlates where facies change makes litho- and biostratigraphy ambiguous; OAE excursions are strong markers in the organic-rich shales & Diagenesis (recrystallisation, meteoric fluids) resets $\delta^{18}$O and Sr; local restriction of the young Atlantic basins can decouple $\delta^{13}$C from the global signal; needs unaltered material \\ \hline
\end{tabularx}
\end{center}

\textbf{5. Integrated case study [4] (hypothetical but realistic).}
In the Rio del Rey Basin, an organic-rich black shale was assigned to the Turonian on \emph{Inoceramus} evidence, while an apparently equivalent shale in the Douala Basin was dated as Cenomanian on ammonites — a correlation conflict. An integrated study resolved it: (i) the $\delta^{13}$C positive excursion of \emph{OAE2} was identified in \emph{both} shales, proving they record the same global event at the Cenomanian--Turonian boundary ($\approx$ 93.9 Ma); (ii) a bentonite within the Douala shale gave a U-Pb zircon age of \SI{93.5}{Ma}, calibrating the excursion; (iii) magnetostratigraphy placed both sections just above the end of the Cretaceous Normal Superchron, consistent with the boundary age. The apparent conflict came from \emph{diachronous first appearances} of the index fossils; the chemostratigraphic and radiometric tie-points provided the correct isochronous framework.

\emph{Examiner expectations:} explicit definitions; the \emph{global synchroneity} argument for both methods; at least three correctly used proxies; balanced evaluation (not a catalogue of advantages only); a coherent integrated example; correct geological vocabulary throughout.
\end{corrige}

\begin{corrige}[Exercise 11 — Multi-Section Correlation and Sequence Stratigraphy (20 marks)]
\textbf{(a) Correlation panel [6 marks].} Fence diagram at 1 cm = 100 m (so 1000 m = 10 cm):

\begin{popfigure}
\begin{center}
\begin{tikzpicture}[scale=0.62]
% vertical scale: 1 unit = 100 m, depth downward
% BH-1
\begin{scope}[shift={(0,0)}]
\draw[fill=yellow!50!poporange] (0,0) rectangle (2,2);
\draw[popred, very thick, decorate, decoration={snake, amplitude=0.05cm, segment length=0.3cm}] (0,2) -- (2,2);
\draw[fill=poporange!50] (0,2) rectangle (2,6);
\draw[fill=poppurple!30, pattern=north west lines] (0,6) rectangle (2,8);
\node[font=\small\bfseries] at (1,-0.5) {BH-1};
\node[font=\tiny, rotate=90] at (1,1) {sands, lignite};
\node[font=\tiny, rotate=90] at (1,4) {red beds, congl.};
\node[font=\tiny, rotate=90] at (1,7) {basement};
\end{scope}
% BH-2
\begin{scope}[shift={(5,0)}]
\draw[fill=popmuted!40] (0,0) rectangle (2,2);
\draw[popred, very thick, decorate, decoration={snake, amplitude=0.05cm, segment length=0.3cm}] (0,2) -- (2,2);
\draw[fill=yellow!50!poporange] (0,2) rectangle (2,4.5);
\draw[fill=poporange!50] (0,4.5) rectangle (2,6);
\draw[fill=poppurple!30, pattern=north west lines] (0,6) rectangle (2,8);
\node[font=\small\bfseries] at (1,-0.5) {BH-2};
\node[font=\tiny, rotate=90] at (1,1) {clays, forams};
\node[font=\tiny, rotate=90] at (1,3.2) {sands, lignite};
\node[font=\tiny, rotate=90] at (1,5.2) {red beds};
\node[font=\tiny, rotate=90] at (1,7) {basement};
\end{scope}
% BH-3
\begin{scope}[shift={(10,0)}]
\draw[fill=popmuted!40] (0,0) rectangle (2,2);
\draw[fill=popblueL!50] (0,2) rectangle (2,6);
\draw[fill=popdark!60] (0,6) rectangle (2,10);
\node[font=\small\bfseries] at (1,-0.5) {BH-3};
\node[font=\tiny, rotate=90] at (1,1) {clays, forams};
\node[font=\tiny, rotate=90] at (1,4) {marls, Orbitolina};
\node[font=\tiny, rotate=90, white] at (1,8) {black shales, Inoceramus};
% ash bed
\draw[popgreen!50!black, very thick] (0,5.5) -- (2,5.5);
\node[popgreen!50!black, right, font=\tiny] at (2,5.5) {ash 93.5 Ma};
\end{scope}
% Correlation lines
\draw[popblue, thick, dashed] (2,0) -- (5,0); 
\draw[popblue, thick, dashed] (7,0) -- (10,0);
\draw[popred, thick, dashed] (2,2) -- (5,2);
\draw[popred, thick, dashed] (7,2) -- (10,2) node[midway, above, font=\tiny] {U1 (Base Tertiary)};
\draw[popblue, thick, dashed] (2,6) -- (5,6);
\draw[popblue, thick, dashed] (7,6) -- (10,6);
% depth axis
\draw[->] (-1,0) -- (-1,10) node[above, font=\small] {depth};
\foreach \d in {0,2,4,6,8,10} \node[left, font=\tiny] at (-1,\d) {\pgfmathparse{int(\d*100)}\pgfmathresult\ m};
\end{tikzpicture}
\end{center}
\legende{Correlation panel (fence diagram) of the three boreholes. Dashed lines: correlated unit boundaries; U1 = Base Tertiary unconformity.}
\end{popfigure}

\emph{Mark scheme:} correct vertical scale and depth axis [1]; correct lithologies per borehole [2]; U1 correctly drawn and correlated between BH-1 and BH-2 and its seismic significance noted (it passes within the BH-3 marl succession as a conformity/condensed interval offshore) [1]; basement correlation [1]; ash-bed marker shown in BH-3 [1].

\textbf{(b) Systems tracts in BH-3 (450--1000 m) [6 marks]:}
\begin{itemize}
\item \textbf{Lowstand Systems Tract (LST) — 600--800 m (lower black shales / basal unit):} following the relative sea-level fall recorded onshore by the red beds and conglomerates of BH-1/BH-2 (fluvial, lowstand deposits), the offshore equivalent is the basal black shale; the coarse clastics onshore = lowstand wedge. [2 marks]
\item \textbf{Transgressive Systems Tract (TST) — $\approx$ 550--600 m and the marls:} the Aptian--Albian eustatic transgression drowned the margin: onshore sands/lignites (BH-2, 200--450 m equivalent) pass offshore into marine marls with \emph{Orbitolina} (Aptian) — deepening-upward, retrogradational stacking; the black shales with TOC 4\% record maximum restriction/anoxia near the \emph{maximum flooding surface} (condensed section). [2 marks]
\item \textbf{Highstand Systems Tract (HST) — 450--550 m (upper marls):} Late Cretaceous highstand: normal marine marls prograding over the condensed black shales; the ash bed at 550 m (\SI{93.5}{Ma}, Cenomanian--Turonian boundary) sits near the TST/HST transition — consistent with the global maximum flooding near that boundary. [2 marks]
\end{itemize}
\emph{Justification must cite:} lithology (shale = deep/anoxic = flooding; marl = open marine; conglomerate/red beds = continental lowstand) and fossils (\emph{Orbitolina} = shallow-marine Aptian; \emph{Inoceramus} = Turonian-Coniacian open marine).

\textbf{(c) Sedimentation rate of the black shale (800--1000 m) [4 marks]:}
The black shale contains \emph{Inoceramus} (Turonian--Coniacian, $\approx$ 93.9--86.3 Ma). Using the ash age \SI{93.5}{Ma} at 550 m as a top anchor for the overlying marls and the biostratigraphic range for the black shale interval:
\begin{itemize}
\item Thickness of black shale interval (800--1000 m): $1000 - 800 = \SI{200}{m}$.
\item Duration (Turonian--Coniacian): $\approx 93.9 - 86.3 = \SI{7.6}{Ma}$.
\item Rate $= \dfrac{200\ \text{m}}{7.6\ \text{Ma}} \approx \SI{26.3}{m/Ma}$ (i.e., $\approx \SI{26}{mm/ka}$, uncompacted).
\end{itemize}
\emph{Accept} any coherent calculation (e.g., using 93.5 Ma for the top and 86 Ma for the base of the unit gives $\approx$ \SI{27}{m/Ma}). Marks: correct thickness [1]; correct time interval justified from fossil range/ash age [2]; correct arithmetic with units [1]. Note the result is an \emph{average, uncompacted} rate.

\textbf{(d) Hydrocarbon potential [4 marks — 1 per element]:}
\begin{itemize}
\item \textbf{Source rock:} the BH-3 black shales (TOC 4\%, Turonian--Coniacian) — an excellent oil-prone source rock, equivalent to the productive Cretaceous source intervals of the West African Atlantic margin.
\item \textbf{Reservoir:} the yellow sands of BH-1/BH-2 (and the basal conglomerates/red-bed sandstones) — porous, laterally continuous nearshore/fluvial sandstones.
\item \textbf{Seal:} the grey marine clays and marls (top unit in all boreholes, and the marls of BH-3) — impermeable cap rocks; the unconformity U1 itself can act as a seal where clays overlie truncated sands.
\item \textbf{Trap:} \emph{stratigraphic pinch-out traps} where the reservoir sands of BH-2 pass laterally into offshore marls (BH-3); \emph{unconformity (truncation) traps} beneath U1; possible structural traps along growth faults typical of the Douala Basin.
\end{itemize}
\emph{Examiner expectation:} each element must be tied to a \emph{specific named unit in a specific borehole}, not generic statements.
\end{corrige}

\newpage

\coursec{Summary sheet}

\begin{popfigure}
\begin{tikzpicture}[
  mindmap,
  concept color=popblue,
  level 1/.style={level distance=3.5cm, sibling angle=360/7},
  level 2/.style={level distance=2.5cm},
  every node/.style={concept, execute at begin node=\hskip0pt},
  root concept/.append style={concept color=popblue, font=\large\bfseries},
  level 1 concept/.append style={font=\normalsize\bfseries},
  level 2 concept/.append style={font=\small}
]
\node[root concept] {Stratigraphy}
  child[concept color=poporange] { node {Scope of Stratigraphy} }
  child[concept color=popgreen] { node {Fundamental Principles} }
  child[concept color=poppurple] { node {Geologic Time Scale} }
  child[concept color=poppink] { node {Dating Methods} }
  child[concept color=popgold] { node {Stratigraphic Contacts} }
  child[concept color=popdark] { node {Sequencing Events} }
  child[concept color=popblue] { node {Correlation} };
\end{tikzpicture}
\legende{Mind map of the Stratigraphy chapter (Upper Sixth GCE).}
\end{popfigure}

\begin{bilan}
\textbf{Key Definitions}
\begin{itemize}
\item \cle{Stratigraphy}: The branch of geology concerned with the study of rock layers (strata), their geometry, distribution, and chronological relationships in time and space.
\item \cle{Principle of Superposition}: In an undeformed sequence of sedimentary or volcanic rocks, the oldest layers are at the bottom and the youngest at the top.
\item \cle{Principle of Original Horizontality}: Sedimentary layers are originally deposited horizontally under gravity.
\item \cle{Principle of Lateral Continuity}: Layers extend laterally in all directions until they thin out, grade into another facies, or meet a barrier.
\item \cle{Principle of Cross-Cutting Relationships}: A geological feature (fault, intrusion, erosion surface) that cuts across another is younger than the feature it cuts.
\item \cle{Principle of Inclusions}: Fragments (xenoliths, clasts) included in a host rock are older than the host rock.
\item \cle{Principle of Faunal Succession}: Fossil assemblages succeed one another in a definite, recognisable order, allowing relative dating and correlation.
\item \cle{Unconformity}: A surface representing a significant gap (hiatus) in the geological record due to erosion or non-deposition; types: angular unconformity, disconformity, nonconformity, paraconformity.
\item \cle{Correlation}: The process of demonstrating the equivalence in age or time of rock units in different geographical areas.
\item \cle{Lithostratigraphic Unit}: A body of rock defined by its lithology (formation, member, group).
\item \cle{Biostratigraphic Unit}: A body of rock defined by its fossil content (biozone).
\item \cle{Chronostratigraphic Unit}: A body of rock formed during a specific interval of geological time (system, series, stage).
\end{itemize}

\textbf{Laws, Formulas \& Theorems to Know}
\begin{itemize}
\item \textbf{Steno's Laws (1669)}: Superposition, Original Horizontality, Lateral Continuity — examinable as the foundation of relative dating.
\item \textbf{Hutton's Cross-Cutting Relationships (1795)} and \textbf{Lyell's Inclusions}: Key for deciphering intrusive and tectonic events.
\item \textbf{Smith's Principle of Faunal Succession (1790s)}: Basis of biostratigraphy and the geological column.
\item \textbf{Radiometric Age Equation}: $t = \frac{1}{\lambda} \ln\left(1 + \frac{D}{P}\right)$, where $\lambda$ = decay constant, $D$ = daughter isotope atoms, $P$ = parent isotope atoms remaining.
\item \textbf{Half-life Relation}: $t_{1/2} = \frac{\ln 2}{\lambda}$. Common systems: \ce{^{238}U -> ^{206}Pb} ($t_{1/2}=4.47$ Ga), \ce{^{235}U -> ^{207}Pb} ($t_{1/2}=0.704$ Ga), \ce{^{40}K -> ^{40}Ar} ($t_{1/2}=1.25$ Ga), \ce{^{87}Rb -> ^{87}Sr} ($t_{1/2}=48.8$ Ga), \ce{^{14}C -> ^{14}N} ($t_{1/2}=5730$ yr).
\item \textbf{Geologic Time Scale (ICS 2024)}: Know the major boundaries: Precambrian/Cambrian (541 Ma), Cambrian/Ordovician (485 Ma), Permian/Triassic (252 Ma), Triassic/Jurassic (201 Ma), Cretaceous/Paleogene (66 Ma), Neogene/Quaternary (2.58 Ma). Be able to place Cameroon events: CVL initiation (~66 Ma), Mount Cameroon recent activity (<1 Ma), Pan-African orogeny (~600–500 Ma).
\end{itemize}

\textbf{Methods}
\begin{enumerate}
\item \textbf{Relative Dating}: Apply stratigraphic principles to a cross-section or block diagram to produce a sequence of events (oldest $\to$ youngest). Always state the principle used for each step.
\item \textbf{Absolute Dating}: Choose appropriate isotopic system based on rock type and age range (e.g., U-Pb on zircon for Precambrian basement, K-Ar/Ar-Ar on volcanic rocks for CVL, \ce{^{14}C} for Holocene sediments). Correct for initial daughter, contamination, and closed-system assumption.
\item \textbf{Stratigraphic Logging}: Measure and describe a section: lithology, thickness, sedimentary structures, fossil content, contact types. Draw a graphic log with standard symbols.
\item \textbf{Unconformity Recognition}: Identify angular discordance (tilted beds below, flat above), erosional relief, palaeosols, basal conglomerates with clasts from underlying unit.
\item \textbf{Correlation Techniques}: 
  \begin{itemize}
  \item \emph{Lithostratigraphy}: Match distinctive rock units (marker beds).
  \item \emph{Biostratigraphy}: Use index fossils (wide geographic range, short time range).
  \item \emph{Chronostratigraphy}: Tie to global standard stages (GSSPs).
  \item \emph{Magnetostratigraphy}: Match magnetic polarity reversal pattern to GPTS.
  \item \emph{Chemostratigraphy}: Correlate isotopic excursions (e.g., $\delta^{13}$C).
  \end{itemize}
\item \textbf{Cross-Section Construction}: Use topographic profile, dip/strike data, and stratigraphic thicknesses. Project contacts vertically or along dip direction. Show vertical exaggeration if used.
\end{enumerate}

\textbf{Common Mistakes}
\begin{itemize}
\item Confusing \emph{relative} (order) and \emph{absolute} (numerical) dating; mixing up principles (e.g., applying superposition to overturned folds without recognising overturning via younging indicators).
\item Treating every contact as an unconformity; a conformable contact shows continuous deposition.
\item Misidentifying a sill as a lava flow (use cross-cutting: sill intrudes both above and below; flow only underlies).
\item Correlating solely on lithology without fossil or radiometric control — lateral facies changes are common in the Benue Trough.
\item Radiometric errors: using wrong decay constant, forgetting to convert half-life to $\lambda$, ignoring initial daughter isotope (use isochron method), or dating a metamorphic event instead of crystallisation.
\item Not knowing the current International Chronostratigraphic Chart; the GCE expects the latest ICS version.
\item In sequencing questions, omitting events (e.g., erosion before an unconformity, tilting before angular unconformity).
\end{itemize}

\textbf{GCE Exam Tips}
\begin{itemize}
\item \textbf{Question Types}: Expect (a) block diagram / cross-section interpretation with sequencing (10–15 marks), (b) radiometric calculation (5–8 marks), (c) correlation of two or more columns using fossils/lithology (8–10 marks), (d) description of unconformity types with diagrams (5 marks), (e) essay on stratigraphic principles or dating methods (10–15 marks).
\item \textbf{Show All Working}: In age calculations, write the formula, substitute values with units, give answer in Ma (million years) or Ga. Use $\lambda = \frac{\ln 2}{t_{1/2}}$.
\item \textbf{Terminology Precision}: Distinguish \emph{angular unconformity}, \emph{disconformity}, \emph{nonconformity}, \emph{paraconformity}. Use \emph{formation}, \emph{member}, \emph{group} correctly.
\item \textbf{Cameroon Context}: Be prepared to apply principles to the Cameroon Volcanic Line (e.g., date a basalt flow from Mount Cameroon using K-Ar), the Congo Craton (Pan-African granites dated by U-Pb on zircon), or the Benue Trough (Cretaceous marine transgressions correlated by ammonites).
\item \textbf{Diagrams}: Draw large, labelled diagrams (cross-sections, stratigraphic columns, fossil ranges). Use a ruler. Indicate younging direction.
\item \textbf{Integration Questions}: Combine principles, dating, and correlation to write a geological history. Structure: 1. Oldest basement, 2. Deformation/metamorphism, 3. Unconformity, 4. Sedimentary cover with fossils, 5. Intrusive/volcanic events, 6. Recent erosion.
\item \textbf{Time Management}: Allocate ~1 minute per mark. For a 15-mark sequencing question, spend 15 minutes: 5 min analysis, 5 min writing sequence with principles, 5 min neat diagram.
\end{itemize}
\end{bilan}

\findecours

\end{document}
